\documentclass[11pt,oneside]{book}

\usepackage[T1]{fontenc}
\usepackage{lmodern}
\usepackage{microtype}
\usepackage{amsmath,amssymb,amsthm,mathtools}
\usepackage[margin=1.02in,headheight=14pt]{geometry}
\usepackage{enumitem}
\usepackage{booktabs,longtable,array}
\usepackage{fancyhdr}
\usepackage{setspace}
\usepackage{xcolor}
\usepackage[colorlinks=true,linkcolor=blue!55!black,citecolor=blue!55!black,urlcolor=blue!55!black]{hyperref}
\usepackage[nameinlink,capitalise,noabbrev]{cleveref}

\setstretch{1.055}
\setlength{\parindent}{1.35em}
\setlength{\parskip}{0.18em}
\setlist{itemsep=0.25em,topsep=0.35em}
\sloppy

\pagestyle{fancy}
\fancyhf{}
\fancyhead[L]{\small Phase retrieval for cubic NLS}
\fancyhead[R]{\small Self-contained proof}
\fancyfoot[C]{\thepage}

\newtheorem{theorem}{Theorem}[chapter]
\newtheorem{proposition}[theorem]{Proposition}
\newtheorem{lemma}[theorem]{Lemma}
\newtheorem{corollary}[theorem]{Corollary}
\newtheorem{claim}[theorem]{Claim}
\theoremstyle{definition}
\newtheorem{definition}[theorem]{Definition}
\newtheorem{convention}[theorem]{Convention}
\theoremstyle{remark}
\newtheorem{remark}[theorem]{Remark}
\newtheorem{checkpoint}[theorem]{Dependency checkpoint}

\newcommand{\R}{\mathbb R}
\newcommand{\C}{\mathbb C}
\newcommand{\T}{\mathbb T}
\newcommand{\N}{\mathbb N}
\newcommand{\Q}{\mathbb Q}
\newcommand{\D}{\mathcal D}
\newcommand{\Ssch}{\mathcal S}
\newcommand{\supp}{\operatorname{supp}}
\newcommand{\1}{\mathbf 1}
\newcommand{\loc}{\mathrm{loc}}
\newcommand{\sgn}{\operatorname{sgn}}
\newcommand{\esssup}{\operatorname*{ess\,sup}}
\newcommand{\Rea}{\operatorname{Re}}
\newcommand{\Ima}{\operatorname{Im}}
\newcommand{\dd}{\,\mathrm d}
\newcommand{\norm}[1]{\left\lVert #1\right\rVert}
\newcommand{\abs}[1]{\left\lvert #1\right\rvert}
\newcommand{\pair}[2]{\left\langle #1,#2\right\rangle}
\newcommand{\FT}{\mathcal F}
\newcommand{\Pfree}{i\partial_t+\Delta}
\newcommand{\NLS}{\mathrm{NLS}}
\newcommand{\e}{\mathrm e}
\newcommand{\iu}{\mathrm i}
\newcommand{\eps}{\varepsilon}
\newcommand{\weakto}{\rightharpoonup}

\title{\textbf{Phase Retrieval for the One-Dimensional Cubic NLS}\\[0.45em]
A self-contained proof in the global $L^2$ solution class\\[0.8em]
\large Expanded formalization-ready edition with endpoint and closure details}
\author{}
\date{July 2026}

\hypersetup{
  pdftitle={Phase Retrieval for the One-Dimensional Cubic NLS},
  pdfsubject={Expanded self-contained proof and autoformalization specification},
  pdfkeywords={nonlinear Schrodinger equation, phase retrieval, Carleman estimate, unique continuation}
}

\begin{document}
\frontmatter
\hypersetup{pageanchor=false}
\maketitle
\clearpage
\hypersetup{pageanchor=true}
\setcounter{page}{1}

\chapter*{Abstract}
\addcontentsline{toc}{chapter}{Abstract}
We prove the following phase-retrieval statement for the one-dimensional
cubic nonlinear Schr\"odinger equation:
\[
  i\partial_tu+\partial_x^2u=\sigma |u|^2u,
  \qquad \sigma\in\R.
\]
If two global mild solutions $u,v\in C(\R;L^2(\R))$ have the same modulus
almost everywhere on a nonempty open time interval, then
$v=\zeta u$ for one constant $\zeta\in\C$ with $|\zeta|=1$; in particular
the conclusion holds when the moduli agree almost everywhere in all of
spacetime.

This expanded edition is organized as a closed mathematical development and
as a no-invention handoff for Lean autoformalization.  In addition to the
PDE argument, it supplies the measure-theoretic endpoint infrastructure that
was compressed in the earlier draft: jointly measurable scalar definitions
of $L^p_tL^\infty_x$, completeness, restriction, Minkowski integration, and
a countable norming predual; compatibility of the $L^1\to L^\infty$ kernel
propagator with the $L^2$ Fourier group; the exact Fourier--Gagliardo identity;
a separate endpoint $TT^*$ proof; full $H^k$ persistence and the
Gelfand-triple mass chain rule; simultaneous critical/uniform approximation;
an extended quadratic-form test algebra; a local $H^2\Rightarrow C^1$ trace
theorem; the time-integrated Neumann difference-quotient limit; and a
step-by-step graph closure for the explicit Carleman inverse.

The only background layer left implicit is the ordinary foundational
analysis pinned in the formalization contract: scalar Lebesgue and Bochner
integration, complete normed spaces, basic Hilbert-space algebra,
Plancherel, standard approximation, the maximum-modulus principle, and Banach's
fixed-point theorem.  No external Strichartz, smoothing, trace, Carleman, or
unique-continuation theorem is used as a premise of the main result.
\chapter*{How to read the phrase ``modulo Mathlib''}
\addcontentsline{toc}{chapter}{How to read the phrase ``modulo Mathlib''}
This document is a mathematical proof manuscript and an autoformalization
specification; it is not itself a Lean proof term.  The reproducible target is
Lean $4.32.0$ with Mathlib tag \texttt{v4.32.0}, release commit
\texttt{81a5d25}.  A translation may use \texttt{import Mathlib}, but it may
leave implicit only the ordinary foundational layer enumerated precisely in
\cref{app:formalization-contract}: scalar measure theory and integration,
complete normed and Hilbert spaces, elementary Fourier/Plancherel theory,
standard approximation, and Banach's fixed-point theorem.

No Schr\"odinger estimate, nonlinear-flow theorem, critical product rule,
trace theorem, Carleman estimate, or unique-continuation theorem is included
in that imported layer.  Those statements are proved in the body of the
manuscript.  Endpoint mixed norms are defined on jointly measurable scalar
representatives rather than as $L^\infty$-valued Bochner spaces.  The
Carleman graph closure is proved by spatial mollification, almost-everywhere
Fourier sections, and a scalar compact-support ODE.  Distributional traces
are constructed by scalarization and a proved local $H^2$ theorem rather than
by an external trace package.

Mathlib's Fourier character has a different normalization from the unitary
convention used in the displayed analysis.  The exact rescaling is stated in
\eqref{eq:mathlib-fourier-rescaling}; a formalizer may either implement that
rescaling or rewrite all multipliers in Mathlib's convention.  The final
appendix gives the version pin, object representations, exact existing
Fourier/Sobolev entry points, and a no-invention gate for the next stage.

Compiling this TeX file checks typography and internal references only.  A
kernel certificate still requires translating the numbered definitions and
proofs and running Lean with no admitted propositions.  The purpose of the
expanded detail is that the translator should face implementation work, not
an unstated mathematical lemma.

\tableofcontents

\mainmatter

\chapter{The theorem and the proof architecture}

\section{Equation and solution class}
Fix $\sigma\in\R$.  We study
\begin{equation}
  i\partial_tu+\partial_x^2u=\sigma |u|^2u,
  \qquad (t,x)\in\R\times\R.
  \label{eq:NLS-main}
\end{equation}
The free Schr\"odinger group is denoted by
\[
  S(t)=\e^{it\partial_x^2}.
\]
With the Fourier convention fixed in \cref{chap:fourier}, its multiplier is
$\e^{-it\xi^2}$.  A function $u\in C(I;L^2(\R))$ is a mild solution on an
interval $I$ if, for one and hence every $t_0\in I$,
\begin{equation}
 u(t)=S(t-t_0)u(t_0)
 -i\sigma\int_{t_0}^{t}S(t-s)(|u|^2u)(s)\,\dd s
 \quad\text{in }L^2(\R).                         \label{eq:mild-main}
\end{equation}
In the fixed-point construction the retarded term is first defined as the
continuous extension of the smooth-forcing operator from
$L^{6/5}_{t,x}$ to $C_tL^2_x\cap L^6_{t,x}$; see
\cref{lem:retarded-continuous}.  After the solution has been constructed,
\cref{prop:endpoint} proves $|u|^2u\in L^1_tL^2_x$, and the same lemma then
identifies the retarded term with the ordinary $L^2$-valued Bochner integral.
Thus \eqref{eq:mild-main} has its usual Bochner meaning a posteriori, without
using that stronger integrability circularly.  The construction below proves
that every $u_0\in L^2$ generates a unique global mild solution and that on
every bounded interval
\begin{equation}
 u\in L^6_{t,x}\cap L^4_tL^\infty_x\cap L^8_tL^4_x. \label{eq:solution-class-preview}
\end{equation}

\begin{theorem}[Main phase-retrieval theorem]\label{thm:main}
Let $u,v\in C(\R;L^2(\R))$ be the global mild solutions of
\eqref{eq:NLS-main} with initial data $u_0,v_0\in L^2(\R)$.  Assume
\begin{equation}
  |u(t,x)|=|v(t,x)|
  \quad\text{for almost every }(t,x)\in\R^2.       \label{eq:equal-modulus-main}
\end{equation}
Then there is a scalar $\zeta\in\C$ such that $|\zeta|=1$ and
\[
  v(t)=\zeta u(t)\quad\text{in }L^2(\R)
  \quad\text{for every }t\in\R.
\]
Equivalently, after choosing measurable representatives,
$v(t,x)=\zeta u(t,x)$ for almost every $(t,x)$.
\end{theorem}

\section{The chain of implications}
The proof is easier to audit when separated into seven logically independent
links.

\begin{enumerate}[label=\textbf{Step \arabic*.},leftmargin=5.2em]
\item Construct the global $L^2$ flow and establish the local spacetime norms
      in \eqref{eq:solution-class-preview}.
\item Use local smoothing to make the current
      $\Ima(\overline u\,u_x)$ meaningful as an $H^{-1/2}$--$H^{1/2}$
      quadratic form.  The continuity equation then shows that equal density
      implies equal current.
\item Prove a critical Wronskian lemma which does not divide by $u$ or $v$ at
      their common zero set.  It yields
      $uv_x-vu_x=0$ for almost every time.
\item Form the exterior product
      $F(t,x,y)=u(t,x)v(t,y)-v(t,x)u(t,y)$.  It solves a two-dimensional
      linear Schr\"odinger equation with a real ridge potential.
\item Show that $F$ has zero Dirichlet and Neumann traces on the diagonal
      $x=y$, even though the traces are only distribution-valued a priori.
\item Extend the restriction of $F$ from one side of the diagonal by zero.
      The two vanishing traces remove both jump distributions.  A
      half-space propagation theorem forces the extension, hence $F$, to
      vanish.
\item A zero Hilbert-space exterior product means that $u(0)$ and $v(0)$ are
      linearly dependent.  Equal norms make the scalar unimodular, and gauge
      covariance propagates the relation to all times.
\end{enumerate}

Every step is stable under smooth approximation.  This stability is not a
cosmetic point: it is what turns identities first computed for Schwartz data
into identities in the global $L^2$ solution class.

\section{Why no pointwise division is used}
At a fixed time, equality of moduli suggests writing $v=q u$ with $|q|=1$.
This is harmless where $u$ is bounded away from zero, but it loses information
at nodal points.  Different connected components of $\{u\ne0\}$ may carry
different constant phases.  For example, the pair $f(x)=x$ and $g(x)=|x|$
shows that a phase jump across a zero is compatible with equal modulus and a
vanishing classical Wronskian away from the zero.  The proof must therefore
use the evolution to glue phases.  The exterior product carries all component
phases at once and is the natural quantity for that task.

\begin{checkpoint}
The main theorem will use no statement not proved later in the document,
except the foundational analysis listed in the front matter.
\end{checkpoint}

\chapter{Measurable representatives and mixed norms}\label{chap:mixed-norms}

The endpoint space $L^4_tL^\infty_x$ occurs in the nonlinear flow and again
in the Carleman argument.  It is not treated here as a Bochner space with
values in the nonseparable Banach space $L^\infty$.  Instead, every mixed norm
is defined directly on jointly measurable scalar representatives.  This
removes all ambiguity about sections, essential suprema, restriction, and
Minkowski integration.

\section{Joint representatives and section norms}

\begin{convention}[Spacetime representatives]\label{conv:joint-representatives}
A scalar spacetime function on $I\times\R^d$ means an equivalence class of
jointly Lebesgue-measurable maps $f:I\times\R^d\to\C$, where two maps are
identified when they agree almost everywhere for product measure.  A
representative is fixed only for the duration of a calculation.  Fubini's
theorem implies that two representatives determine the same spatial section
for almost every time, so all section norms below are independent of the
choice of representative.
\end{convention}

For $1\le q<\infty$ put
\[
 m_qf(t)=\left(\int_{\R^d}|f(t,x)|^q\,\dd x\right)^{1/q},
\]
and put
\[
 m_\infty f(t)=\esssup_{x\in\R^d}|f(t,x)|.
\]

\begin{lemma}[Measurability of the endpoint section norm]
\label{lem:measurable-esssup-section}
If $f$ is jointly measurable, then $t\mapsto m_\infty f(t)$ is measurable.
\end{lemma}

\begin{proof}
Choose once and for all a strictly positive integrable weight, for example
$\omega(x)=c_d\exp(-|x|)$ with $\int\omega=1$.  For every rational $a\ge0$,
Tonelli's theorem shows that
\[
 t\longmapsto A_a(t)
 :=\int_{\R^d}\omega(x)\1_{\{|f(t,x)|>a\}}\,\dd x
\]
is measurable.  Since $\omega>0$ almost everywhere,
\[
 m_\infty f(t)>a \quad\Longleftrightarrow\quad A_a(t)>0.
\]
Thus the inverse image of $(a,\infty]$ is measurable for every rational
$a\ge0$, which proves the claim.
\end{proof}

\begin{definition}[Scalar mixed norm]\label{def:scalar-mixed-norm}
For $1\le p\le\infty$ and $1\le q\le\infty$, define
\[
 \norm f_{L^p_tL^q_x(I\times\R^d)}
 :=\norm{m_qf}_{L^p(I)}.
\]
The resulting space is the quotient of the jointly measurable functions of
finite norm by the zero-norm relation.  If $E\subset I\times\R^d$ is
measurable, the restricted norm means the norm of $\1_Ef$.
\end{definition}

For finite $q$, measurability of $m_qf$ follows directly from Tonelli.  The
previous lemma supplies the only endpoint case.  Fubini also shows that the
zero-norm relation is exactly equality almost everywhere on the product
space.

\begin{lemma}[Completeness of the scalar endpoint space]
\label{lem:scalar-endpoint-complete}
For $1\le p\le\infty$, the scalar space $L^p_tL^\infty_x$ of
\cref{def:scalar-mixed-norm} is complete.
\end{lemma}

\begin{proof}
Let $(f_n)$ be Cauchy.  Choose a subsequence $(f_{n_k})$ such that
\[
 \norm{f_{n_{k+1}}-f_{n_k}}_{L^p_tL^\infty_x}\le2^{-k}.
\]
Set
\[
 a_k(t)=m_\infty(f_{n_{k+1}}-f_{n_k})(t).
\]
Minkowski's inequality for the scalar $L^p(I)$ norm gives
\[
 \norm{\sum_{k=1}^\infty a_k}_{L^p_t}
 \le\sum_{k=1}^\infty\norm{a_k}_{L^p_t}<\infty,
\]
with the same statement in the essential-supremum sense when $p=\infty$.
Hence $\sum_ka_k(t)<\infty$ for almost every $t$.

For every $k$, the inequality
\[
 |f_{n_{k+1}}(t,x)-f_{n_k}(t,x)|\le a_k(t)
\]
holds for almost every $(t,x)$.  After deleting the countable union of the
exceptional sets, the telescoping series converges absolutely for almost
every $(t,x)$.  Define $f$ to be its sum there and zero elsewhere.  It is
jointly measurable.  Moreover,
\[
 m_\infty(f-f_{n_k})(t)\le\sum_{j\ge k}a_j(t)
\]
for almost every $t$, so $f_{n_k}\to f$ in the mixed norm.  Since the
original sequence is Cauchy, the whole sequence converges to $f$.
\end{proof}

\begin{lemma}[Restriction and scalar Minkowski]
\label{lem:scalar-mixed-minkowski}
Let $(S,\nu)$ be a $\sigma$-finite measure space and let
$F:S\times I\times\R^d\to\C$ be jointly measurable.  Assume
\[
 \int_S\norm{F(s,\cdot,\cdot)}_{L^p_tL^\infty_x}\,\dd\nu(s)<\infty.
\]
Then the scalar integral $H(t,x)=\int_SF(s,t,x)\,\dd\nu(s)$ is defined for
almost every $(t,x)$ and
\begin{equation}
 \norm H_{L^p_tL^\infty_x}
 \le\int_S\norm{F(s)}_{L^p_tL^\infty_x}\,\dd\nu(s).
 \label{eq:scalar-minkowski-endpoint}
\end{equation}
Multiplication by an indicator of a measurable subset has norm at most one in
every scalar mixed norm.
\end{lemma}

\begin{proof}
Put $B(t)=\int_Sm_\infty(F(s))(t)\,\dd\nu(s)$.  Ordinary scalar
Minkowski gives $B\in L^p(I)$.  After removing one product-null set by
Tonelli, $|F(s,t,x)|\le m_\infty(F(s))(t)$ for almost every
$(s,t,x)$; hence the scalar integral defining $H$ is absolutely convergent
for almost every $(t,x)$.  For almost every $t$, the definition of essential
supremum then gives
\[
 m_\infty H(t)
 \le\int_Sm_\infty(F(s))(t)\,\dd\nu(s)=B(t).
\]
Taking the scalar $L^p(I)$ norm proves \eqref{eq:scalar-minkowski-endpoint}.  The restriction statement is pointwise.
\end{proof}

\section{A norming predual at the spatial endpoint}

We never identify the full Banach dual of $L^4_tL^\infty_x$.  Only a
norming family of $L^{4/3}_tL^1_x$ functions is needed, and that statement has
an elementary measurable proof.

Choose a countable dense subset $(g_n)_{n\ge1}$ of the closed unit ball of
$L^1(\R^d)$.  One may take compactly supported step functions on rational
boxes with coefficients in $\Q+i\Q$, followed by normalization.

\begin{lemma}[Countable norming of $L^\infty$]\label{lem:countable-norming-Linfty}
For every $h\in L^\infty(\R^d)$,
\begin{equation}
 \norm h_\infty
 =\sup_{n\ge1}\left|\int_{\R^d}h(x)g_n(x)\,\dd x\right|.
 \label{eq:countable-norming-Linfty}
\end{equation}
\end{lemma}

\begin{proof}
The right side is at most $\norm h_\infty$.  Conversely, let
$M=\norm h_\infty$ and $\eps>0$.  The set
$E=\{|h|>M-\eps\}$ has positive measure.  Intersect it with a sufficiently
large ball and, if necessary, with a subset of finite positive measure.  The
function
\[
 g=\frac{\1_E\overline{h}/|h|}{|E|},
\]
with value zero where $h=0$, belongs to the unit ball of $L^1$ and satisfies
$|\int hg|\ge M-\eps$.  Approximate $g$ in $L^1$ by one of the $g_n$ and let
$\eps\downarrow0$.
\end{proof}

\begin{proposition}[Endpoint norming identity]
\label{prop:endpoint-norming}
Let $1<p<\infty$ and $p'=p/(p-1)$.  For every jointly measurable
$H\in L^p_tL^\infty_x$,
\begin{equation}
 \norm H_{L^p_tL^\infty_x}
 =\sup\left\{
 \left|\iint H(t,x)G(t,x)\,\dd x\,\dd t\right|:
 \norm G_{L^{p'}_tL^1_x}\le1
 \right\}.
 \label{eq:endpoint-norming}
\end{equation}
It is enough in the supremum to use jointly measurable functions which, at
each time, are a scalar multiple of one of the countably many $g_n$.
\end{proposition}

\begin{proof}
H\"older gives the upper bound.  Put
\[
 a_n(t)=\int H(t,x)g_n(x)\,\dd x,
 \qquad M_N(t)=\max_{1\le n\le N}|a_n(t)|.
\]
Every $a_n$ is measurable by Fubini, and by
\cref{lem:countable-norming-Linfty},
$M_N(t)\uparrow m_\infty H(t)$ for almost every $t$.
Let $n_N(t)$ be the least index realizing the finite maximum; its level sets
are measurable.  Define
\[
 \theta_N(t)=
 \begin{cases}
  \overline{a_{n_N(t)}(t)}/|a_{n_N(t)}(t)|,&M_N(t)>0,\\
  0,&M_N(t)=0.
 \end{cases}
\]
If $M_N\not\equiv0$, put
\[
 b_N(t)=\frac{M_N(t)^{p-1}}{\norm{M_N}_{L^p}^{p-1}},
 \qquad
 G_N(t,x)=b_N(t)\theta_N(t)g_{n_N(t)}(x).
\]
Then $G_N$ is jointly measurable,
$\norm{G_N}_{L^{p'}_tL^1_x}\le1$, and
\[
 \iint HG_N=\int b_N(t)M_N(t)\,\dd t=\norm{M_N}_{L^p}.
\]
The case $M_N\equiv0$ is trivial.  Monotone convergence gives
$\norm{M_N}_p\uparrow\norm H_{L^pL^\infty}$, proving the reverse inequality.
\end{proof}

\begin{remark}
\Cref{prop:endpoint-norming} is a norming statement, not a claim that every
continuous functional on $L^{p'}_tL^1_x$ is represented by a strongly
measurable $L^\infty$-valued map.  This distinction is exactly what is needed
at the endpoint Strichartz step.
\end{remark}

\section{Almost-everywhere uniform convergence of continuous slices}

\begin{lemma}[Mixed-norm subsequence lemma]
\label{lem:mixed-subsequence-uniform}
Let $1\le p<\infty$ and suppose $f_n\to f$ in
$L^p(I;L^\infty(\R^d))$ in the scalar sense.  There is a subsequence
$(f_{n_k})$ such that for almost every $t$,
\[
 \sum_{k=1}^\infty
 \norm{f_{n_{k+1}}(t,\cdot)-f_{n_k}(t,\cdot)}_{L^\infty_x}<\infty.
\]
If every spatial slice $f_n(t,\cdot)$ is continuous, then the essential
supremum in the series equals the ordinary supremum and the subsequence is
uniformly Cauchy in $x$ for almost every $t$.
\end{lemma}

\begin{proof}
Choose the subsequence so that the $k$th mixed-norm difference is at most
$2^{-k}$.  The proof of \cref{lem:scalar-endpoint-complete} shows that the
sum of the section essential suprema is finite for almost every $t$.  A
continuous function whose absolute value exceeds a number at one point
exceeds it on a set of positive measure, so its supremum and essential
supremum agree.
\end{proof}

\begin{lemma}[Identification of the continuous limit]
\label{lem:C0-limit-identification}
Suppose, in addition, that $f_n(t,\cdot)\in C_0(\R^d)$, that
$f_n\to f$ in $C(I;L^2_x)$, and that a subsequence converges uniformly in
$x$ for a given time $t$.  Then its uniform limit belongs to $C_0$ and
represents the $L^2$ class $f(t)$.
\end{lemma}

\begin{proof}
The space $C_0$ is closed in the uniform norm.  The same subsequence converges
to $f(t)$ in $L^2$, while uniform convergence implies convergence locally in
$L^2$.  Uniqueness of an $L^2$ limit on every bounded set identifies the two
limits almost everywhere; the continuous uniform limit is therefore a
representative of $f(t)$.
\end{proof}

\section{Sum and intersection spaces}

The Carleman argument uses sums and intersections of scalar mixed-norm
spaces.  If $A$ and $B$ are Banach spaces continuously embedded in the same
Hausdorff topological vector space, define
\[
 A+B=\{a+b:a\in A,\ b\in B\},
 \qquad
 \norm h_{A+B}=\inf_{h=a+b}(\norm a_A+\norm b_B),
\]
and
\[
 A\cap B=\{h:h\in A\text{ and }h\in B\},
 \qquad
 \norm h_{A\cap B}=\max(\norm h_A,\norm h_B).
\]

\begin{lemma}[Completeness of sums and intersections]
\label{lem:sum-intersection-complete}
With the preceding hypotheses, $A+B$ and $A\cap B$ are Banach spaces.
Multiplication by an indicator is contractive whenever it is contractive on
each component.
\end{lemma}

\begin{proof}
The sum is the quotient of the Banach space $A\oplus_1B$ by the closed kernel
of the continuous addition map $(a,b)\mapsto a+b$.  The intersection is
isometric to the closed diagonal subspace
$\{(a,b)\in A\oplus_\infty B:a=b\}$.  The restriction statement follows by
applying the componentwise contractions before taking an infimum or a
maximum.
\end{proof}

\section{Formal meaning of the mild equation}

\begin{lemma}[Joint representative for a continuous $L^2$ path]
\label{lem:joint-representative-L2}
Let $I$ be an interval and let $u:I\to L^2(\R^d)$ be continuous.  There is a
jointly measurable scalar function $U:I\times\R^d\to\C$ such that, for almost
every $t$, the section $U(t,\cdot)$ represents the $L^2$ class $u(t)$.  On a
bounded interval $I$, the representative can be obtained as an
$L^2(I\times\R^d)$ limit of jointly measurable simple-in-time functions whose
sections converge to $u(t)$ uniformly in the $L^2$ norm.
\end{lemma}

\begin{proof}
First let $I=[a,b]$.  Uniform continuity of $u$ gives, for every $n$, a
finite measurable partition $I=\bigcup_jE_{n,j}$ and points $t_{n,j}\in I$
such that the simple $L^2$-valued map
\[
 u_n(t)=\sum_j\1_{E_{n,j}}(t)u(t_{n,j})
\]
satisfies $\sup_t\norm{u_n(t)-u(t)}_2\le2^{-n}$.  Choose a measurable scalar
representative $g_{n,j}$ of each of the finitely many $L^2$ classes
$u(t_{n,j})$, and put
\[
 U_n(t,x)=\sum_j\1_{E_{n,j}}(t)g_{n,j}(x).
\]
This is jointly measurable and represents $u_n(t)$ for every section.  The
uniform estimate implies
\[
 \norm{U_n-U_m}_{L^2(I\times\R^d)}
 \le |I|^{1/2}(2^{-n}+2^{-m}),
\]
so $U_n$ converges in product $L^2$ to a jointly measurable function $U$.
Choose a subsequence, still denoted $U_n$, whose successive product-$L^2$
differences have summable norms.  Minkowski then shows that the sum of the
section $L^2_x$ norms of those differences is finite for almost every $t$.
Thus $U_n(t,\cdot)\to U(t,\cdot)$ in $L^2_x$ for almost every $t$.  For every
$t$, however, the classes represented by $U_n(t,\cdot)$ converge in $L^2$ to
$u(t)$.  Uniqueness of the $L^2$ limit therefore identifies the section of
$U$ with $u(t)$ outside one null set of times.

For an arbitrary interval, apply the construction on a countable increasing
exhaustion by bounded intervals.  On overlaps the representatives agree
product-almost everywhere because their sections represent the same $L^2$
classes.  Modify on the countable union of overlap null sets and paste the
representatives along disjoint annular pieces of the exhaustion.
\end{proof}

A map $u:I\to L^2(\R)$ is always understood as a continuous map into the
Mathlib $L^2$ quotient.  We use the jointly measurable representative
supplied by \cref{lem:joint-representative-L2}.  If
$u\in L^6(I\times\R)$, the scalar function $|u|^2u$ is independent of the
chosen representative and defines an element of $L^{6/5}(I\times\R)$ by the
estimate proved later.  The retarded Duhamel term is first defined for smooth compactly supported
forcing and then completed in the norm supplied by
\cref{thm:inhom-str}; \cref{lem:retarded-continuous} proves that the completed
object has a unique continuous $L^2$ representative, is independent of the
approximating sequence, and satisfies the base-time identity.  Consequently
\eqref{eq:mild-main} is an equality in the $L^2$ quotient for every time; no
pointwise representative is part of the definition.  When the forcing also
belongs to $L^1_tL^2_x$, the completed operator agrees with the ordinary
Bochner integral.

\begin{checkpoint}
All uses of $L^4_tL^\infty_x$, endpoint duality, restriction, Minkowski
integration, and almost-everywhere uniform convergence now refer to explicit
scalar constructions.  No nonseparable Bochner-space identification is used
anywhere below.
\end{checkpoint}

\chapter{Fourier analysis and Sobolev spaces}\label{chap:fourier}

\section{Fourier convention}
For $f\in L^1(\R^d)$ define
\begin{equation}
 \widehat f(\xi)=(2\pi)^{-d/2}\int_{\R^d}\e^{-ix\cdot\xi}f(x)\,\dd x.
 \label{eq:FT-convention}
\end{equation}
The inverse transform uses $\e^{ix\cdot\xi}$ with the same normalization.
The transform extends by density to a unitary operator on $L^2(\R^d)$.
We use Plancherel in the form
\[
  \norm{f}_{L^2}=\norm{\widehat f}_{L^2},
  \qquad
  \pair{f}{g}_{L^2}=\pair{\widehat f}{\widehat g}_{L^2}.
\]
For Schwartz functions,
$\widehat{\partial_jf}=i\xi_j\widehat f$ and
$\widehat{\Delta f}=-|\xi|^2\widehat f$.

\section{Bessel and homogeneous multipliers}
For $s\in\R$ let
\[
  \langle D\rangle^sf
  =\FT^{-1}(\langle\xi\rangle^s\widehat f),
  \qquad
  |D|^sf=\FT^{-1}(|\xi|^s\widehat f),
  \qquad \langle\xi\rangle=(1+|\xi|^2)^{1/2}.
\]
Define
\[
  H^s(\R^d)=\{f\in\Ssch'(\R^d):\langle\xi\rangle^s\widehat f\in L^2\},
  \qquad
  \norm{f}_{H^s}=\norm{\langle D\rangle^sf}_2.
\]
Only $s=\pm\tfrac12$ and $s=1$ are needed.  Local membership means
$\chi f\in H^s$ for every $\chi\in C_c^\infty$.

For $0<s<1$ set
\[
 c_s=\int_{\R}\frac{|\e^{ih}-1|^2}{|h|^{1+2s}}\,\dd h.
\]
The integral is finite and strictly positive: near zero the numerator is
$O(h^2)$, while at infinity it is bounded.

\begin{theorem}[Exact Fourier--Gagliardo identity]
\label{thm:fourier-gagliardo}
For $0<s<1$ and every $f\in L^2(\R)$, the following are equivalent:
\begin{enumerate}
\item $|\xi|^s\widehat f(\xi)\in L^2_\xi$;
\item
\[
 \iint_{\R^2}\frac{|f(x)-f(y)|^2}{|x-y|^{1+2s}}\,\dd x\,\dd y<\infty.
\]
\end{enumerate}
Whenever either condition holds,
\begin{equation}
 \iint_{\R^2}\frac{|f(x)-f(y)|^2}{|x-y|^{1+2s}}\,\dd x\,\dd y
 =c_s\int_\R|\xi|^{2s}|\widehat f(\xi)|^2\,\dd\xi.
 \label{eq:exact-gagliardo-fourier}
\end{equation}
In particular, the norm obtained from the $L^2$ norm and the Gagliardo
seminorm is equivalent to the Bessel-potential $H^s$ norm.
\end{theorem}

\begin{proof}
First let $f$ be Schwartz.  Put $h=y-x$.  Tonelli and Plancherel give
\begin{align*}
 \iint\frac{|f(x+h)-f(x)|^2}{|h|^{1+2s}}\,\dd x\,\dd h
 &=\int_\R\frac1{|h|^{1+2s}}
   \int_\R|\e^{ih\xi}-1|^2|\widehat f(\xi)|^2\,\dd\xi\,\dd h\\
 &=\int_\R|\widehat f(\xi)|^2
   \left(\int_\R\frac{|\e^{ih\xi}-1|^2}{|h|^{1+2s}}\,\dd h\right)\dd\xi.
\end{align*}
For $\xi\ne0$, the substitution $k=h\xi$ makes the inner integral equal to
$c_s|\xi|^{2s}$; it is zero at $\xi=0$.  This proves
\eqref{eq:exact-gagliardo-fourier} for Schwartz functions.

Now suppose $|\xi|^s\widehat f\in L^2$.  Approximate $\widehat f$ in the
weighted norm
\[
 \int(1+|\xi|^{2s})|\widehat f(\xi)|^2\,\dd\xi
\]
by functions in $C_c^\infty$ and take inverse Fourier transforms $f_n$.
Then $f_n$ is Schwartz, $f_n\to f$ in $L^2$, and
$|D|^s(f_n-f_m)\to0$ in $L^2$.  The identity already proved shows that
\[
 \frac{f_n(x)-f_n(y)}{|x-y|^{(1+2s)/2}}
\]
is Cauchy in $L^2(\R^2)$.  After passing to a subsequence, $f_n\to f$ almost
everywhere, so the $L^2(\R^2)$ limit is the displayed difference quotient
of $f$.  Passing to the limit proves the identity.

Conversely, suppose the Gagliardo integral of $f$ is finite.  Let
$f_\eps=\varphi_\eps*f$ with a compactly supported smooth approximate
identity.  Jensen's inequality and translation invariance give
\[
 [f_\eps]_{\dot H^s}\le[f]_{\dot H^s}.
\]
We now justify the Schwartz approximation in both norms rather than taking it
as a density theorem.  Choose $\chi\in C_c^\infty(\R)$ with
$0\le\chi\le1$ and $\chi=1$ near zero, and put
$\chi_R(x)=\chi(x/R)$ and $g_R=\chi_Rf_\eps$.  Then
$g_R\in C_c^\infty(\R)\subset\mathcal S(\R)$ and
$g_R\to f_\eps$ in $L^2$.  For the Gagliardo seminorm, write
\begin{align*}
 g_R(x)-g_R(y)
 &=\chi_R(x)(f_\eps(x)-f_\eps(y))
   +f_\eps(y)(\chi_R(x)-\chi_R(y)).
\end{align*}
The first summand converges to $f_\eps(x)-f_\eps(y)$ in the weighted
$L^2(\R^2)$ space by dominated convergence.  For the second, scaling and
the Lipschitz bound for $\chi$ give, uniformly in $y$,
\begin{align}
 \int_\R\frac{|\chi_R(x)-\chi_R(y)|^2}{|x-y|^{1+2s}}\,\dd x
 &=R^{-2s}\int_\R
   \frac{|\chi(y/R+z)-\chi(y/R)|^2}{|z|^{1+2s}}\,\dd z\notag\\
 &\le C_{\chi,s}R^{-2s}.                         \label{eq:cutoff-gagliardo-error}
\end{align}
Indeed, on $|z|\le1$ the numerator is $O(|z|^2)$, and on
$|z|>1$ it is bounded.  Therefore the squared weighted norm of the second
summand is at most
$C_{\chi,s}R^{-2s}\norm{f_\eps}_2^2$, which tends to zero.  Hence
$g_R\to f_\eps$ in $L^2$ and in the Gagliardo seminorm.

Apply the Schwartz identity to $g_R-g_{R'}$.  The preceding convergence shows
that $|\xi|^s\widehat{g_R}$ is Cauchy in $L^2_\xi$, while ordinary
Plancherel gives $\widehat{g_R}\to\widehat{f_\eps}$ in $L^2$.  On every
annulus $\{a\le|\xi|\le b\}$ the two limits identify, so the weighted
limit is $|\xi|^s\widehat{f_\eps}$.  Passing to the limit in the Schwartz
identity yields
\[
 c_s\int|\xi|^{2s}|\widehat{f_\eps}(\xi)|^2\,\dd\xi
 =[f_\eps]_{\dot H^s}^2
 \le[f]_{\dot H^s}^2.
\]
With the unitary convention \eqref{eq:FT-convention}, the convolution
formula contains the factor $(2\pi)^{1/2}$.  Thus
\[
 \widehat{f_\eps}(\xi)=m_\eps(\xi)\widehat f(\xi),
 \qquad
 m_\eps(\xi)=(2\pi)^{1/2}\widehat\varphi(\eps\xi).
\]
Because $\int\varphi=1$, one has $m_\eps(\xi)\to1$ pointwise.
Fatou's lemma therefore yields
$|\xi|^s\widehat f\in L^2$ with the same bound.  The identity for $f$
then follows from the already proved implication from the weighted Fourier
condition to the Gagliardo condition.
\end{proof}

\begin{remark}[Normalization relative to Mathlib]\label{rem:mathlib-fourier-normalization}
Mathlib's standard Fourier character uses the phase $\e^{-2\pi i x\cdot\eta}$.
If $\FT_{\mathrm M}$ denotes that transform and $\FT$ denotes
\eqref{eq:FT-convention}, then in dimension $d$
\begin{equation}
 (\FT f)(\xi)=(2\pi)^{-d/2}
 (\FT_{\mathrm M}f)\!\left(\frac{\xi}{2\pi}\right).
 \label{eq:mathlib-fourier-rescaling}
\end{equation}
The change of variables $\xi=2\pi\eta$ shows that the right side is unitary
on $L^2$.  A Lean implementation may either define this rescaled transform or
retain Mathlib's convention and replace every occurrence of $\xi$ by
$2\pi\eta$.  No analytic estimate depends on which choice is made.
\end{remark}

\begin{lemma}[Smooth multipliers on $H^{1/2}$]\label{lem:smooth-multiplier}
If $\chi\in C_c^\infty(\R)$, multiplication by $\chi$ is bounded on
$H^{1/2}(\R)$ and on $H^{-1/2}(\R)$.
\end{lemma}

\begin{proof}
For $H^{1/2}$, the $L^2$ term is immediate.  From
\[
 \chi(x)f(x)-\chi(y)f(y)
 =\chi(x)(f(x)-f(y))+f(y)(\chi(x)-\chi(y))
\]
and $|\chi(x)-\chi(y)|\le \norm{\chi'}_\infty|x-y|$, the first term is
controlled by $\norm{\chi}_\infty[f]_{\dot H^{1/2}}$.  In the second term,
the region $|x-y|\le1$ is controlled by $\norm{\chi'}_\infty\norm f_2$ and
the region $|x-y|>1$ by boundedness and compact support of $\chi$.  The
$H^{-1/2}$ statement follows by duality.
\end{proof}

\begin{lemma}[Translation difference quotients]\label{lem:translation-dq}
For $f\in H^{1/2}(\R)$,
\[
 \frac{f(\cdot+h)-f}{h}\longrightarrow f_x
 \quad\text{strongly in }H^{-1/2}(\R)
 \quad(h\to0).
\]
Moreover the difference quotients are uniformly bounded from
$H^{1/2}$ to $H^{-1/2}$.
\end{lemma}

\begin{proof}
The Fourier multiplier of the difference quotient is
$(\e^{ih\xi}-1)/h$.  After multiplying by $\langle\xi\rangle^{-1/2}$, its
absolute value is bounded by $C\langle\xi\rangle^{1/2}$.  It converges
pointwise to $i\xi$.  Dominated convergence in $L^2_\xi$ proves the claim.
\end{proof}

\begin{lemma}[Difference quotients in $W^{1,p}$]
\label{lem:Wp-difference-quotient}
Let $1\le p<\infty$ and $f\in W^{1,p}(\R)$.  Then
\begin{equation}
 \frac{f(\cdot+h)-f}{h}\longrightarrow f'
 \quad\text{strongly in }L^p(\R)\quad(h\to0),       \label{eq:Wp-difference-quotient}
\end{equation}
and the difference quotients have norm at most $\norm{f'}_p$.
The same assertion holds after integration in an auxiliary variable.
\end{lemma}

\begin{proof}
Use the locally absolutely continuous representative of $f$.  For almost
every $x$,
\[
 \frac{f(x+h)-f(x)}h=\int_0^1 f'(x+\theta h)\,\dd\theta.
\]
Minkowski gives the uniform bound.  Subtract $f'(x)$, apply Minkowski again,
and use strong continuity of translations in scalar $L^p$:
\[
 \left\|\frac{\tau_hf-f}{h}-f'\right\|_p
 \le\int_0^1\norm{\tau_{\theta h}f'-f'}_p\,\dd\theta\longrightarrow0
\]
by dominated convergence in $\theta$.  For an auxiliary variable, apply the
same identity to almost every section and use the identical domination in
the product $L^p$ norm.
\end{proof}

\begin{lemma}[One-dimensional $H^1$ representative]
\label{lem:H1-C0}
Every $f\in H^1(\R)$ has a unique continuous representative in $C_0(\R)$.
It is locally absolutely continuous and satisfies
\begin{equation}
 \norm f_\infty^2\le2\norm f_2\norm{f'}_2.          \label{eq:H1-Linfty}
\end{equation}
More precisely, for every $x$,
\begin{align}
 |f(x)|^2&\le2\norm f_{L^2(x,\infty)}\norm{f'}_{L^2(x,\infty)},\label{eq:H1-tail-plus}\\
 |f(x)|^2&\le2\norm f_{L^2(-\infty,x)}\norm{f'}_{L^2(-\infty,x)}.\label{eq:H1-tail-minus}
\end{align}
\end{lemma}

\begin{proof}
The one-dimensional weak fundamental theorem gives a locally absolutely
continuous representative.  Since $f\in L^2$, the numbers
$\int_n^{n+1}|f|^2$ tend to zero.  Choose a Lebesgue point
$R_n\in[n,n+1]$ with
$|f(R_n)|^2\le\int_n^{n+1}|f|^2$; then $f(R_n)\to0$.  The same construction
on $[-n-1,-n]$ gives $L_n\to-\infty$ with $f(L_n)\to0$.  For fixed $x$
and $R_n>x$, the ordinary product rule gives
\[
 |f(R_n)|^2-|f(x)|^2
 =2\Rea\int_x^{R_n}f'(y)\overline{f(y)}\,\dd y.
\]
Let $n\to\infty$ and apply Cauchy--Schwarz to obtain
\eqref{eq:H1-tail-plus}.  Integrating from $L_n$ to $x$ gives
\eqref{eq:H1-tail-minus}; either estimate implies \eqref{eq:H1-Linfty}.
The tail norms in \eqref{eq:H1-tail-plus} and
\eqref{eq:H1-tail-minus} tend to zero at the corresponding infinity, so the
representative lies in $C_0$.  Two continuous representatives agreeing
almost everywhere agree everywhere.
\end{proof}

\section{Tensor products}
For $f,g\in L^2(\R)$, the function $(f\otimes g)(x,y)=f(x)g(y)$ belongs to
$L^2(\R^2)$ and
\begin{equation}
  \norm{f\otimes g}_{L^2(\R^2)}=\norm f_2\norm g_2. \label{eq:tensor-norm}
\end{equation}
This is Tonelli followed by factorization of the integral.  Consequently the
map $(f,g)\mapsto f\otimes g$ is jointly continuous from $L^2\times L^2$ to
$L^2(\R^2)$ on bounded sets.  The orthogonal coordinate change
\[
 r=\frac{x+y}{\sqrt2},\qquad s=\frac{x-y}{\sqrt2}
\]
preserves Lebesgue measure and the Laplacian:
$\partial_x^2+\partial_y^2=\partial_r^2+\partial_s^2$.

\chapter{The free Schr\"odinger group and Strichartz estimates}

\section{The propagator}
For $f\in L^2(\R)$ define
\begin{equation}
 \widehat{S(t)f}(\xi)=\e^{-it\xi^2}\widehat f(\xi). \label{eq:free-group}
\end{equation}
Plancherel gives $\norm{S(t)f}_2=\norm f_2$, the group law, and strong
continuity in $t$.  For $f\in\Ssch$, differentiation under the Fourier
integral gives $(i\partial_t+\partial_x^2)S(t)f=0$.


\begin{lemma}[Complex Gaussian and the free kernel]
\label{lem:complex-gaussian-kernel}
Let $z\in\C$ satisfy $\Rea z>0$, and choose the holomorphic square root on
the right half-plane which is positive on $(0,\infty)$.  For every
$a\in\R$,
\begin{equation}
 \int_\R \e^{-z\xi^2+ia\xi}\,\dd\xi
 =\sqrt\pi\,z^{-1/2}\exp\!\left(-\frac{a^2}{4z}\right).
 \label{eq:complex-gaussian}
\end{equation}
Consequently, for $t\ne0$ and $f\in\mathcal S(\R)$,
\begin{equation}
 S(t)f(x)=(4\pi it)^{-1/2}\int_\R
 \exp\!\left(\frac{i|x-y|^2}{4t}\right)f(y)\,\dd y,
 \label{eq:free-kernel-explicit}
\end{equation}
where the square root is the boundary value from $\Rea z>0$.
\end{lemma}

\begin{proof}
For $\Rea z>0$, the integral
$I(z,0)=\int_\R\e^{-z\xi^2}\,\dd\xi$ is holomorphic in $z$: on every
compact subset of the right half-plane, the integrand and its $z$ derivative
are dominated by a polynomial times $\e^{-c\xi^2}$.  For positive real $z$,
the substitution $y=\sqrt z\,\xi$ and the real Gaussian integral give
$I(z,0)=\sqrt\pi z^{-1/2}$.  The identity theorem extends this equality to
the whole right half-plane.

Fix such a $z$ and put
$I_z(a)=\int\e^{-z\xi^2+ia\xi}\,\dd\xi$.  Differentiation under the integral
is legitimate.  Integration by parts, with vanishing boundary terms because
$\Rea z>0$, gives
\[
 0=\int_\R(-2z\xi+ia)\e^{-z\xi^2+ia\xi}\,\dd\xi,
 \qquad
 I_z'(a)=i\int_\R\xi\e^{-z\xi^2+ia\xi}\,\dd\xi
        =-\frac a{2z}I_z(a).
\]
Solving this scalar ODE and using the value at $a=0$ proves
\eqref{eq:complex-gaussian}.

Now take $z=\eps+it$ with $\eps>0$.  Insert the Fourier definition of
$\widehat f$, use Fubini (the Gaussian damping makes the double integral
absolutely convergent), and apply \eqref{eq:complex-gaussian} with
$a=x-y$.  This gives
\[
 (2\pi)^{-1}\iint
  \e^{i(x-y)\xi-(\eps+it)\xi^2}f(y)\,\dd y\,\dd\xi
 =(4\pi(\eps+it))^{-1/2}\int
  \exp\!\left(-\frac{|x-y|^2}{4(\eps+it)}\right)f(y)\,\dd y.
\]
On the Fourier side, dominated convergence against
$|\widehat f|\in L^1$ gives the inverse transform defining $S(t)f$, uniformly
in $x$.  On the physical side, the exponential has modulus at most one and
the prefactor stays bounded as $\eps\downarrow0$; dominated convergence
against $|f(y)|$ therefore gives the right side of
\eqref{eq:free-kernel-explicit}.  The boundary value of the chosen square
root fixes the phase for both signs of $t$.
\end{proof}

\begin{lemma}[Dispersive estimate]\label{lem:dispersive-1d}
For $t\ne0$ and $f\in L^1(\R)$,
\[
  \norm{S(t)f}_{L^\infty}\le (4\pi|t|)^{-1/2}\norm f_{L^1}.
\]
\end{lemma}

\begin{proof}
For Schwartz $f$, the formula is \eqref{eq:free-kernel-explicit}.  Its
kernel has constant modulus $(4\pi|t|)^{-1/2}$, which gives the estimate.
For general $f\in L^1$, choose Schwartz $f_n\to f$ in $L^1$.  The kernel
integrals converge uniformly in $x$, so they define the unique
$L^1\to L^\infty$ extension and preserve the same bound.
\end{proof}

\begin{lemma}[Compatibility of the $L^2$ and kernel propagators]
\label{lem:propagator-compatibility}
For $t\ne0$, let $S_2(t)$ be the $L^2$ Fourier multiplier
\eqref{eq:free-group}, and let $S_{1,\infty}(t)$ be the integral operator with
kernel
\[
 K_t(x-y)=(4\pi it)^{-1/2}\exp\!\left(\frac{i|x-y|^2}{4t}\right).
\]
If $f\in L^1\cap L^2$, then the two outputs agree almost everywhere.  Thus
they are compatible endpoint realizations of one linear operator.
\end{lemma}

\begin{proof}
The equality holds for Schwartz functions by \cref{lem:complex-gaussian-kernel}.  Given $f\in L^1\cap L^2$, choose Schwartz $f_n$ converging
to $f$ simultaneously in $L^1$ and $L^2$; convolution followed by spatial
cutoff gives such a sequence.  The kernel formula gives uniform convergence
$S_{1,\infty}(t)f_n\to S_{1,\infty}(t)f$, while unitarity gives convergence
$S_2(t)f_n\to S_2(t)f$ in $L^2$.  A subsequence of the latter converges almost
everywhere.  The two limits therefore agree almost everywhere.
\end{proof}


\begin{lemma}[Hadamard three-lines estimate]\label{lem:three-lines-explicit}
Let $F$ be bounded and continuous on the closed strip
$\{z:0\le\Rea z\le1\}$ and holomorphic in its interior.  Put
\[
 M_j=\sup_{y\in\R}|F(j+iy)|\qquad(j=0,1).
\]
Then, for $0\le\theta\le1$,
\begin{equation}
 |F(\theta)|\le M_0^{1-\theta}M_1^\theta.          \label{eq:three-lines-explicit}
\end{equation}
\end{lemma}

\begin{proof}
The boundary cases are tautological.  Fix $0<\theta<1$.  First suppose
$M_0,M_1>0$, and use the real logarithms of these positive numbers to define
\[
 G(z)=F(z)\exp\big((z-1)\log M_0-z\log M_1\big).
\]
Thus $|G(iy)|\le1$ and $|G(1+iy)|\le1$.  Let
$B=\sup_{0\le\Rea z\le1}|G(z)|<\infty$.  For $\delta>0$ set
\[
 G_\delta(z)=G(z)\e^{\delta z^2}.
\]
On the two vertical sides of the rectangle
$[0,1]\times[-R,R]$, its modulus is at most $\e^\delta$.  On a horizontal
side $z=x\pm iR$,
\[
 |G_\delta(z)|\le B\exp\big(\delta(x^2-R^2)\big)
 \le B\exp\big(\delta(1-R^2)\big),
\]
which is at most $\e^\delta$ once $R$ is sufficiently large.  The maximum
modulus principle on this rectangle therefore gives
$|G_\delta(\theta)|\le\e^\delta$.  Since
$|G_\delta(\theta)|=|G(\theta)|\e^{\delta\theta^2}$, letting
$\delta\downarrow0$ yields $|G(\theta)|\le1$, which is exactly
\eqref{eq:three-lines-explicit}.

If one of the $M_j$ is zero, fix $\eps>0$ and repeat the same rectangle
argument with the positive boundary majorants
$A_j=M_j+\eps$ in place of $M_j$.  Since $|F(j+iy)|\le A_j$, it gives
$|F(\theta)|\le A_0^{1-\theta}A_1^\theta$.  Letting
$\eps\downarrow0$ proves the asserted inequality, including the case in
which the right side is zero.
\end{proof}

\begin{lemma}[Riesz--Thorin interpolation in the form used here]
\label{lem:riesz-thorin-used}
Let a linear operator $T$ be defined on simple functions and have compatible
extensions
\[
 T:L^{p_0}\to L^{q_0},\qquad \norm T\le M_0,
 \quad\text{and}\quad
 T:L^{p_1}\to L^{q_1},\qquad \norm T\le M_1,
\]
where $1\le p_j<\infty$ and $1\le q_j\le\infty$.  If
$0<\theta<1$ and
\[
 \frac1{p_\theta}=\frac{1-\theta}{{p_0}}+\frac\theta{p_1},
 \qquad
 \frac1{q_\theta}=\frac{1-\theta}{q_0}+\frac\theta{q_1},
\]
then
\[
 \norm{Tf}_{q_\theta}
 \le M_0^{1-\theta}M_1^\theta\norm f_{p_\theta}.
\]
\end{lemma}

\begin{proof}
We first reduce the interpolation estimate to the proved three-lines lemma.  It is enough to
consider simple $f$ and simple dual test functions $g$, normalized by
$\norm f_{p_\theta}=\norm g_{q_\theta'}=1$.  On the strip
$0\le\Re z\le1$, define
\[
 \alpha(z)=\frac{1-z}{p_0}+\frac z{p_1},
 \qquad
 \gamma(z)=\frac{1-z}{q_0'}+\frac z{q_1'},
\]
and, on the nonzero level sets of the simple functions,
\[
 f_z=|f|^{p_\theta\alpha(z)}\frac f{|f|},
 \qquad
 g_z=|g|^{q_\theta'\gamma(z)}\frac g{|g|}.
\]
Zero values are left zero.  The finite-sum expression
\[
 \Phi(z)=\int (Tf_z)g_z
\]
is analytic in the open strip and continuous on its closure.  On
$\Re z=j$, direct calculation gives
$\norm{f_z}_{p_j}=\norm{g_z}_{q_j'}=1$, hence
$|\Phi(z)|\le M_j$.  \Cref{lem:three-lines-explicit} yields
$|\Phi(\theta)|\le M_0^{1-\theta}M_1^\theta$.  At $z=\theta$ we have
$f_z=f$ and $g_z=g$.  Taking the supremum over simple dual test functions and then completing
proves the estimate.  If $q_j=\infty$, then $q_j'=1$ and the corresponding
boundary estimate is simply the $L^\infty$--$L^1$ H\"older inequality; the
same formulas remain literal because $1/q_j=0$.  Thus no limiting argument
at the output endpoint is being suppressed.  Input exponents are finite in
the stated form, which is exactly the case used below.
\end{proof}

Applying \cref{lem:riesz-thorin-used} to the compatible operators in
\cref{lem:propagator-compatibility}, with
$(p_0,q_0)=(1,\infty)$, $(p_1,q_1)=(2,2)$, and
$\theta=2/r$, yields, for $2\le r\le\infty$,
\begin{equation}
 \norm{S(t)f}_{L^r}
 \le C_r|t|^{-(1/2-1/r)}\norm f_{L^{r'}}.
 \label{eq:disp-interp}
\end{equation}
The extension on $L^{r'}$ is unique because $L^1\cap L^2$ is dense in
$L^{r'}$ when $1<r'<2$, and it is compatible with both endpoint
realizations.

\section{A one-dimensional fractional-integration estimate}
We need one scalar convolution theorem.  A proof is included to keep the
Strichartz argument closed.

\begin{lemma}[Hardy--Littlewood--Sobolev on the time line]\label{lem:HLS-time}
Let $0<\alpha<1$, $1<p<q<\infty$, and
$1/q=1/p-(1-\alpha)$.  Then
\[
 \left\|\int_\R |t-s|^{-\alpha}F(s)\,\dd s\right\|_{L^q_t}
 \le C_{p,\alpha}\norm F_{L^p_t}.
\]
\end{lemma}

\begin{proof}
We include the maximal-function input.  For $F\in L^1_{\loc}(\R)$ set
\[
 MF(t)=\sup_{R>0}\frac1{2R}\int_{t-R}^{t+R}|F(s)|\,\dd s.
\]
The centered maximal operator is of weak type $(1,1)$:
\begin{equation}
 \big|\{MF>\lambda\}\big|\le \frac{C}{\lambda}\norm F_1.
 \label{eq:maximal-weak11}
\end{equation}
Indeed, let $K$ be a compact subset of $\{MF>\lambda\}$.  For every
$t\in K$ choose an interval centered at $t$ on which the average of $|F|$
is larger than $\lambda$.  Compactness gives a finite cover.  Select from it,
greedily in decreasing order of length, a pairwise disjoint subfamily.  Every
unselected interval is contained in the triple of a selected one, so the
triples cover $K$.  Therefore
\[
 |K|\le3\sum_j|I_j|
 <\frac3\lambda\sum_j\int_{I_j}|F|
 \le\frac3\lambda\norm F_1.
\]
Inner regularity gives \eqref{eq:maximal-weak11}.

For $1<p<\infty$, this weak estimate implies the strong bound
\begin{equation}
 \norm{MF}_p\le C_p\norm F_p.                      \label{eq:maximal-strongp}
\end{equation}
To see this without invoking an interpolation theorem, fix $\lambda>0$ and
write
$F=F\1_{\{|F|>\lambda/2\}}+F\1_{\{|F|\le\lambda/2\}}$.
The maximal function of the second summand is at most $\lambda/2$, hence
\[
 |\{MF>\lambda\}|
 \le\frac C\lambda\int_{\{|F|>\lambda/2\}}|F(s)|\,\dd s.
\]
Multiply by $p\lambda^{p-1}$, integrate in $\lambda$, and use Tonelli:
\begin{align*}
 \norm{MF}_p^p
 &\le Cp\int_0^\infty\lambda^{p-2}
       \int_{\{|F|>\lambda/2\}}|F(s)|\,\dd s\,\dd\lambda\\
 &=Cp\int_\R |F(s)|
      \int_0^{2|F(s)|}\lambda^{p-2}\,\dd\lambda\,\dd s
 \le C_p\norm F_p^p.
\end{align*}

Now fix $R>0$ and split the fractional integral into
$|t-s|<R$ and $|t-s|\ge R$.  Dyadic decomposition of the near part gives
\[
 \int_{|t-s|<R}|t-s|^{-\alpha}|F(s)|\,\dd s
 \le C R^{1-\alpha}MF(t).
\]
For the far part, H\"older gives
\[
 \int_{|t-s|\ge R}|t-s|^{-\alpha}|F(s)|\,\dd s
 \le C R^{1-\alpha-1/p}\norm F_p.
\]
Choose $R=(\norm F_p/MF(t))^p$ when $MF(t)>0$; the cases
$MF(t)=0$ and $\norm F_p=0$ are immediate.  Since
$p(1-\alpha)=1-p/q$, the two estimates yield the Hedberg inequality
\[
 I_{1-\alpha}F(t)
 \le C(MF(t))^{p/q}\norm F_p^{1-p/q}.
\]
Raise to the $q$th power, integrate, and apply
\eqref{eq:maximal-strongp}.
\end{proof}

\section{Homogeneous estimates}
A pair $(q,r)$ is called one-dimensional Schr\"odinger admissible if
\begin{equation}
  2\le q,r\le\infty,
  \qquad \frac2q+\frac1r=\frac12.                 \label{eq:admissible}
\end{equation}

\begin{theorem}[Homogeneous Strichartz estimates]\label{thm:hom-strichartz}
For every admissible pair $(q,r)$,
\begin{equation}
  \norm{S(t)f}_{L^q_t(\R;L^r_x)}\le C_{q,r}\norm f_2. \label{eq:hom-str}
\end{equation}
In particular,
\[
 \norm{S(t)f}_{L^6_{t,x}}+\norm{S(t)f}_{L^4_tL^\infty_x}
 \le C\norm f_2.
\]
\end{theorem}

\begin{proof}
We first treat $2<q<\infty$ and $r<\infty$.  On compactly supported smooth
spacetime functions define
\[
 Tf(t)=S(t)f,
 \qquad
 T^*F=\int_\R S(-s)F(s)\,\dd s.
\]
The latter is an $L^2$ Bochner integral for such $F$.  The full convolution is
\[
 TT^*F(t)=\int_\R S(t-s)F(s)\,\dd s.
\]
By \eqref{eq:disp-interp} and Minkowski,
\[
 \norm{TT^*F(t)}_{L^r_x}
 \le C\int_\R|t-s|^{-(1/2-1/r)}
              \norm{F(s)}_{L^{r'}_x}\,\dd s.
\]
Admissibility says $1/2-1/r=2/q$.  The time-line HLS estimate with
$p=q'$ and $\alpha=2/q$ therefore gives
\begin{equation}
 \norm{TT^*F}_{L^q_tL^r_x}
 \le C\norm F_{L^{q'}_tL^{r'}_x}.                 \label{eq:TTstar-finite}
\end{equation}
Consequently,
\[
 \norm{T^*F}_2^2
 =\left|\iint (TT^*F)(t,x)\overline{F(t,x)}\,\dd x\,\dd t\right|
 \le C\norm F_{L^{q'}L^{r'}}^2.
\]
Thus $T^*$ extends boundedly to $L^{q'}L^{r'}$, and ordinary finite-exponent
duality yields \eqref{eq:hom-str}.

It remains to justify $(q,r)=(4,\infty)$ without identifying a nonseparable
Bochner dual.  For smooth compactly supported $F$, the dispersive estimate
gives directly
\[
 m_\infty(TT^*F)(t)
 \le C\int_\R|t-s|^{-1/2}\norm{F(s)}_{L^1_x}\,\dd s.
\]
HLS therefore yields the scalar endpoint estimate
\begin{equation}
 \norm{TT^*F}_{L^4_tL^\infty_x}
 \le C\norm F_{L^{4/3}_tL^1_x}.                  \label{eq:TTstar-endpoint}
\end{equation}
The same quadratic identity as above gives
\[
 \norm{T^*F}_2\le C\norm F_{L^{4/3}_tL^1_x}.
\]
Now let $f$ be Schwartz.  Before applying the norming identity, note that
$S(t)f$ already has finite scalar endpoint norm: Fourier inversion gives
\[
 \norm{S(t)f}_\infty\le(2\pi)^{-1/2}\norm{\widehat f}_1
 \quad(|t|\le1),
\]
while \cref{lem:dispersive-1d} gives
$\norm{S(t)f}_\infty\le C|t|^{-1/2}\norm f_1$ for $|t|>1$.
The fourth power of the latter bound is integrable at infinity.  Thus
$S(t)f\in L^4_tL^\infty_x$ in the scalar sense.  For every smooth compactly
supported $G$,
\[
 \left|\iint S(t)f(x)\overline{G(t,x)}\,\dd x\,\dd t\right|
 =|\pair f{T^*G}_{L^2}|
 \le C\norm f_2\norm G_{L^{4/3}_tL^1_x}.
\]
Smooth compactly supported functions are dense in the norming class of
\cref{prop:endpoint-norming}; approximation in $L^{4/3}_tL^1_x$ and the last
bound extend the inequality to that class.  The norming identity then gives
\[
 \norm{S(t)f}_{L^4_tL^\infty_x}\le C\norm f_2.
\]

For a general $f\in L^2$, choose Schwartz $f_n\to f$ in $L^2$.  The estimate
makes $S(t)f_n$ Cauchy in the complete scalar space
$L^4_tL^\infty_x$, so it converges there to some jointly measurable $H$.
On every bounded time interval, unitarity also gives
\[
 \int_I\norm{S(t)(f_n-f)}_2^2\,\dd t
 =|I|\norm{f_n-f}_2^2\longrightarrow0.
\]
Choose a subsequence whose successive differences are summable both in
$L^4_tL^\infty_x$ and in $L^2(I\times\R)$.  By
\cref{lem:mixed-subsequence-uniform}, for almost every $t$ this subsequence
converges in the spatial essential-supremum norm to the section $H(t)$.  By
the ordinary summable-subsequence argument in $L^2(I\times\R)$, after a
further subsequence it converges pointwise almost everywhere to the canonical
spacetime representative of $S(t)f$.  The further subsequence retains the
sectionwise uniform convergence.  Hence the two limits agree for almost every
$(t,x)$.  This identifies $H$ with $S(t)f$ and proves the endpoint estimate
for all $f\in L^2$.  A countable exhaustion of the time axis makes the
identification global.

Finally, $(q,r)=(\infty,2)$ is unitarity.  These cases exhaust the admissible
pairs, with the finite pairs handled by \eqref{eq:TTstar-finite} and the sole
spatial endpoint handled separately above.
\end{proof}

\section{Retarded estimates}
We first isolate the elementary time-ordering lemma.

\begin{lemma}[Christ--Kiselev time ordering]\label{lem:CK}
Let $1\le p<q\le\infty$ and let
$K:L^p(I;A)\to L^q(I;B)$ be a bounded integral operator
\[
 KF(t)=\int_I K(t,s)F(s)\,\dd s
\]
between Banach-valued function spaces.  Then the retarded operator
\[
 K_{<}F(t)=\int_{s<t}K(t,s)F(s)\,\dd s
\]
is bounded from $L^p(I;A)$ to $L^q(I;B)$, with norm depending only on
$p,q$ and the norm of $K$.
\end{lemma}

\begin{proof}
Normalize $\norm F_{L^p}=1$ and define
$\Phi(t)=\int_{I\cap(-\infty,t)}\norm{F(s)}_A^p\,\dd s$.  Up to null sets,
$\Phi$ is nondecreasing and takes values in $[0,1]$.  At generation $n$,
pair the two sibling dyadic intervals
\[
 E^-_{n,j}=[2j2^{-n},(2j+1)2^{-n}),
 \qquad
 E^+_{n,j}=[(2j+1)2^{-n},(2j+2)2^{-n}),
\]
and set $I^\pm_{n,j}=\Phi^{-1}(E^\pm_{n,j})$.  Every pair
$(s,t)$ with $s<t$ and $\Phi(s)<\Phi(t)$ belongs to exactly one rectangle
$I^-_{n,j}\times I^+_{n,j}$: take the first binary generation at which the
two values lie in different children.  The remaining pairs lie on level sets
of $\Phi$, where $F=0$ almost everywhere, so they do not affect the integral.
Thus the retarded operator is the sum over $n,j$ of
\[
 \1_{I^+_{n,j}}(t)
 \int_{I^-_{n,j}}K(t,s)F(s)\,\dd s.
\]
For fixed $n$ the output sets are disjoint, and every input piece has
$L^p$ norm at most $2^{-n/p}$.  If the norm of the full operator is $M$, the
$L^q$ norm of the $n$th generation is therefore at most
\[
 M\left(\sum_j 2^{-nq/p}\right)^{1/q}
 \le CM2^{n/q}2^{-n/p},
\]
When $q=\infty$, disjointness is replaced by taking the maximum over $j$,
and the same bound reads $M2^{-n/p}$ because $2^{n/q}=1$.  Summing in $n$
gives the convergent geometric series
$\sum_n2^{-n(1/p-1/q)}$.
\end{proof}

\begin{theorem}[Inhomogeneous estimates used below]\label{thm:inhom-str}
For every interval $I$ and every $F\in L^{6/5}(I\times\R)$,
\begin{align}
 \left\|\int_{t_0}^tS(t-s)F(s)\,\dd s\right\|_{L^6(I\times\R)}
 &\le C\norm F_{L^{6/5}(I\times\R)},                 \label{eq:inhom66}\\
 \left\|\int_{t_0}^tS(t-s)F(s)\,\dd s\right\|_{L^\infty_tL^2_x(I)}
 &\le C\norm F_{L^{6/5}(I\times\R)}.               \label{eq:inhomL2}
\end{align}
The constants do not depend on $I$ or $t_0\in I$.
\end{theorem}

\begin{proof}
The full convolution estimate from $L^{6/5}_{t,x}$ to $L^6_{t,x}$ is the
$TT^*$ estimate in the proof of \cref{thm:hom-strichartz}.  Apply
\cref{lem:CK} with time exponents $6/5<6$ to restrict to $s<t$; reversing
time handles $t<t_0$.  For the $L^\infty_tL^2_x$ estimate, duality and the
homogeneous $(6,6)$ estimate give
\[
 \left\|\int_I S(-s)F(s)\,\dd s\right\|_2
 \le C\norm F_{6/5}.
\]
The corresponding full operator
\[
 F\longmapsto S(t)\int_I S(-s)F(s)\,\dd s
\]
maps $L^{6/5}_{t,x}$ to $L^\infty_tL^2_x$.  Apply
\cref{lem:CK} with input time exponent $6/5$ and output exponent $\infty$
to restrict the integral to $s<t$; reversing time handles the other
orientation around $t_0$.  This proves \eqref{eq:inhomL2}.
\end{proof}

\begin{lemma}[Continuous retarded operator, base time, and Bochner compatibility]
\label{lem:retarded-continuous}
Let $I$ be a bounded interval and $t_0\in I$.  For smooth compactly supported
$F$ put
\[
 (\mathcal R_{t_0}F)(t)=\int_{t_0}^tS(t-s)F(s)\,\dd s.
\]
This operator extends uniquely to a bounded map
\begin{equation}
 \mathcal R_{t_0}:L^{6/5}(I\times\R)
 \longrightarrow C(I;L^2(\R))\cap L^6(I\times\R),               \label{eq:retarded-continuous-map}
\end{equation}
with the estimates of \cref{thm:inhom-str}.  If $t_1\in I$, then
\begin{equation}
 \mathcal R_{t_0}F(t)
 =\mathcal R_{t_1}F(t)+S(t-t_1)\mathcal R_{t_0}F(t_1)             \label{eq:retarded-base-time}
\end{equation}
for every $t\in I$.  If in addition
$F\in L^1(I;L^2)\cap L^{6/5}(I\times\R)$, the completed operator is the
ordinary $L^2$-valued Bochner integral displayed above.
\end{lemma}

\begin{proof}
For smooth compactly supported $F$, the displayed integral is an ordinary
Bochner integral in $L^2$.  Strong continuity of $S(t)$ and dominated
convergence show that it is continuous in $t$.  Let $F_n$ be smooth compactly
supported and Cauchy in $L^{6/5}$.  By \cref{thm:inhom-str},
$\mathcal R_{t_0}F_n$ is Cauchy in both $L^6_{t,x}$ and
$L^\infty_tL^2_x$.  The difference of two such outputs is a continuous
$L^2$-valued function, so the essential supremum of its $L^2$ norm equals the
ordinary supremum.  Hence the sequence is uniformly Cauchy in
$C(I;L^2)$ and converges there to a continuous limit.  Density and the two
inhomogeneous estimates prove existence, uniqueness, and
\eqref{eq:retarded-continuous-map}.

For smooth forcing, split the oriented integral at $t_1$ and use the group
law:
\begin{align*}
 \int_{t_0}^tS(t-s)F(s)\,\dd s
 &=\int_{t_1}^tS(t-s)F(s)\,\dd s
   +S(t-t_1)\int_{t_0}^{t_1}S(t_1-s)F(s)\,\dd s.
\end{align*}
Approximation in $L^{6/5}$, continuity of evaluation on $C(I;L^2)$, and
strong continuity of $S(t-t_1)$ give \eqref{eq:retarded-base-time} for all
forcing in the completed domain.

Finally suppose $F$ belongs also to $L^1_tL^2_x$.  Approximate it by smooth
compactly supported functions in the intersection norm
$\norm\cdot_{L^{6/5}}+\norm\cdot_{L^1L^2}$.  Unitarity and scalar Minkowski
give
\[
 \sup_{t\in I}\left\|\int_{t_0}^tS(t-s)(F_n-F)(s)\,\dd s\right\|_2
 \le\norm{F_n-F}_{L^1_tL^2_x}.
\]
Thus the Bochner integrals converge uniformly in $L^2$ to the same continuous
limit as the Strichartz completion.
\end{proof}

\begin{checkpoint}
The free estimates required for the nonlinear flow have now been proved from
Plancherel, the dispersive kernel, scalar fractional integration, and time
ordering.  No endpoint or inhomogeneous Strichartz theorem remains external.
\end{checkpoint}

\chapter{Global \texorpdfstring{$L^2$}{L2} well-posedness for the cubic flow}\label{chap:wellposed}

\section{The fixed-point space}
Let $I=[t_0-T,t_0+T]$.  Set
\[
  \mathcal X(I)=C(I;L^2(\R))\cap L^6(I\times\R),
  \qquad
  \norm u_{\mathcal X(I)}
  =\norm u_{L^\infty_tL^2_x(I)}+\norm u_{L^6_{t,x}(I)}.
\]
The space is Banach.  Indeed, it is the closed diagonal subspace of
$C(I;L^2)\oplus_1L^6(I\times\R)$ after identifying the two components as
distributions; equality of the two distributional limits follows from
convergence against compactly supported smooth tests.
The nonlinear map associated with initial data $u_0$ is
\begin{equation}
 \Phi_{u_0}(u)(t)
 =S(t-t_0)u_0-i\sigma\int_{t_0}^tS(t-s)(|u|^2u)(s)\,\dd s.  \label{eq:Phi}
\end{equation}
Here the integral denotes the completed retarded operator of
\cref{lem:retarded-continuous}; after \cref{prop:endpoint} it is also the
ordinary Bochner integral.  The key smallness is supplied by the length of
$I$, not by the mass.

\begin{lemma}[Cubic estimate]\label{lem:cubic-65}
For $u\in\mathcal X(I)$,
\begin{equation}
 \norm{|u|^2u}_{L^{6/5}_{t,x}(I)}
 \le |I|^{1/2}\norm u_{L^\infty_tL^2_x(I)}
                    \norm u_{L^6_{t,x}(I)}^2.       \label{eq:cubic65}
\end{equation}
If $u,v\in\mathcal X(I)$, then
\begin{align}
 \norm{|u|^2u-|v|^2v}_{L^{6/5}_{t,x}(I)}
 \le C|I|^{1/2}
  (&\norm u_{\mathcal X(I)}^2+\norm v_{\mathcal X(I)}^2)
  \norm{u-v}_{\mathcal X(I)}.                      \label{eq:cubic-lip}
\end{align}
\end{lemma}

\begin{proof}
At each time, H\"older in space gives
\[
 \norm{|u|^2u}_{L^{6/5}_x}
 \le \norm u_2\norm u_6^2,
 \qquad \frac56=\frac12+\frac16+\frac16.
\]
The function $t\mapsto\norm{u(t)}_6^2$ belongs to $L^3(I)$.  Embedding
$L^3(I)$ into $L^{6/5}(I)$ costs the factor
$|I|^{5/6-1/3}=|I|^{1/2}$, which proves \eqref{eq:cubic65}.
For the difference, use
\[
 ||u|^2u-|v|^2v|
 \le C(|u|^2+|v|^2)|u-v|
\]
and assign one of the three factors to $L^2_x$ and the other two to
$L^6_x$.  The same time estimate completes the proof.
\end{proof}

\begin{theorem}[Local existence, uniqueness, and Lipschitz dependence]
\label{thm:LWP}
For every $M>0$ there is $T_M>0$ such that whenever
$\norm{u_0}_2\le M$, equation \eqref{eq:NLS-main} has a unique mild solution
on $[t_0-T_M,t_0+T_M]$.  It belongs to $\mathcal X(I)$, and the data-to-solution
map is locally Lipschitz from the $L^2$ ball of radius $M$ to
$\mathcal X(I)$.  One may take $T_M=c(1+|\sigma|M^2)^{-2}$ with a sufficiently
small absolute $c$.
\end{theorem}

\begin{proof}
By \cref{thm:hom-strichartz,thm:inhom-str},
\[
 \norm{\Phi_{u_0}(u)}_{\mathcal X(I)}
 \le C\norm{u_0}_2+C|\sigma|\norm{|u|^2u}_{6/5}.
\]
Fix $R=2CM$ and consider the closed ball
$B_R\subset\mathcal X(I)$.  By \cref{lem:cubic-65},
\[
 \norm{\Phi_{u_0}(u)}_{\mathcal X(I)}
 \le CM+C|\sigma||I|^{1/2}R^3.
\]
Choose $|I|$ so that $C|\sigma||I|^{1/2}R^2\le1/4$.  The same choice and
\eqref{eq:cubic-lip} make $\Phi_{u_0}$ a strict contraction on $B_R$.
Banach's theorem gives a unique fixed point in $B_R$.  The difference
estimate with two initial data gives
\[
 \norm{u-v}_{\mathcal X(I)}
 \le C\norm{u_0-v_0}_2+\tfrac12\norm{u-v}_{\mathcal X(I)},
\]
which proves local Lipschitz dependence for the constructed solutions.
\Cref{lem:retarded-continuous} shows directly that the fixed point belongs to
$C(I;L^2)$.

It remains to record uniqueness among all mild solutions in
$\mathcal X(I)$, not only fixed points lying in the chosen ball.  Let
$u$ and $v$ be two such solutions with the same value at some time $s\in I$.
Their $\mathcal X(I)$ norms are finite.  Choose a subinterval $J\subset I$
containing $s$ so short that
\[
 C|\sigma||J|^{1/2}
   (\norm u_{\mathcal X(I)}^2+\norm v_{\mathcal X(I)}^2)<\frac12.
\]
The base-time identity \eqref{eq:retarded-base-time} writes both mild
formulas with base time $s$.  The cubic difference estimate on $J$ then gives
$\norm{u-v}_{\mathcal X(J)}\le\tfrac12\norm{u-v}_{\mathcal X(J)}$, so the
solutions agree on $J$.  Repeating this argument along a finite chain of
short overlapping intervals covers $I$.  This proves unconditional
uniqueness in the stated solution class and also justifies the phrase ``for
one and hence every base time'' in the definition of a mild solution.
\end{proof}

\section{Local persistence of one derivative}
We first prove persistence on one mass-controlled fixed-point interval.  Mass
conservation will then globalize the statement without any circular use of
higher regularity.

\begin{lemma}[Derivative estimate for the cubic map]
\label{lem:cubic-derivative-map}
Let $J$ be bounded and let $u,w\in\mathcal X(J)$.  Define the real-linear
expression
\[
 DN(u)[w]=2\Rea(\overline u\,w)u+|u|^2w.
\]
Then
\begin{equation}
 \norm{DN(u)[w]}_{L^{6/5}_{t,x}(J)}
 \le C|J|^{1/2}\norm u_{\mathcal X(J)}^2
                      \norm w_{\mathcal X(J)}.     \label{eq:DN-bound}
\end{equation}
Moreover,
\begin{align}
 \norm{DN(u)[w]-DN(\widetilde u)[\widetilde w]}_{6/5}
 \le C|J|^{1/2}\Big[&
 (\norm u_{\mathcal X}^2+\norm{\widetilde u}_{\mathcal X}^2)
 \norm{w-\widetilde w}_{\mathcal X}\notag\\
 &+(\norm u_{\mathcal X}+\norm{\widetilde u}_{\mathcal X})
 (\norm w_{\mathcal X}+\norm{\widetilde w}_{\mathcal X})
 \norm{u-\widetilde u}_{\mathcal X}\Big].          \label{eq:DN-Lipschitz}
\end{align}
\end{lemma}

\begin{proof}
Pointwise, $|DN(u)[w]|\le3|u|^2|w|$.  At each time assign one $u$
factor to $L^2_x$, the other $u$ factor to $L^6_x$, and $w$ to
$L^6_x$; the alternative allocation with $w$ in $L^2_x$ is also
available.  The spatial exponents add to $5/6$.  The two $L^6_t$
factors give an $L^3_t$ product, and
$L^3(J)\hookrightarrow L^{6/5}(J)$ costs $|J|^{1/2}$.  This proves
\eqref{eq:DN-bound}.  Expanding the difference and using
\[
 |DN(u)[w]-DN(\widetilde u)[\widetilde w]|
 \le C(|u|^2+|\widetilde u|^2)|w-\widetilde w|
  +C(|u|+|\widetilde u|)(|w|+|\widetilde w|)|u-\widetilde u|
\]
gives \eqref{eq:DN-Lipschitz} by the same allocation.
\end{proof}

\begin{lemma}[Local $H^1$ persistence]
\label{lem:H1-persistence}
Let $J$ be a local fixed-point interval centered at $t_0$ on which
$u\in\mathcal X(J)$ and
\begin{equation}
 C|\sigma||J|^{1/2}\norm u_{\mathcal X(J)}^2\le\frac14. \label{eq:H1-smallness}
\end{equation}
If $u(t_0)\in H^1(\R)$, then
\[
 u\in C(J;H^1)\cap L^6(J;W^{1,6}).
\]
The equation holds in $C(J;H^{-1})$.  The size of $J$ needed here depends
only on an upper bound for the $L^2$ mass, not on the $H^1$ norm.
\end{lemma}

\begin{proof}
For $z\in\mathcal X(J)$ define
\begin{equation}
 \Psi_u(z)(t)=S(t-t_0)u(t_0)'
 -i\sigma\int_{t_0}^tS(t-s)DN(u(s))[z(s)]\,\dd s.  \label{eq:derivative-fixed-point}
\end{equation}
By \cref{thm:inhom-str,lem:cubic-derivative-map} and
\eqref{eq:H1-smallness}, this is a strict contraction on
$\mathcal X(J)$.  It has a unique fixed point $z$ and
\[
 \norm z_{\mathcal X(J)}\le 2C\norm{u(t_0)'}_2.
\]

We identify $z$ with the distributional derivative without assuming in
advance that smooth data generate smooth solutions.  Write
$\tau_hf(x)=f(x+h)$ and
\[
 q_h=\frac{\tau_hu-u}{h},\qquad
 q_{0,h}=\frac{\tau_hu(t_0)-u(t_0)}h.
\]
Translation invariance and uniqueness imply that $\tau_hu$ is the solution
with datum $\tau_hu(t_0)$.  Since translations are strongly continuous on
$L^2$, local flow stability gives
\begin{equation}
 \tau_hu\longrightarrow u\quad\text{in }\mathcal X(J). \label{eq:translation-flow-convergence}
\end{equation}
The pointwise fundamental theorem of calculus for the polynomial map
$N(a)=|a|^2a$, viewed as a $C^1$ map on the real plane, gives
\[
 \frac{N(\tau_hu)-N(u)}h=A_h[q_h],
 \qquad
 A_h[w]=\int_0^1DN\bigl(u+\theta(\tau_hu-u)\bigr)[w]\,\dd\theta.
\]
Consequently $q_h$ satisfies
\begin{equation}
 q_h(t)=S(t-t_0)q_{0,h}
 -i\sigma\int_{t_0}^tS(t-s)A_h[q_h](s)\,\dd s.      \label{eq:qh-equation}
\end{equation}
For every $\theta\in[0,1]$, spatial translation invariance and the
triangle inequality give
\[
 \norm{(1-\theta)u+\theta\tau_hu}_{\mathcal X(J)}
 \le\norm u_{\mathcal X(J)}.
\]
Thus \eqref{eq:DN-bound} and \eqref{eq:H1-smallness} make the linear part of
\eqref{eq:qh-equation} a contraction with constant at most $1/4$, uniformly
in $h$.  Also $q_{0,h}\to u(t_0)'$ in $L^2$.  Subtract
\eqref{eq:derivative-fixed-point} from \eqref{eq:qh-equation}; then
\eqref{eq:DN-Lipschitz} and \eqref{eq:translation-flow-convergence} give
\[
 \norm{q_h-z}_{\mathcal X(J)}
 \le C\norm{q_{0,h}-u(t_0)'}_2
   +\frac12\norm{q_h-z}_{\mathcal X(J)}+o(1).
\]
Thus $q_h\to z$ in $\mathcal X(J)$.  On the other hand, translation
difference quotients converge to $u_x$ in distributions.  Hence $z=u_x$,
which proves $u\in C H^1\cap L^6W^{1,6}$.

Finally,
$u_t=i u_{xx}-i\sigma|u|^2u$.  We have $u_{xx}\in C H^{-1}$, while the
one-dimensional estimate in \cref{lem:H1-C0} gives
$H^1\hookrightarrow L^\infty$, so $|u|^2u\in C L^2\subset C H^{-1}$.
This proves the last assertion.
\end{proof}

\section{Mass conservation and globalization}

\begin{lemma}[Gelfand-triple chain rule]\label{lem:gelfand-chain-rule}
Let $V\hookrightarrow H\hookrightarrow V^*$ be a dense continuous Gelfand
triple, with $H$ identified with its anti-dual in the usual way.  Suppose
$w\in C([a,b];V)$ and its distributional time derivative belongs to
$C([a,b];V^*)$.  Then $t\mapsto\norm{w(t)}_H^2$ is continuously
differentiable and
\begin{equation}
 \frac{\dd}{\dd t}\norm{w(t)}_H^2
 =2\Rea\pair{w_t(t)}{w(t)}_{V^*,V}.                \label{eq:gelfand-chain-rule}
\end{equation}
The same conclusion, with absolute continuity instead of continuous
differentiability, holds when $w\in L^p(a,b;V)$ and
$w_t\in L^{p'}(a,b;V^*)$ in the usual dual-exponent formulation and $w$
has its standard continuous $H$ representative.
\end{lemma}

\begin{proof}
Extend $w$ constantly a short distance beyond each endpoint and extend
$w_t$ by zero there.  Because the extended $w$ is continuous, its
distributional derivative is exactly the zero extension of $w_t$; no endpoint
Dirac mass occurs.  Convolve both in time with a smooth approximate identity.
The mollified map $w_\eps$ is smooth with values in $V$, and its classical
derivative agrees in $V^*$ with $(w_t)_\eps$.  The Hilbert-space product rule
gives
\[
 \norm{w_\eps(t)}_H^2-\norm{w_\eps(s)}_H^2
 =2\Rea\int_s^t
   \pair{(w_t)_\eps(\tau)}{w_\eps(\tau)}_{V^*,V}\,\dd\tau.
\]
Uniform continuity on compact subintervals gives
$w_\eps\to w$ uniformly in $V$ and $(w_t)_\eps\to w_t$ uniformly in
$V^*$.  The dual pairings converge uniformly, so passage to the limit gives
the integrated form of \eqref{eq:gelfand-chain-rule}; continuity of its
right side gives the pointwise derivative.  In the integrable case, first
work on compact subintervals and use convergence in the corresponding
Bochner spaces together with H\"older.  The integrated identity supplies the
continuous $H$ representative and absolute continuity.
\end{proof}

\begin{lemma}[Mass conservation for local $H^1$ solutions]\label{lem:mass-H1}
If $u\in C(J;H^1)$ solves \eqref{eq:NLS-main} on $J$, then
\[
  \norm{u(t)}_2=\norm{u(t_0)}_2\qquad(t,t_0\in J).
\]
\end{lemma}

\begin{proof}
Apply \cref{lem:gelfand-chain-rule} to
$H^1\hookrightarrow L^2\hookrightarrow H^{-1}$.  Since
$u_t=i u_{xx}-i\sigma|u|^2u$ in $H^{-1}$,
\begin{align*}
 \frac{\dd}{\dd t}\norm u_2^2
 &=2\Rea\pair{i u_{xx}-i\sigma|u|^2u}{u}_{H^{-1},H^1}\\
 &=2\Rea\left(-i\norm{u_x}_2^2-i\sigma\norm u_4^4\right)=0.
\end{align*}
Integration by parts in the first pairing is the defining weak-derivative
identity.  Both quantities in parentheses are real.
\end{proof}

\begin{proposition}[Mass conservation in $L^2$]\label{prop:mass-L2}
Every $L^2$ mild solution satisfies
\begin{equation}
  \norm{u(t)}_2=\norm{u(t_0)}_2\qquad(t,t_0\text{ in its lifespan}).
  \label{eq:mass-conservation}
\end{equation}
\end{proposition}

\begin{proof}
Fix a mass-controlled local interval centered at $t_0$.  Choose
$u_{0,n}\in H^1$ with $u_{0,n}\to u(t_0)$ in $L^2$ and with uniformly
bounded $L^2$ norms.  The local construction gives all approximating
solutions on the same interval.  By \cref{lem:H1-persistence} they are
$H^1$ there, and by \cref{lem:mass-H1} their masses are constant.  Local
flow stability gives $u_n\to u$ in $C_tL^2_x$, so passing to the limit proves
\eqref{eq:mass-conservation} on that interval.  Restarting the same argument
at any point in the lifespan and chaining overlapping intervals proves the
identity throughout the lifespan.
\end{proof}

\begin{theorem}[Global $L^2$ well-posedness]\label{thm:GWP}
For every $u_0\in L^2(\R)$ and every $\sigma\in\R$, there is a unique global
mild solution
\[
 u\in C(\R;L^2)\cap L^6_{\loc}(\R^2).
\]
The flow map $\Phi_t:u_0\mapsto u(t)$ is continuous, obeys the group law, and
is gauge covariant:
\begin{equation}
  \Phi_t(\zeta u_0)=\zeta\Phi_t(u_0)
  \qquad(|\zeta|=1).                               \label{eq:gauge-cov}
\end{equation}
\end{theorem}

\begin{proof}
The local lifespan in \cref{thm:LWP} depends only on an upper bound for the
$L^2$ norm.  By \cref{prop:mass-L2}, that norm remains equal to
$\norm{u_0}_2$.  The local construction can therefore be iterated indefinitely
in both time directions.  Uniqueness makes overlapping pieces agree and gives
the group law.  Continuity follows by iterating the local Lipschitz estimate
on any fixed compact time interval.  Finally, if $u$ solves the equation, then
$\zeta u$ solves it with initial data $\zeta u_0$ whenever $|\zeta|=1$;
uniqueness gives \eqref{eq:gauge-cov}.
\end{proof}

\section{Persistence of all integer derivatives}

For an integer $k\ge0$ and a bounded interval $J$, write
\[
 \mathcal X_k(J)=C(J;H^k(\R))\cap L^6(J;W^{k,6}(\R)),
 \qquad
 \norm u_{\mathcal X_k(J)}
 =\sum_{j=0}^k\norm{\partial_x^ju}_{\mathcal X(J)}.
\]
This is a Banach space because each weak derivative is a closed operator and
$\mathcal X(J)$ is Banach.

\begin{lemma}[Cubic Leibniz decomposition and lower-order bound]
\label{lem:cubic-Hk-leibniz}
Let $k\ge1$.  For a smooth complex-valued $u$,
\begin{equation}
 \partial_x^k(|u|^2u)=DN(u)[\partial_x^ku]+R_k(u), \label{eq:top-derivative-split}
\end{equation}
where
\begin{equation}
 R_k(u)=
 \sum_{\substack{a+b+c=k\\0\le a,b,c\le k-1}}
 \frac{k!}{a!b!c!}
 (\partial_x^au)(\partial_x^b\overline u)(\partial_x^cu).       \label{eq:Rk-explicit}
\end{equation}
If $u\in\mathcal X_{k-1}(J)$, the right side of
\eqref{eq:Rk-explicit} is well-defined in $L^{6/5}(J\times\R)$ and
\begin{equation}
 \norm{R_k(u)}_{L^{6/5}_{t,x}(J)}
 \le C_k|J|^{1/2}\norm u_{\mathcal X_{k-1}(J)}^3.               \label{eq:Rk-bound}
\end{equation}
The map $R_k:\mathcal X_{k-1}(J)\to L^{6/5}(J\times\R)$ is locally
Lipschitz.
\end{lemma}

\begin{proof}
The multinomial Leibniz rule gives
\[
 \partial_x^k(u\overline u u)
 =\sum_{a+b+c=k}\frac{k!}{a!b!c!}
   (\partial_x^au)(\partial_x^b\overline u)(\partial_x^cu).
\]
The three terms in which one index is $k$ are
\[
 |u|^2\partial_x^ku
 +u^2\partial_x^k\overline u
 +|u|^2\partial_x^ku
 =DN(u)[\partial_x^ku],
\]
and every remaining index is at most $k-1$, which proves
\eqref{eq:top-derivative-split}--\eqref{eq:Rk-explicit}.

For one remaining monomial choose any one factor for
$L^\infty_tL^2_x$ and the other two for $L^6_tL^6_x$.  At each time the
spatial reciprocal exponents add to
$1/2+1/6+1/6=5/6$, while in time the two $L^6$ factors give an $L^3$
function.  The finite-measure embedding
$L^3(J)\hookrightarrow L^{6/5}(J)$ has norm $|J|^{1/2}$.  Hence
\[
 \norm{f_1f_2f_3}_{L^{6/5}_{t,x}}
 \le |J|^{1/2}\norm{f_1}_{L^\infty_tL^2_x}
                   \norm{f_2}_{L^6_tL^6_x}
                   \norm{f_3}_{L^6_tL^6_x}.
\]
Summing the finitely many monomials proves \eqref{eq:Rk-bound}.  Expanding
a difference of two cubic monomials into three terms, with one factor
replaced by its difference, proves local Lipschitz continuity by the same
estimate.
\end{proof}

\begin{proposition}[Global $H^k$ persistence]
\label{prop:Hk-persistence}
Let $k\ge1$ be an integer.  If $u_0\in H^k(\R)$, then on every compact time
interval
\[
 u\in C_tH^k_x\cap L^6_tW^{k,6}_x.
\]
If $u_0$ is Schwartz, then $u$ is smooth in $(t,x)$ on every compact time
interval.
\end{proposition}

\begin{proof}
The case $k=1$ is \cref{lem:H1-persistence}.  Its local interval can be
chosen from the mass-controlled construction in \cref{thm:LWP}; the length
and the bound for $\norm u_{\mathcal X(J)}$ depend only on the conserved
$L^2$ norm.  Finitely many such intervals cover every compact time interval.

Assume now that persistence has been proved through order $k-1$, where
$k\ge2$, and work on one of those mass-controlled intervals $J$ centered at
$t_0$.  By reducing its length by a fixed mass-dependent factor if necessary,
we retain
\begin{equation}
 C|\sigma||J|^{1/2}\norm u_{\mathcal X(J)}^2\le\frac14.          \label{eq:Hk-linear-smallness}
\end{equation}
The induction hypothesis and \cref{lem:cubic-Hk-leibniz} give
$R_k(u)\in L^{6/5}_{t,x}(J)$.  Define the affine map on
$\mathcal X(J)$ by
\begin{align}
 \Theta_k(z)(t)=S(t-t_0)\partial_x^ku(t_0)
 -i\sigma\int_{t_0}^tS(t-s)
 \bigl(DN(u(s))[z(s)]+R_k(u)(s)\bigr)\,\dd s.       \label{eq:Hk-affine-map}
\end{align}
The inhomogeneous Strichartz estimate, \eqref{eq:DN-bound}, and
\eqref{eq:Hk-linear-smallness} make the linear part a contraction.  Hence
there is a unique $z_k\in\mathcal X(J)$ satisfying
\begin{equation}
 z_k=\Theta_k(z_k).                                  \label{eq:zk-fixed}
\end{equation}

We identify this candidate with the $k$th weak derivative.  Put
$y=\partial_x^{k-1}u$, which belongs to $\mathcal X(J)$ by induction, and
for $h\ne0$ set
\[
 q_h=\frac{\tau_hy-y}{h},
 \qquad q_{0,h}=\frac{\tau_h\partial_x^{k-1}u(t_0)
                         -\partial_x^{k-1}u(t_0)}h.
\]
Then $q_{0,h}\to\partial_x^ku(t_0)$ in $L^2$.  Differentiate the mild
equation $k-1$ times, translate it, subtract, and divide by $h$.  Applying
the scalar fundamental theorem of calculus to every translated cubic
monomial yields an equation
\begin{equation}
 q_h=S(t-t_0)q_{0,h}
 -i\sigma\int_{t_0}^tS(t-s)
       \bigl(A_h(s)[q_h(s)]+R_{k,h}(s)\bigr)\,\dd s.              \label{eq:Hk-dq-equation}
\end{equation}
To make the coefficients explicit, differentiate the mild equation only
$k-1$ times before taking the quotient.  By
\eqref{eq:top-derivative-split} at order $k-1$,
\[
 \partial_x^{k-1}N(u)=DN(u)[y]+R_{k-1}(u).
\]
Consequently one may take
\begin{align*}
 A_h[w]&=DN(\tau_hu)[w],\\
 R_{k,h}
 &=\frac{DN(\tau_hu)-DN(u)}h[y]
   +\frac{\tau_hR_{k-1}(u)-R_{k-1}(u)}h.
\end{align*}
Every term in $R_{k,h}$ is a finite product of derivatives of order at most
$k-1$ and a translation difference quotient of a derivative of order at most
$k-2$; no occurrence of $q_h$ remains there.  Termwise differentiation of
the explicit multinomial formula shows that its formal limit is precisely
$R_k(u)$.  Quantitatively,
\begin{align}
 \sup_{\norm w_{\mathcal X(J)}\le1}
 \norm{A_h[w]-DN(u)[w]}_{L^{6/5}_{t,x}(J)}
 &\longrightarrow0,                                \label{eq:Ah-to-DN}\\
 \norm{R_{k,h}-R_k(u)}_{L^{6/5}_{t,x}(J)}
 &\longrightarrow0.                                \label{eq:Rkh-to-Rk}
\end{align}
Indeed, the induction hypothesis places every derivative through order
$k-1$ in both $C_tL^2_x$ and $L^6_tL^6_x$.  Translation is strongly
continuous in both spaces.  If $j\le k-2$, then
\[
 \frac{\tau_h\partial_x^ju-\partial_x^ju}{h}
 \longrightarrow\partial_x^{j+1}u
\]
in $L^6_tL^6_x$ by
\cref{lem:Wp-difference-quotient}, applied to the product spacetime
norm.  The
convergence in $C_tL^2_x$ is uniform in time: the compact set
$\{\partial_x^{j+1}u(t):t\in J\}\subset L^2$ is acted on by the uniformly
bounded Fourier multipliers $(\e^{ih\xi}-1)/(ih\xi)$, which converge strongly
to one.  Expand every coefficient or product difference one factor at a time
and apply the three-factor estimate in \eqref{eq:Rk-bound}.  This proves
\eqref{eq:Ah-to-DN}--\eqref{eq:Rkh-to-Rk} without assuming a derivative of
order $k$.

For small $h$, the linear operator in \eqref{eq:Hk-dq-equation} has
contraction norm at most $1/2$, by
\eqref{eq:Hk-linear-smallness} and translation continuity.  Subtract
\eqref{eq:zk-fixed} from \eqref{eq:Hk-dq-equation}, apply the inhomogeneous
estimate, and absorb the $q_h-z_k$ term.  The result is
\[
 \norm{q_h-z_k}_{\mathcal X(J)}
 \le C\norm{q_{0,h}-\partial_x^ku(t_0)}_2
   +C\norm{R_{k,h}-R_k(u)}_{6/5}
   +o(1)\norm{z_k}_{\mathcal X(J)}\longrightarrow0.
\]
Translation difference quotients of $y=\partial_x^{k-1}u$ converge to
$\partial_xy=\partial_x^ku$ in $\D'(J\times\R)$, so the strong
$\mathcal X$ limit satisfies $z_k=\partial_x^ku$.  This proves $u\in\mathcal X_k(J)$.  Finite iteration
on the mass-controlled intervals proves the assertion on every compact time
interval.

If the datum is Schwartz, it belongs to every $H^k$, so the preceding result
gives all spatial derivatives locally in time.  The equation
$u_t=i u_{xx}-i\sigma|u|^2u$ then gives the first time derivative with
arbitrarily high spatial regularity; differentiating the equation repeatedly
in time gives every mixed derivative.  Thus $u$ is smooth.  No persistence
of spatial moments is required.
\end{proof}

\section{Endpoint spacetime regularity}
The next estimate is stronger in the spatial variable than the fixed-point
norm.  It is also the source of almost-everywhere continuous spatial slices.

\begin{proposition}[Endpoint bound and integrable forcing]\label{prop:endpoint}
Let $I$ be bounded and let $u$ be a global $L^2$ solution.  Then
\begin{align}
 u&\in L^4(I;L^\infty(\R)),                         \label{eq:L4Linfty}\\
 u&\in L^8(I;L^4(\R)),                              \label{eq:L8L4}\\
 |u|^2u&\in L^1(I;L^2(\R)).                         \label{eq:forcing-L1L2}
\end{align}
Moreover, if $u_n$ are the solutions with data $u_{0,n}\to u_0$ in $L^2$,
then on every bounded $I$
\[
 u_n\to u\quad\text{in }L^4_tL^\infty_x,
 \qquad
 |u_n|^2u_n\to|u|^2u\quad\text{in }L^1_tL^2_x.
\]
\end{proposition}

\begin{proof}
First,
\begin{equation}
 \norm{|u|^2u}_{L^1_tL^2_x(I)}
 =\int_I\norm{u(t)}_6^3\,\dd t
 \le |I|^{1/2}\norm u_{L^6(I\times\R)}^3.           \label{eq:L1L2-from-L6}
\end{equation}
Use the Duhamel formula and the homogeneous endpoint estimate
$(q,r)=(4,\infty)$.  For each fixed $s$, restriction of the time variable can
only decrease the scalar $L^4_tL^\infty_x$ norm, so
\cref{lem:scalar-mixed-minkowski} gives
\begin{align*}
 \left\|\int_{t_0}^{t}S(t-s)F(s)\,\dd s\right\|_{L^4_tL^\infty_x(I)}
 &\le\int_I
   \norm{\1_{\{s\text{ lies between }t_0\text{ and }t\}}
             S(t-s)F(s)}_{L^4_tL^\infty_x(I)}\,\dd s\\
 &\le C\int_I\norm{F(s)}_2\,\dd s.
\end{align*}
With $F=|u|^2u$, \eqref{eq:L1L2-from-L6} proves \eqref{eq:L4Linfty}.
The $L^8_tL^4_x$ estimate follows without invoking an interpolation-space
identification.  Pointwise in time,
\[
 \norm{u(t)}_4^4\le\norm{u(t)}_2^2\norm{u(t)}_\infty^2.
\]
Square, integrate in time, and use the $L^\infty_tL^2_x$ and
$L^4_tL^\infty_x$ bounds to obtain \eqref{eq:L8L4}.

For convergence of the nonlinear forcing, use
\[
 \norm{|u_n|^2u_n-|u|^2u}_{L^2_x}
 \le C(\norm{u_n}_6^2+\norm u_6^2)\norm{u_n-u}_6.
\]
The first factor belongs uniformly to $L^3_t$, the last converges in $L^6_t$,
and an additional H\"older inequality on the finite interval yields
convergence in $L^1_t$.  The endpoint Duhamel estimate just proved then gives
convergence in $L^4_tL^\infty_x$.  Since the cubic forcing belongs to both
$L^{6/5}_{t,x}$ and $L^1_tL^2_x$, the compatibility assertion of
\cref{lem:retarded-continuous} shows that every mild formula used from this
point onward is an ordinary $L^2$-valued Bochner integral.
\end{proof}

\begin{corollary}[Continuous spatial representatives at almost every time]
\label{cor:C0-slices}
For every bounded interval $I$, there is a full-measure set $E_I\subset I$
such that, for $t\in E_I$, the $L^2$ class $u(t)$ has a representative in
$C_0(\R)$.
\end{corollary}

\begin{proof}
Choose Schwartz initial data $u_{0,n}\to u_0$ and let $u_n$ be the
corresponding solutions.  By \cref{prop:Hk-persistence}, every
$u_n(t)$ belongs to $H^1$.  The explicit one-dimensional representative
theorem \cref{lem:H1-C0}, including its tail estimates, therefore gives
$u_n(t,\cdot)\in C_0(\R)$ for every time.

By \cref{prop:endpoint}, $u_n\to u$ in the scalar space
$L^4(I;L^\infty_x)$.  Apply \cref{lem:mixed-subsequence-uniform}.  For almost
every $t$, a subsequence is uniformly Cauchy in $x$ and hence converges to a
function in $C_0$.  At every time the same sequence converges to $u(t)$ in
$L^2$ because the flow convergence is in $C_tL^2_x$.  Therefore
\cref{lem:C0-limit-identification} identifies the uniform limit with a
representative of $u(t)$.  A countable exhaustion of the time axis and an
intersection of the resulting full-measure sets give one common set.
\end{proof}

\begin{checkpoint}
The global flow, its stability, mass conservation, gauge covariance, and all
spacetime norms later used in the proof have been established.  In
particular, the cubic forcing is locally integrable with values in $L^2$.
\end{checkpoint}

\chapter{Local smoothing by one half derivative}\label{chap:smoothing}

\section{The homogeneous identity}
The local smoothing gain is proved directly from the Fourier representation.

\begin{lemma}[Pointwise-in-space smoothing]\label{lem:point-smoothing}
For every $f\in L^2(\R)$ and every $x\in\R$,
\begin{equation}
 \int_\R\big| |D|^{1/2}S(t)f(x)\big|^2\,\dd t
 \le C\norm f_2^2.                                 \label{eq:point-smoothing}
\end{equation}
The constant is independent of $x$.
\end{lemma}

\begin{proof}
Begin with $f\in\Ssch$.  Write
\[
 |D|^{1/2}S(t)f(x)
 =(2\pi)^{-1/2}\int_\R
    \e^{ix\xi-it\xi^2}|\xi|^{1/2}\widehat f(\xi)\,\dd\xi.
\]
Split into positive and negative frequencies and put $\lambda=\xi^2$.  Up to
one fixed normalization constant, the expression is the inverse time Fourier
transform of
\begin{equation}
 B_xf(\lambda)=\1_{\{\lambda>0\}}\lambda^{-1/4}
 \left(\e^{ix\sqrt\lambda}\widehat f(\sqrt\lambda)
      +\e^{-ix\sqrt\lambda}\widehat f(-\sqrt\lambda)\right).
 \label{eq:smoothing-Bx}
\end{equation}
Plancherel in $t$ and $(a+b)^2\le2a^2+2b^2$ give
\[
 \norm{B_xf}_{L^2_\lambda}^2
 \le C\int_0^\infty \lambda^{-1/2}
 \left(|\widehat f(\sqrt\lambda)|^2
      +|\widehat f(-\sqrt\lambda)|^2\right)\,\dd\lambda
 =C\norm f_2^2,
\]
where the final equality is the substitution $\lambda=\xi^2$.  The bound is
independent of $x$.

For general $f\in L^2$, formula \eqref{eq:smoothing-Bx} still defines an
$L^2_\lambda$ function, and hence by inverse time Plancherel a canonical
$A_xf\in L^2_t$.  The preceding estimate proves
$\norm{A_xf}_{L^2_t}\le C\norm f_2$.  Moreover
$x\mapsto B_xf$ is continuous into $L^2_\lambda$: the phase factors converge
pointwise when $x_n\to x$, and the displayed integrable majorant is
independent of $n$.  Thus $x\mapsto A_xf$ is continuous into $L^2_t$ and has
a jointly measurable representative.

Choose Schwartz $f_n\to f$ in $L^2$.  For every compact spatial set $K$,
\[
 \int_K\norm{A_x(f_n-f)}_{L^2_t}^2\,\dd x
 \le C|K|\norm{f_n-f}_2^2\longrightarrow0.
\]
For the Schwartz approximants, $A_xf_n=|D|^{1/2}S(t)f_n(x)$.  The latter
converges in distributions to the Fourier-multiplier distribution
$|D|^{1/2}S(t)f$, while the preceding estimate gives the local
$L^2_{t,x}$ limit $A_xf$.  The two limits coincide.  We use this canonical
representative for the expression in \eqref{eq:point-smoothing}; its norm
bound is exactly the asserted estimate.
\end{proof}

\begin{corollary}[Localized homogeneous smoothing]\label{cor:hom-smoothing}
For $\chi\in C_c^\infty(\R)$ and a bounded interval $I$,
\begin{equation}
 \norm{\chi S(t)f}_{L^2(I;H^{1/2}_x)}
 \le C_{\chi,I}\norm f_2.                          \label{eq:hom-local-smoothing}
\end{equation}
\end{corollary}

\begin{proof}
Integrating \eqref{eq:point-smoothing} against $|\chi(x)|^2$ gives
\[
 \norm{\chi|D|^{1/2}S(t)f}_{L^2_{t,x}(\R^2)}
 \le C\norm\chi_2\norm f_2.
\]
It remains to commute $|D|^{1/2}$ through $\chi$.  In Fourier variables the
commutator has kernel
\[
 (|\xi|^{1/2}-|\eta|^{1/2})\widehat\chi(\xi-\eta).
\]
Since $||\xi|^{1/2}-|\eta|^{1/2}|\le|\xi-\eta|^{1/2}$ and
$|\theta|^{1/2}\widehat\chi(\theta)\in L^1$, Young's inequality gives
\[
 \norm{[|D|^{1/2},\chi]g}_2\le C_\chi\norm g_2.
\]
The $L^2$ part of the $H^{1/2}$ norm contributes
$|I|^{1/2}\norm\chi_\infty\norm f_2$ by unitarity.
\end{proof}

\section{Duhamel smoothing}

\begin{proposition}[Local smoothing for forced solutions]\label{prop:forced-smoothing}
Let $I$ be bounded, $t_0\in I$, $f\in L^2$, and
$F\in L^1(I;L^2)$.  Then
\[
 w(t)=S(t-t_0)f+\int_{t_0}^tS(t-s)F(s)\,\dd s
\]
satisfies, for every $\chi\in C_c^\infty(\R)$,
\begin{equation}
 \norm{\chi w}_{L^2(I;H^{1/2})}
 \le C_{\chi,I}\left(\norm f_2+\norm F_{L^1_tL^2_x(I)}\right). \label{eq:forced-smoothing}
\end{equation}
\end{proposition}

\begin{proof}
The homogeneous term is \cref{cor:hom-smoothing}.  For the Duhamel term,
Minkowski and time translation give
\begin{align*}
 &\left\|\chi\int_{t_0}^tS(t-s)F(s)\,\dd s\right\|_{L^2_tH^{1/2}_x(I)}\\
 &\qquad\le\int_I
 \norm{\chi\1_{\{s\text{ between }t_0\text{ and }t\}}
       S(t-s)F(s)}_{L^2_tH^{1/2}_x(I)}\,\dd s
 \le C_{\chi,I}\int_I\norm{F(s)}_2\,\dd s.
\end{align*}
Restriction of the time domain again only decreases the norm.
\end{proof}

\begin{corollary}[Local smoothing for the cubic flow]\label{cor:NLS-smoothing}
Every global $L^2$ solution satisfies
\begin{equation}
  \chi u\in L^2(I;H^{1/2}(\R))
  \qquad(\chi\in C_c^\infty,\ I\Subset\R).          \label{eq:NLS-local-smoothing}
\end{equation}
If $u_n$ are smooth-data solutions converging to $u$ in the $L^2$ flow, then
\begin{equation}
  \chi u_n\to\chi u
  \quad\text{in }L^2(I;H^{1/2}).                   \label{eq:smoothing-convergence}
\end{equation}
\end{corollary}

\begin{proof}
Apply \cref{prop:forced-smoothing} to the mild equation and use
\eqref{eq:forcing-L1L2}.  For differences, the initial data converge in
$L^2$ and the cubic forcings converge in $L^1_tL^2_x$ by
\cref{prop:endpoint}.
\end{proof}

\section{A useful countable full-measure set}
Fix an increasing sequence of compact time intervals $I_m$ and compact
spatial cutoffs $\chi_m$ with $\chi_m=1$ on $[-m,m]$.  From
\eqref{eq:NLS-local-smoothing}, for almost every $t$ we have
$\chi_m u(t),\chi_m v(t)\in H^{1/2}$ for every $m$.  Intersect this set with
the full-measure set in \cref{cor:C0-slices}.  We obtain one set
$E\subset\R$ of full measure such that, for every $t\in E$,
\begin{equation}
 u(t),v(t)\in H^{1/2}_{\loc}(\R)\cap C_0(\R).       \label{eq:good-slices}
\end{equation}
Later we shall shrink $E$ by countably many distributional identities without
changing its measure.

\begin{checkpoint}
The exact critical regularity needed for the current and diagonal traces is
now available, and smooth approximants converge in that norm.
\end{checkpoint}

\chapter{Critical \texorpdfstring{$H^{1/2}$}{H1/2} calculus}\label{chap:critical-calculus}

The endpoint proof repeatedly multiplies an $H^{1/2}$ function by a bounded
$H^{1/2}$ function and composes with Lipschitz maps.  At the critical index,
these operations are valid only when the $L^\infty$ information is retained.
This chapter records the exact estimates and proves the approximation facts
used at the common zero set.

\section{The product estimate}

\begin{lemma}[$H^{1/2}\cap L^\infty$ is an algebra]\label{lem:Hhalf-algebra}
If $f,g\in H^{1/2}(\R)\cap L^\infty(\R)$, then $fg\in H^{1/2}$ and
\begin{equation}
 \norm{fg}_{H^{1/2}}
 \le C\left(\norm f_\infty\norm g_{H^{1/2}}
           +\norm g_\infty\norm f_{H^{1/2}}\right). \label{eq:Hhalf-product}
\end{equation}
The corresponding statement is local after multiplication by a smooth
cutoff.
\end{lemma}

\begin{proof}
The $L^2$ part follows from either factor in $L^\infty$.  For the Gagliardo
seminorm,
\[
 f(x)g(x)-f(y)g(y)
 =f(x)(g(x)-g(y))+g(y)(f(x)-f(y)).
\]
Square, use $(a+b)^2\le2a^2+2b^2$, divide by $|x-y|^2$, and integrate.
The local statement follows by applying the global estimate to cutoffs and
using \cref{lem:smooth-multiplier}.
\end{proof}

\begin{lemma}[Lipschitz composition]\label{lem:Lip-composition}
Let $F:\C\to\C$ be Lipschitz and satisfy $F(0)=0$.  If
$f\in H^{1/2}(\R)$, then $F\circ f\in H^{1/2}$ and
\[
 \norm{F(f)}_{H^{1/2}}\le C\operatorname{Lip}(F)\norm f_{H^{1/2}}.
\]
If $F$ is merely Lipschitz on the range of a locally bounded function, the
same conclusion holds locally after subtracting the constant $F(0)$.
\end{lemma}

\begin{proof}
The pointwise inequalities
$|F(f(x))|\le\operatorname{Lip}(F)|f(x)|$ and
$|F(f(x))-F(f(y))|\le\operatorname{Lip}(F)|f(x)-f(y)|$ control the $L^2$
norm and the Gagliardo seminorm.
\end{proof}

\section{Mollification in both critical and uniform norms}

\begin{lemma}[Simultaneous approximation]\label{lem:simultaneous-mollification}
Let $f\in H^{1/2}(\R)\cap C_0(\R)$.  If $\varphi_\delta$ is a standard smooth
approximate identity, then
\[
 \varphi_\delta*f\to f
 \quad\text{in }H^{1/2}(\R)\text{ and in }L^\infty(\R).
\]
The same is true locally after inserting a cutoff.
\end{lemma}

\begin{proof}
With the unitary Fourier normalization, convolution by
$\varphi_\delta$ has multiplier
$m_\delta(\xi)=(2\pi)^{1/2}\widehat\varphi(\delta\xi)$.
It is uniformly bounded and converges pointwise to one because
$\int\varphi=1$.  Dominated convergence therefore gives strong
convergence in $H^{1/2}$.
A continuous function vanishing at infinity is uniformly continuous, so the
standard approximate-identity argument gives uniform convergence.  For the
local version, first multiply by a cutoff which is identically one on the
compact set of interest.
\end{proof}

\section{The zero-set regularizers}
For $\eps>0$ define
\begin{equation}
 T_\eps(z)=\frac{\eps z}{|z|^2+\eps},
 \qquad z\in\C.                                    \label{eq:Teps-def}
\end{equation}

\begin{lemma}[Uniform Lipschitz and strong critical convergence]
\label{lem:Teps}
The maps $T_\eps:\C\to\C$ are uniformly Lipschitz, satisfy
$T_\eps(0)=0$, and obey
\begin{equation}
 |T_\eps(z)|\le\frac{\sqrt\eps}{2}.                \label{eq:Teps-sup}
\end{equation}
If $f\in H^{1/2}_{\loc}(\R)\cap C(\R)$ is locally bounded, then
\begin{equation}
 T_\eps(f)\longrightarrow0
 \quad\text{strongly in }H^{1/2}_{\loc}(\R).       \label{eq:Teps-strong}
\end{equation}
\end{lemma}

\begin{proof}
Viewed as a map on $\R^2$, $T_\eps$ is radial.  Its tangential derivative is
$\eps/(r^2+\eps)$ and its radial derivative is
$\eps(\eps-r^2)/(r^2+\eps)^2$; both have absolute value at most one.  The
maximum of $\eps r/(r^2+\eps)$ occurs at $r=\sqrt\eps$ and equals
$\sqrt\eps/2$.  We also have the domination
\begin{equation}
 |T_\eps(z)|=\frac{\eps|z|}{|z|^2+\eps}\le |z|.    \label{eq:Teps-domination}
\end{equation}

After localizing, assume $f\in H^{1/2}\cap C_0$.  For every fixed $z$,
$T_\eps(z)\to0$ as $\eps\downarrow0$.  By
\eqref{eq:Teps-domination}, $|T_\eps(f)|^2\le|f|^2$, so dominated convergence
proves convergence in $L^2(\R)$.  For the Gagliardo seminorm, pointwise
convergence and the uniform Lipschitz bound give
\[
 \frac{|T_\eps(f(x))-T_\eps(f(y))|^2}{|x-y|^2}
 \le \frac{|f(x)-f(y)|^2}{|x-y|^2}.
\]
The right side is integrable by
\cref{thm:fourier-gagliardo}, and dominated convergence on $\R^2$ proves
convergence of the seminorm.  This yields strong $H^{1/2}$ convergence.  The
local statement follows by multiplying $f$ by a cutoff equal to one on the
compact set under consideration.
\end{proof}

\section{Quadratic-form distributions}
If $f\in H^{1/2}_{\loc}$, define the local distribution
$\overline f f_x$ by
\begin{equation}
 \pair{\overline f f_x}{\phi}
 :=\pair{f_x}{\phi\overline f}_{H^{-1/2},H^{1/2}},
 \qquad \phi\in C_c^\infty(\R).                   \label{eq:quad-form}
\end{equation}
The pairing is complex bilinear, extending $\int T\phi$; it is not the
sesquilinear Hilbert inner product.  By
\cref{lem:smooth-multiplier}, the test factor belongs to $H^{1/2}$, and the
definition is independent of the cutoff used to globalize $f$.  Notice that
local boundedness is not needed for this quadratic form; it will enter only
later, when the zero-set regularizers are multiplied into test factors.

\begin{definition}[Multiplier test algebra]\label{def:multiplier-test-algebra}
For a compact interval $K\subset\R$, let $\mathcal A(K)$ be the space of
functions $\psi\in H^{1/2}(\R)\cap L^\infty(\R)$ supported in $K$, with norm
\[
 \norm\psi_{\mathcal A}=\norm\psi_{H^{1/2}}+\norm\psi_\infty.
\]
Smooth compactly supported functions are dense in $\mathcal A(K)$ for the
simultaneous topology whenever $\psi$ has a continuous representative; the
mollification statement is \cref{lem:simultaneous-mollification}.
\end{definition}

\begin{lemma}[Extension of the logarithmic quadratic form]
\label{lem:quadratic-form-extension}
Let $f\in H^{1/2}(\R)\cap L^\infty(\R)$.  The map initially defined on
smooth $\psi$ by
\[
 \mathcal Q_f(\psi)=\pair{f_x}{\psi\overline f}_{H^{-1/2},H^{1/2}}
\]
extends continuously to $\mathcal A(K)$, and
\begin{equation}
 |\mathcal Q_f(\psi)|
 \le C\norm f_{H^{1/2}}
 \left(\norm\psi_\infty\norm f_{H^{1/2}}
       +\norm f_\infty\norm\psi_{H^{1/2}}\right).  \label{eq:quadratic-extension-bound}
\end{equation}
If $f_n\to f$ in both $H^{1/2}$ and $L^\infty$, and
$\psi_n\to\psi$ in both norms, then
$\mathcal Q_{f_n}(\psi_n)\to\mathcal Q_f(\psi)$.
The same assertions hold locally after one fixed cutoff.
\end{lemma}

\begin{proof}
The derivative map is bounded from $H^{1/2}$ to $H^{-1/2}$.  By the algebra
estimate,
\[
 \norm{\psi\overline f}_{H^{1/2}}
 \le C\left(\norm\psi_\infty\norm f_{H^{1/2}}
       +\norm f_\infty\norm\psi_{H^{1/2}}\right).
\]
Pairing the two Sobolev spaces proves \eqref{eq:quadratic-extension-bound}.
For convergence, add and subtract
$\pair{(f_n)_x}{\psi_n\overline f}$ and
$\pair{(f_n)_x}{\psi\overline f}$, then apply the same product estimate to
each difference.  Uniform boundedness of the convergent sequences completes
the proof.  Localization is obtained by inserting a cutoff equal to one near
$K$; two choices give the same value because the distributional derivative is
local.
\end{proof}

\begin{lemma}[Stability of quadratic forms]\label{lem:quad-stability}
Suppose on a compact set $K$ that
$f_n\to f$ in $H^{1/2}$ after multiplication by one cutoff equal to one near
$K$.  Then
\[
 \overline{f_n}(f_n)_x\longrightarrow\overline f f_x
 \quad\text{in }\D'(K).
\]
The same conclusion holds after integration in time when the localized
$H^{1/2}$ convergence is in $L^2_t$.
\end{lemma}

\begin{proof}
For a test function $\phi$ supported in $K$, subtract the two pairings and
write
\begin{align*}
 \pair{(f_n)_x-f_x}{\phi\overline{f_n}}
 +\pair{f_x}{\phi(\overline{f_n}-\overline f)}.
\end{align*}
By boundedness of multiplication by the fixed smooth function $\phi$, the
absolute value is at most
\[
 C_\phi\norm{f_n-f}_{H^{1/2}}
 \bigl(\norm{f_n}_{H^{1/2}}+\norm f_{H^{1/2}}\bigr),
\]
after localization.  This tends to zero.  In time, integrate the same bound
and use Cauchy--Schwarz for the two $H^{1/2}$ factors.
\end{proof}

\begin{lemma}[Real part of the logarithmic form]\label{lem:real-log-form}
For $f\in H^{1/2}_{\loc}$,
\begin{equation}
  \Rea(\overline f f_x)=\frac12\partial_x|f|^2
  \quad\text{in }\D'(\R).                         \label{eq:real-log-form}
\end{equation}
\end{lemma}

\begin{proof}
Mollify after a cutoff.  For smooth $f_n$ the identity is the product rule.
The left side converges by \cref{lem:quad-stability}.  The right side
converges because $f_n\to f$ in $L^2_{\loc}$ and
\[
 \norm{|f_n|^2-|f|^2}_{L^1(K)}
 \le \norm{f_n-f}_{L^2(K)}
     \bigl(\norm{f_n}_{L^2(K)}+\norm f_{L^2(K)}\bigr).
\]
\end{proof}

\begin{checkpoint}
All critical products, compositions, and limiting operations needed below are
now justified by explicit estimates.  In particular, no division at a zero
will be used.
\end{checkpoint}

\chapter{The density determines the endpoint current}\label{chap:current}

\section{Density equality at every time}

\begin{lemma}[Continuity of the density map]\label{lem:density-continuity}
If $f\in C(I;L^2(\R))$, then $t\mapsto|f(t)|^2$ is continuous from $I$ to
$L^1(\R)$.
\end{lemma}

\begin{proof}
The elementary factorization gives
\[
 \norm{|f(t)|^2-|f(s)|^2}_{L^1}
 \le\norm{f(t)-f(s)}_2\big(\norm{f(t)}_2+\norm{f(s)}_2\big).
\]
\end{proof}

\begin{corollary}[All-time equality of densities]\label{cor:all-time-density}
Under the hypothesis \eqref{eq:equal-modulus-main},
\begin{equation}
 |u(t)|^2=|v(t)|^2\quad\text{in }L^1(\R)
 \quad\text{for every }t\in\R.                    \label{eq:all-time-density}
\end{equation}
\end{corollary}

\begin{proof}
Fubini gives equality in $L^1$ for almost every time.  The difference of the
two density maps is continuous by \cref{lem:density-continuity}; a continuous
map which is zero almost everywhere is zero everywhere.
\end{proof}

\section{The smooth continuity equation}
For a smooth solution set
\[
 \rho_u=|u|^2,
 \qquad j_u=\Ima(\overline u\,u_x).
\]

\begin{lemma}[Continuity equation]\label{lem:continuity-smooth}
Every $H^1$ solution satisfies
\begin{equation}
  \partial_t\rho_u+2\partial_xj_u=0
  \quad\text{in }\D'(I\times\R).                  \label{eq:continuity-eq}
\end{equation}
Consequently, if $\phi\in C_c^\infty(\R)$ and
$a(x)=\int_{-\infty}^x\phi(y)\,\dd y$, then
\begin{equation}
 \frac{\dd}{\dd t}\int_\R a(x)|u(t,x)|^2\,\dd x
 =2\int_\R\phi(x)j_u(t,x)\,\dd x                 \label{eq:localized-mass-smooth}
\end{equation}
in the distributional sense in time.
\end{lemma}

\begin{proof}
We justify the product rule at the stated $H^1$ regularity.  If
$b\in C_c^\infty(\R)$ is real, multiplication by $b$ is bounded on $H^1$.
The time-mollification proof of \cref{lem:gelfand-chain-rule}, applied to the
bounded sesquilinear form
$B_b(f,g)=\int b f\overline g$, gives
\[
 \frac{\dd}{\dd t}\int b|u|^2
 =2\Rea\pair{u_t}{bu}_{H^{-1},H^1}
\]
in distributions in time.  Substitute
$u_t=i u_{xx}-i\sigma|u|^2u$.  The nonlinear pairing is purely imaginary,
while weak integration by parts gives
\begin{align*}
 2\Rea\pair{i u_{xx}}{bu}
 &=2\Ima\int \overline u\,u_x b'
 =2\int b'\Ima(\overline u\,u_x).
\end{align*}
Testing also in time proves
$\partial_t|u|^2+2\partial_x\Ima(\overline u u_x)=0$ in spacetime
distributions.

For the primitive $a(x)=\int_{-\infty}^x\phi$, insert a compact cutoff
$b_R=a\chi_R$ in the preceding identity.  The term with $a\chi_R'$ tends to
zero because $a$ is bounded, $u\,u_x\in L^1$, and the cutoffs may be chosen
with uniformly bounded derivative converging pointwise to zero.  The term
with $a'\chi_R=\phi\chi_R$ converges to the right side of
\eqref{eq:localized-mass-smooth}.  This proves the localized identity.
\end{proof}

\section{Endpoint current as a spacetime distribution}
For a solution with only $L^2$ data, define for almost every time
\begin{equation}
 \pair{\mathfrak j_u(t)}{\phi}
 =\Ima\pair{u_x(t)}{\phi\overline{u(t)}}_{H^{-1/2},H^{1/2}}. \label{eq:endpoint-current-def}
\end{equation}
The local smoothing norm makes this an $L^1_{\loc}$ family in time.
Indeed, after a cutoff containing $\supp\phi$,
\[
 |\pair{\mathfrak j_u(t)}{\phi}|
 \le C_\phi\norm{\chi u(t)}_{H^{1/2}}^2,
\]
and the right side is integrable in time.

\begin{proposition}[Weak localized mass identity]\label{prop:weak-localized-mass}
Let $u$ be a global $L^2$ solution.  For
$\eta\in C_c^\infty(I)$, $\phi\in C_c^\infty(\R)$, and the bounded primitive
$a$ above,
\begin{equation}
 -\int_I\eta'(t)\int_\R a(x)|u(t,x)|^2\,\dd x\,\dd t
 =2\int_I\eta(t)\pair{\mathfrak j_u(t)}{\phi}\,\dd t. \label{eq:weak-localized-mass}
\end{equation}
\end{proposition}

\begin{proof}
Take Schwartz initial data $u_{0,n}\to u_0$ and let $u_n$ be the smooth
solutions.  Identity \eqref{eq:weak-localized-mass} holds for $u_n$ by
\cref{lem:continuity-smooth}.  The left side converges because
$u_n\to u$ in $C_tL^2_x$ and $a$ is bounded.  The right side converges by
\cref{cor:NLS-smoothing,prop:endpoint,lem:quad-stability}.  Pass to the limit.
\end{proof}

\begin{theorem}[Endpoint recovery of the logarithmic derivative]
\label{thm:current-recovery}
Under the hypotheses of \cref{thm:main},
\begin{equation}
 \overline u\,u_x=\overline v\,v_x
 \quad\text{as a local quadratic-form distribution on }\R_t\times\R_x.
 \label{eq:log-derivative-equality}
\end{equation}
There is a full-measure set of times on which the same equality holds as a
spatial distribution.
\end{theorem}

\begin{proof}
Apply \cref{prop:weak-localized-mass} to $u$ and $v$.  By
\eqref{eq:all-time-density}, the left sides agree, so the current forms agree
for every tensor-product test function $\eta(t)\phi(x)$.

We record the continuity needed to pass to arbitrary tests.  If
$\Psi\in C_c^\infty(I\times\R)$ and $\chi$ is one on its spatial support,
then the algebra estimate with the smooth multiplier $\Psi(t,\cdot)$ gives
\begin{align}
 \left|\int_I
 \pair{u_x(t)}{\Psi(t)\overline{u(t)}}\,\dd t\right|
 &\le C_K\max_{|\alpha|\le1}\norm{\partial_x^\alpha\Psi}_\infty
       \norm{\chi u}_{L^2_tH^{1/2}_x}^2.          \label{eq:current-test-continuity}
\end{align}
The same bound holds for $v$.  By \cref{lem:tensor-test-density} with $M=1$, finite sums of separated
tests converge to $\Psi$ in the seminorm on the right of
\eqref{eq:current-test-continuity}.  Hence the imaginary parts of the two
forms agree against every spacetime test.  Their
real parts agree by \cref{lem:real-log-form} and equality of the densities.
This proves \eqref{eq:log-derivative-equality} as a spacetime distribution.

For slicing, fix compact intervals $I_m,K_m$ exhausting spacetime and a
cutoff $\chi_m$ equal to one near $K_m$.  For smooth $\phi$ supported in
$K_m$, the multiplier proof gives the concrete estimate
\begin{equation}
 |\Lambda_t(\phi)|
 \le C_m\bigl(\norm{\chi_m u(t)}_{H^{1/2}}^2
              +\norm{\chi_m v(t)}_{H^{1/2}}^2\bigr)
       \bigl(\norm\phi_\infty+\norm{\phi'}_\infty\bigr),
 \label{eq:sliced-current-test-bound}
\end{equation}
where
\[
 \Lambda_t(\phi)=
 \pair{u_x(t)}{\phi\overline{u(t)}}
 -\pair{v_x(t)}{\phi\overline{v(t)}}.
\]
The coefficient in parentheses is integrable on $I_m$.  Choose a countable
family $\{\phi_{m,n}\}_n\subset C_c^\infty(K_m)$ dense, for the
$C^1$ norm, in the smooth functions supported in $K_m$; piecewise-polynomial
approximants with rational data followed by a fixed smoothing construction
provide one.  The spacetime identity implies
$\int\eta(t)\Lambda_t(\phi_{m,n})\,\dd t=0$ for every scalar test $\eta$.
The scalar fundamental lemma gives
$\Lambda_t(\phi_{m,n})=0$ outside a null set depending on $(m,n)$.  Remove
the countable union of these null sets.  Estimate
\eqref{eq:sliced-current-test-bound} extends equality from the dense family
to every smooth test supported in $K_m$.  Finally take the countable union
of the compact exhaustions.  This produces one full-measure set on which the
spatial quadratic forms agree as distributions.  At the good times, their
continuous extensions to the multiplier algebra are then equal as well, by
\cref{lem:quadratic-form-extension}.
\end{proof}

\begin{checkpoint}
Equal modulus has now produced equal complex logarithmic-derivative forms at
the critical regularity.  The remaining issue is to convert that identity to
an unconjugated Wronskian without dividing by the common density.
\end{checkpoint}

\chapter{The critical Wronskian lemma}\label{chap:wronskian}

\begin{theorem}[Zero-set-safe Wronskian lemma]\label{thm:critical-wronskian}
Let $f,g\in H^{1/2}_{\loc}(\R)\cap C(\R)$ be locally bounded.  Assume
\begin{equation}
 |f|=|g|\quad\text{almost everywhere}               \label{eq:W-density}
\end{equation}
and
\begin{equation}
 \overline f f'-\overline g g'=0                   \label{eq:W-log}
\end{equation}
as a local $H^{-1/2}$--$H^{1/2}$ quadratic-form distribution.  Then
\begin{equation}
  g f'-f g'=0\quad\text{in }\D'(\R).               \label{eq:W-zero}
\end{equation}
\end{theorem}

\begin{proof}
\textbf{Step 1: localization.}
Take a real cutoff $\chi$.  The derivative terms created by $\chi$ cancel:
\begin{align*}
 &\overline{\chi f}(\chi f)'-\overline{\chi g}(\chi g)'\\
 &\qquad=\chi^2(\overline f f'-\overline g g')
          +\chi\chi'(|f|^2-|g|^2)=0,
\end{align*}
while
$W(\chi f,\chi g)=\chi^2W(f,g)$.  It is therefore enough to prove the
identity after a compactly supported cutoff.  We now assume
$f,g\in H^{1/2}\cap L^\infty$.  Because the chosen representatives are
continuous, the almost-everywhere identity $|f|=|g|$ holds in fact at every
point.

\textbf{Step 2: a legal test function.}
Put
\[
 \rho=|f|^2=|g|^2,
 \qquad s_\eps=\frac{\rho}{\rho+\eps},
 \qquad h_\eps=\frac{fg}{\rho+\eps}.
\]
We spell out why these are legal critical multipliers.  The function
$r\mapsto r/(r+\eps)$ is bounded, Lipschitz on the bounded range of
$\rho$, and vanishes at zero.  Hence
$s_\eps\in H^{1/2}\cap L^\infty$ by
\cref{lem:Lip-composition}, and the algebra lemma gives
$s_\eps f,s_\eps g\in H^{1/2}\cap L^\infty$.
For $h_\eps$, the reciprocal $r\mapsto(r+\eps)^{-1}$ does not vanish at
zero and therefore is not itself an $L^2$ function on the line.  Choose a
fixed cutoff $\kappa$ equal to one on the supports of $f$ and $g$.  The
function
\[
 r\longmapsto\frac1{r+\eps}-\frac1\eps
\]
is Lipschitz on the range of $\rho$ and vanishes at zero, so
\[
 \kappa(\rho+\eps)^{-1}
 =\frac{\kappa}{\eps}
  +\kappa\left((\rho+\eps)^{-1}-\eps^{-1}\right)
 \in H^{1/2}\cap L^\infty.
\]
Since $fg$ is supported where $\kappa=1$, the algebra estimate now gives
\[
 h_\eps=fg\,\kappa(\rho+\eps)^{-1}
 \in H^{1/2}\cap L^\infty.
\]
We may therefore test \eqref{eq:W-log} against $h_\eps\phi$, with
$\phi\in C_c^\infty$.  The equality of logarithmic forms extends from
smooth tests to this multiplier by
\cref{lem:quadratic-form-extension}.  Concretely, $h_\eps\phi$ is
continuous and compactly supported; mollify it after one fixed cutoff.  The
mollifications converge in both $H^{1/2}$ and $L^\infty$ by
\cref{lem:simultaneous-mollification}, and
\eqref{eq:quadratic-extension-bound} makes both quadratic pairings converge.

Since
\[
 \overline f h_\eps=s_\eps g,
 \qquad
 \overline g h_\eps=s_\eps f,
\]
we obtain
\begin{equation}
 \pair{f'}{s_\eps g\phi}-\pair{g'}{s_\eps f\phi}=0. \label{eq:s-eps-W}
\end{equation}

\textbf{Step 3: concentrate the error on the common zero set.}
Let $D=gf'-fg'$.  Subtract \eqref{eq:s-eps-W} from its defining pairing:
\begin{align}
 \pair D\phi
 &=\pair{f'}{(1-s_\eps)g\phi}
   -\pair{g'}{(1-s_\eps)f\phi}\\
 &=\pair{f'}{T_\eps(g)\phi}
   -\pair{g'}{T_\eps(f)\phi},                      \label{eq:W-Teps}
\end{align}
where the last identity uses $|f|=|g|$ and
$T_\eps(z)=\eps z/(|z|^2+\eps)$.

\textbf{Step 4: remove the regularizer.}
By \cref{lem:Teps}, both $T_\eps(f)$ and $T_\eps(g)$ converge strongly to
zero in local $H^{1/2}$.  Multiplication by $\phi$ preserves that convergence.
Since $f',g'\in H^{-1/2}_{\loc}$, the right side of \eqref{eq:W-Teps} tends
to zero.  Thus $\pair D\phi=0$ for every test function, which is
\eqref{eq:W-zero}.
\end{proof}

\begin{corollary}[Diagonal Wronskian for the NLS pair]\label{cor:diagonal-W}
There is a full-measure set $E\subset\R$ such that, for every $t\in E$,
\begin{equation}
 u(t)v_x(t)-v(t)u_x(t)=0
 \quad\text{in }\D'(\R_x).                         \label{eq:time-W}
\end{equation}
The set $E$ may be chosen so that \eqref{eq:good-slices} and
\eqref{eq:log-derivative-equality} also hold at every $t\in E$.
\end{corollary}

\begin{proof}
Use the good slices from \cref{chap:smoothing}, the sliced identity in
\cref{thm:current-recovery}, and apply \cref{thm:critical-wronskian} at each
such time.  Intersecting finitely many full-measure sets preserves full
measure.
\end{proof}

\begin{remark}[What the Wronskian does and does not prove]
On an interval on which $u(t)$ never vanishes, \eqref{eq:time-W} implies that
$v/u$ is constant.  It does not by itself show that constants on different
components of $\{u(t)\ne0\}$ agree.  The two-particle argument in the next
chapters is precisely the dynamical gluing mechanism.
\end{remark}

\begin{checkpoint}
The current identity has been converted into zero normal derivative for the
exterior product on the diagonal.  The argument is valid at common zeros and
uses no quotient.
\end{checkpoint}

\chapter{The two-particle exterior product}\label{chap:exterior}

\section{Definition and algebraic equation}
Define
\begin{equation}
 F(t,x,y)=u(t,x)v(t,y)-v(t,x)u(t,y).                \label{eq:F-def}
\end{equation}
By \eqref{eq:tensor-norm} and continuity of $u,v$ in $L^2$,
\begin{equation}
 F\in C(\R;L^2(\R^2)).                             \label{eq:F-C-L2}
\end{equation}
It is antisymmetric:
\begin{equation}
 F(t,y,x)=-F(t,x,y).                                \label{eq:F-antisym}
\end{equation}
Let
\[
  \rho(t,x)=|u(t,x)|^2=|v(t,x)|^2.
\]

\begin{proposition}[Exterior-product equation]\label{prop:F-equation}
In distributions on $\R_t\times\R_x^2$,
\begin{equation}
 iF_t+F_{xx}+F_{yy}
 =\sigma\bigl(\rho(t,x)+\rho(t,y)\bigr)F.          \label{eq:F-equation}
\end{equation}
On every bounded interval $I$,
\begin{align}
 F&\in L^4(I\times\R^2),                           \label{eq:F-L4}\\
 \bigl(\rho(x)+\rho(y)\bigr)F&\in L^2(I\times\R^2). \label{eq:QF-L2}
\end{align}
\end{proposition}

\begin{proof}
Write $N(z)=|z|^2z$.  For a smooth pair, with no equality-of-moduli
assumption imposed on the pair, direct differentiation gives the unfactored
identity
\begin{align}
 &(i\partial_t+\partial_x^2+\partial_y^2)
   \bigl(u(x)v(y)-v(x)u(y)\bigr) \notag\\
 &\quad=\sigma\bigl[
   N(u)(x)v(y)+u(x)N(v)(y)
  -N(v)(x)u(y)-v(x)N(u)(y)\bigr].                 \label{eq:F-unfactored}
\end{align}
This is the identity that is stable under arbitrary smooth approximation;
the approximating pair need not have equal moduli.

Choose smooth initial data converging to the two given data and denote the
corresponding solutions by $u_n,v_n$.  On every compact interval,
$u_n\to u$ and $v_n\to v$ in $C_tL^2_x$, while
$N(u_n)\to N(u)$ and $N(v_n)\to N(v)$ in $L^1_tL^2_x$.
Consequently every tensor on the right of \eqref{eq:F-unfactored} converges
in $L^1_tL^2_{x,y}$.  For example,
\begin{align*}
 &\norm{N(u_n)\otimes v_n-N(u)\otimes v}_{L^1_tL^2_{x,y}}\\
 &\quad\le
   \norm{N(u_n)-N(u)}_{L^1_tL^2_x}
   \norm{v_n}_{L^\infty_tL^2_y}
  +\norm{N(u)}_{L^1_tL^2_x}
   \norm{v_n-v}_{L^\infty_tL^2_y}\longrightarrow0.
\end{align*}
The left side converges in distributions by the $C_tL^2_{x,y}$ convergence
of the exterior products.  Thus \eqref{eq:F-unfactored} holds for the given
rough solutions.  Only now use the exact hypothesis
$|u|^2=|v|^2=\rho$: its four terms factor algebraically as
\[
 \rho(x)\bigl(u(x)v(y)-v(x)u(y)\bigr)
 +\rho(y)\bigl(u(x)v(y)-v(x)u(y)\bigr),
\]
which proves \eqref{eq:F-equation}.

For the $L^4$ bound, a typical tensor term satisfies
\begin{align*}
 \norm{u(t,x)v(t,y)}_{L^4_{t,x,y}}^4
 &=\int_I\norm{u(t)}_4^4\norm{v(t)}_4^4\,\dd t\\
 &\le\norm u_{L^8_tL^4_x(I)}^4
       \norm v_{L^8_tL^4_x(I)}^4.
\end{align*}
For the forcing, one typical term in the factored right side is
$|u(x)|^2u(x)v(y)$.  Its squared $L^2$ norm is
\[
 \int_I\norm{u(t)}_6^6\norm{v(t)}_2^2\,\dd t,
\]
which is finite by mass conservation and the $L^6_{t,x}$ bound.  The other
terms are identical, proving \eqref{eq:QF-L2}.
\end{proof}

\section{Normal and tangential coordinates}
Set
\begin{equation}
 r=\frac{x+y}{\sqrt2},\qquad s=\frac{x-y}{\sqrt2}. \label{eq:rs-coordinates}
\end{equation}
We use the same letter $F$ after this orthogonal change of variables.  Then
\begin{equation}
 iF_t+F_{rr}+F_{ss}=QF,                             \label{eq:F-rs}
\end{equation}
where
\begin{equation}
 Q(t,r,s)=\sigma\left[
  \rho\!\left(t,\frac{r+s}{\sqrt2}\right)
 +\rho\!\left(t,\frac{r-s}{\sqrt2}\right)
 \right].                                          \label{eq:Q-ridge}
\end{equation}
The function $Q$ is real and even in $s$, while $F$ is odd in $s$.

\begin{lemma}[Ridge-potential bounds]\label{lem:Q-bounds}
On every bounded time interval $I$,
\begin{align}
 Q&\in L^1(I;L^\infty(\R^2)),                      \label{eq:Q-L1Linfty}\\
 \norm Q_{L^1_tL^\infty_{r,s}(I)}
 &\le2|\sigma||I|^{1/2}\norm u_{L^4_tL^\infty_x(I)}^2. \label{eq:Q-L1-bound}
\end{align}
Furthermore, for every strip
$S_{a,h}=\{(r,s):a<s<a+h\}$,
\begin{equation}
 \norm{\1_{S_{a,h}}Q}_{L^2(I\times\R^2)}
 \le C|\sigma|h^{1/2}\norm u_{L^4_tL^4_x(I)}^2,   \label{eq:Q-strip}
\end{equation}
uniformly in $a$.
\end{lemma}

\begin{proof}
At each time,
$\norm Q_\infty\le2|\sigma|\norm{u(t)}_\infty^2$.
Integrate and use Cauchy--Schwarz in time to obtain
\eqref{eq:Q-L1-bound}.  For one ridge term, the change of variables in $r$
gives
\begin{align*}
 &\int_I\int_a^{a+h}\int_\R
 \rho\!\left(t,\frac{r+s}{\sqrt2}\right)^2
 \,\dd r\,\dd s\,\dd t\\
 &\qquad=\sqrt2h\int_I\norm{u(t)}_4^4\,\dd t.
\end{align*}
The second ridge is identical, and the triangle inequality proves
\eqref{eq:Q-strip}.
\end{proof}

\begin{remark}
The full spacetime $L^2$ norm of $Q$ is generally infinite because each
one-body density is repeated along a transverse direction.  The proof below
can use either the global $L^1_tL^\infty_{r,s}$ control on short time
intervals or the uniform strip estimate.  We prove the half-space theorem in
a form that accommodates both decompositions.
\end{remark}

\chapter{Distribution-valued diagonal traces}\label{chap:traces}

The equal-density condition gives the Dirichlet trace $F|_{s=0}=0$ formally,
and the Wronskian gives the Neumann trace.  At $L^2$ regularity neither trace
is initially classical.  We now construct them by testing in the tangential
variables and using the equation as an ordinary differential equation in the
normal variable.

\section{The one-dimensional trace theorem used below}

\begin{lemma}[Interior first-derivative estimate]
\label{lem:interior-first-derivative}
Let $J'\Subset J$ be bounded intervals.  There is $C=C(J',J)$ such that every
smooth $h$ on $J$ satisfies
\begin{equation}
 \norm{h'}_{L^2(J')}
 \le C\left(\norm h_{L^2(J)}+\norm{h''}_{L^2(J)}\right). \label{eq:interior-H1-estimate}
\end{equation}
\end{lemma}

\begin{proof}
Choose $\chi\in C_c^\infty(J)$ with $\chi=1$ on $J'$.  Integration by parts
gives
\[
 \int\chi^2|h'|^2
 =-\Rea\int\chi^2h''\overline h
   -2\Rea\int\chi\chi'h'\overline h.
\]
Cauchy--Schwarz and $2ab\le\tfrac12a^2+2b^2$ absorb half of the left side and
give
\[
 \norm{\chi h'}_2^2
 \le C_\chi\left(\norm h_2^2+\norm{h''}_2^2\right).
\]
Since $\chi=1$ on $J'$, this is \eqref{eq:interior-H1-estimate}.
\end{proof}

\begin{theorem}[Local $H^2$ characterization and $C^1$ representative]
\label{thm:local-H2-characterization}
Let $h\in\D'(J)$.  If $h$ and its distributional second derivative are
represented by functions in $L^2_{\loc}(J)$, then
$h\in H^2_{\loc}(J)$.  It has a unique $C^1$ representative, and on every
$J'\Subset J$
\begin{equation}
 \norm h_{H^2(J')}
 \le C_{J',J''}
 \left(\norm h_{L^2(J'')}+\norm{h''}_{L^2(J'')}\right) \label{eq:local-H2-bound}
\end{equation}
for any intermediate interval $J'\Subset J''\Subset J$.
\end{theorem}

\begin{proof}
Mollify $h$ inside $J''$.  The mollifications $h_\eps$ and
$h_\eps''$ converge in $L^2$ on every smaller interval and are uniformly
bounded there.  Apply \cref{lem:interior-first-derivative} on nested
intervals to obtain a uniform $L^2$ bound for $h_\eps'$.  Weak compactness
and the distributional identity identify its weak limit with $h'$, so
$h'\in L^2_{\loc}$.  Since $h''\in L^2_{\loc}$ by hypothesis, this proves
$h\in H^2_{\loc}$ and the displayed estimate.

For a one-dimensional $H^1$ function $g$, the absolutely continuous
representative satisfies
\[
 |g(x)-g(y)|\le|x-y|^{1/2}\norm{g'}_{L^2([x,y])}.
\]
Apply this first to $h'$ and then to $h$.  Thus $h'$ and $h$ are continuous,
so $h$ has a $C^1$ representative.  Two continuous representatives agreeing
almost everywhere agree everywhere.
\end{proof}

\section{Scalarization in the normal variable}
Fix a bounded interval $I$ and
$\psi\in C_c^\infty(I\times\R_r)$.  Define a scalar distribution in $s$ by
\begin{equation}
 H_\psi(s)=\pair{F(\cdot,\cdot,s)}{\psi}_{t,r}.     \label{eq:Hpsi-def}
\end{equation}

\begin{lemma}[$H^2$ regularity of scalarized traces]\label{lem:Hpsi-H2}
For every $\psi$ as above,
\[
 H_\psi\in H^2_{\loc}(\R_s).
\]
In particular $H_\psi(0)$ and $H_\psi'(0)$ are well-defined.
\end{lemma}

\begin{proof}
Since $F\in L^2(I\times\R^2)$, Cauchy--Schwarz in $(t,r)$ gives
$H_\psi\in L^2_s$ locally.  Pair \eqref{eq:F-rs} with $\psi$ and integrate by
parts in $t$ and $r$:
\begin{equation}
 H_\psi''(s)
 =\pair{QF}{\psi}_{t,r}
  +\pair F{i\psi_t-\psi_{rr}}_{t,r}.               \label{eq:Hpsi-second}
\end{equation}
Both right-hand terms belong to $L^2_s$ locally because
$QF,F\in L^2(I\times\R^2)$ and the test functions are compactly supported.
Thus $H_\psi''\in L^2_{\loc}$ in the distributional sense.  The local characterization in
\cref{thm:local-H2-characterization} then gives
$H_\psi\in H^2_{\loc}$ together with its $C^1$ representative.
\end{proof}

\section{The Dirichlet trace}

\begin{lemma}[Zero Dirichlet trace]\label{lem:Dirichlet-trace}
For every $\psi\in C_c^\infty(I\times\R)$,
\[
 H_\psi(0)=0.
\]
Equivalently, $F|_{s=0}=0$ as a distribution in $(t,r)$.
\end{lemma}

\begin{proof}
Antisymmetry \eqref{eq:F-antisym} becomes $F(t,r,-s)=-F(t,r,s)$ in the new
coordinates.  Hence $H_\psi$ is odd.  Its continuous representative, supplied
by $H^2_{\loc}\hookrightarrow C^1_{\loc}$ in one dimension, vanishes at zero.
\end{proof}

\section{The Neumann trace}
The derivative requires the half-derivative smoothing and the translation
difference quotient from \cref{lem:translation-dq}.

\begin{proposition}[Zero Neumann trace]\label{prop:Neumann-trace}
For every $\psi\in C_c^\infty(I\times\R)$,
\[
 H_\psi'(0)=0.
\]
Equivalently, $F_s|_{s=0}=0$ as a distribution in $(t,r)$.
\end{proposition}

\begin{proof}
Put $z=(r-s)/\sqrt2$.  Then $r=\sqrt2z+s$, $\dd r=\sqrt2\,\dd z$, and
\eqref{eq:F-def} gives
\begin{align}
 H_\psi(s)=\sqrt2\iint
 &\big[u(t,z+\sqrt2s)v(t,z)-v(t,z+\sqrt2s)u(t,z)\big]\notag\\
 &\hspace{25mm}\times\psi(t,\sqrt2z+s)\,\dd z\,\dd t. \label{eq:Hpsi-expanded}
\end{align}
The bracket vanishes at $s=0$.

Fix $s_0>0$.  The set of all $z$ for which
$\psi(t,\sqrt2z+s)$ can be nonzero for some $t$ and $|s|\le s_0$ is compact.
Choose $\chi_0,\chi_1\in C_c^\infty(\R)$ so that $\chi_0=1$ on that set and
$\chi_1=1$ on an open neighbourhood of
\[
 \supp\chi_0+[-\sqrt2s_0,\sqrt2s_0].
\]
Thus whenever $|s|\le s_0$ and $\psi(t,\sqrt2z+s)\ne0$, both
$\chi_1(z)$ and $\chi_1(z+\sqrt2s)$ are equal to one.  Local smoothing gives
\[
 \chi_1u,\chi_1v\in L^2(I;H^{1/2}(\R)).
\]

For $s\ne0$ define the genuinely localized difference quotients
\[
 \widetilde D_su
   =\frac{(\chi_1u)(\,\cdot+\sqrt2s)-\chi_1u}{s},
 \qquad
 \widetilde D_sv
   =\frac{(\chi_1v)(\,\cdot+\sqrt2s)-\chi_1v}{s}.
\]
By \cref{lem:translation-dq}, for almost every $t$,
\[
 \widetilde D_su(t)\longrightarrow\sqrt2(\chi_1u(t))_x,
 \qquad
 \widetilde D_sv(t)\longrightarrow\sqrt2(\chi_1v(t))_x
 \quad\text{in }H^{-1/2}.
\]
The same lemma gives the pointwise bounds
\[
 \norm{\widetilde D_su(t)}_{H^{-1/2}}
 \le C\norm{\chi_1u(t)}_{H^{1/2}},
 \qquad
 \norm{\widetilde D_sv(t)}_{H^{-1/2}}
 \le C\norm{\chi_1v(t)}_{H^{1/2}},
\]
uniformly for $0<|s|\le s_0$.  Squaring and applying dominated convergence
in $t$ therefore yields
\begin{align}
 \widetilde D_su&\longrightarrow\sqrt2(\chi_1u)_x,
 &\widetilde D_sv&\longrightarrow\sqrt2(\chi_1v)_x
 &&\text{in }L^2(I;H^{-1/2}).                 \label{eq:DQ-time-convergence}
\end{align}
This formulation avoids a cutoff commutator: on the support of every test
function used below, the localized quotients agree exactly with the
unlocalized quotients.

Write
\[
 \psi_s(t,z)=\psi(t,\sqrt2z+s),
 \qquad \psi_0(t,z)=\psi(t,\sqrt2z).
\]
Multiplication by $\psi_s(t,\cdot)$ maps $H^{1/2}$ to itself.  The proof of
\cref{lem:smooth-multiplier}, combined with the mean-value theorem for the
first two spatial derivatives of $\psi$, gives the uniform operator estimate
\begin{equation}
 \norm{(\psi_s-\psi_0)f}_{H^{1/2}}
 \le C_\psi|s|\norm f_{H^{1/2}}
 \qquad(0<|s|\le s_0).                         \label{eq:translated-test-multiplier}
\end{equation}
Because $\chi_0=1$ on the spatial support of every $\psi_s$,
local smoothing and the multiplier lemma imply
\[
 \chi_0u,\chi_0v\in L^2(I;H^{1/2}),
 \qquad
 v\psi_s,\ u\psi_s\in L^2(I;H^{1/2}).
\]

Dividing \eqref{eq:Hpsi-expanded} by $s$, using the support property of
$\chi_1$, and interpreting the products by the
$H^{-1/2}$--$H^{1/2}$ pairing gives the exact identity
\begin{align}
 \frac{H_\psi(s)}s
 =\sqrt2\int_I\Big(
   &\pair{\widetilde D_su(t)}{v(t)\psi_s(t)}
    -\pair{\widetilde D_sv(t)}{u(t)\psi_s(t)}\Big)\,\dd t.       \label{eq:Hpsi-dq-exact}
\end{align}
Replace $\psi_s$ by $\psi_0$.  By Cauchy--Schwarz in time,
\eqref{eq:translated-test-multiplier}, and the uniform bounds following
\eqref{eq:DQ-time-convergence}, the absolute value of the resulting error is
at most
\[
 C|s|\Big(
   \norm{\chi_1u}_{L^2_tH^{1/2}}
   \norm{\chi_0v}_{L^2_tH^{1/2}}
  +\norm{\chi_1v}_{L^2_tH^{1/2}}
   \norm{\chi_0u}_{L^2_tH^{1/2}}
 \Big),
\]
which tends to zero.

We may now pass to the limit in the two pairings by
\eqref{eq:DQ-time-convergence}.  Since $\chi_1$ is identically one on a
neighbourhood of the support of $\psi_0$, distributional differentiation
shows
\[
 \pair{(\chi_1u)_x}{v\psi_0}=\pair{u_x}{v\psi_0},
 \qquad
 \pair{(\chi_1v)_x}{u\psi_0}=\pair{v_x}{u\psi_0}.
\]
Consequently
\begin{equation}
 H_\psi'(0)
 =2\int_I\pair{vu_x-uv_x}{\psi(t,\sqrt2\,\cdot)}\,\dd t. \label{eq:Hpsi-prime-W}
\end{equation}
The integrand belongs to $L^1(I)$ by Cauchy--Schwarz, local smoothing, and
the smooth-multiplier estimate.  It vanishes for almost every $t$ by
\cref{cor:diagonal-W}.  Hence $H_\psi'(0)=0$.
\end{proof}

\section{The zero-extension jump formula}
Let
\begin{equation}
  G(t,r,s)=\1_{\{s>0\}}F(t,r,s).                   \label{eq:G-def}
\end{equation}
Multiplication by a half-space is bounded on $L^2$, hence
$G\in C(I;L^2_{r,s})\cap L^4(I\times\R^2)$.

\begin{lemma}[One-dimensional jump identity]\label{lem:jump-formula}
Let $h\in H^2_{\loc}(\R)$ and write $h_+=\1_{\{s>0\}}h$.  Then
\begin{equation}
 h_+''=\1_{\{s>0\}}h''+h'(0)\delta_0+h(0)\delta_0'
 \quad\text{in }\D'(\R).                           \label{eq:jump-formula}
\end{equation}
\end{lemma}

\begin{proof}
For a test function $\varphi$, integrate $\int_0^\infty h\varphi''$ twice by
parts.  The boundary terms are $h'(0)\varphi(0)-h(0)\varphi'(0)$, which are
exactly the pairings with $h'(0)\delta_0+h(0)\delta_0'$.
\end{proof}

\begin{proposition}[The zero extension solves the same equation]
\label{prop:G-equation}
On $I\times\R^2$,
\begin{equation}
 iG_t+G_{rr}+G_{ss}=QG                            \label{eq:G-equation}
\end{equation}
in distributions, and $G=0$ on the open half-space $s<0$.
\end{proposition}

\begin{proof}
Test first in $(t,r)$ and use the scalar functions $H_\psi$ from
\eqref{eq:Hpsi-def}.  By \cref{lem:jump-formula}, the only possible defect in
$\partial_s^2G$ is
$H_\psi'(0)\delta_0+H_\psi(0)\delta_0'$.  Both coefficients vanish by
\cref{lem:Dirichlet-trace,prop:Neumann-trace}.  The remaining terms are
simply the restriction of \eqref{eq:F-rs} to $s>0$.  This proves the identity
first against tests of the form $\psi(t,r)\varphi(s)$.  Finite sums of such tests approximate an
arbitrary compactly supported smooth test together with two derivatives by
\cref{lem:tensor-test-density}; continuity of all $L^2$ pairings completes
the passage to general tests.  Vanishing on $s<0$ is the definition.
\end{proof}

\begin{checkpoint}
The nonlinear phase-retrieval problem has now been reduced to a linear
half-space uniqueness statement in two spatial dimensions.  All boundary
traces and jump terms have been handled in the exact low-regularity solution
class.
\end{checkpoint}

\chapter{A self-contained half-space Carleman theorem}\label{chap:carleman}

We prove the linear uniqueness statement needed for $G$.  The proof has three
parts.  First, a uniform exponential-weight estimate is obtained by writing an
explicit inverse for the conjugated free Schr\"odinger operator.  Second, a
small-potential absorption lemma converts that estimate into exponential
separation.  Third, a moving spatial interface makes the time-cutoff errors
occur strictly inside the already known zero side.  Quantitative iteration
propagates a half-space zero through any prescribed finite distance.

\section{The solution and source spaces}
On a time interval $J$ and in two spatial dimensions, define
\begin{align}
 X(J)&=L^1(J;L^2(\R^2))+L^{4/3}(J\times\R^2),       \label{eq:X-space}\\
 X'(J)&=L^\infty(J;L^2(\R^2))\cap L^4(J\times\R^2),\label{eq:Xprime-space}\\
 Y(J)&=L^1(J;L^\infty(\R^2))+L^2(J\times\R^2).     \label{eq:Y-space}
\end{align}
Every component is the scalar mixed-norm space of
\cref{def:scalar-mixed-norm}.  The sum space has its infimum norm and the
intersection has the maximum norm; completeness and restriction are supplied
by \cref{lem:sum-intersection-complete}.  When the time interval is the whole
line, it is omitted.

\begin{lemma}[Critical multiplier estimate]\label{lem:XY-multiplier}
There is an absolute $C$ such that
\begin{equation}
  \norm{WZ}_{X(J)}\le C\norm W_{Y(J)}\norm Z_{X'(J)}. \label{eq:XY-multiplier}
\end{equation}
\end{lemma}

\begin{proof}
Write $W=W_1+W_2$ with
$W_1\in L^1_tL^\infty_z$ and $W_2\in L^2_{t,z}$.  Then
\[
 \norm{W_1Z}_{L^1_tL^2_z}
 \le\norm{W_1}_{L^1_tL^\infty_z}\norm Z_{L^\infty_tL^2_z},
\]
while
\[
 \norm{W_2Z}_{L^{4/3}_{t,z}}
 \le\norm{W_2}_{L^2_{t,z}}\norm Z_{L^4_{t,z}}.
\]
Take the infimum over decompositions of $W$.
\end{proof}

\section{The conjugated free operator}
Put
\[
  P=i\partial_t+\Delta_z,
  \qquad z=(z_1,z_2)\in\R^2.
\]
For $\beta>0$,
\begin{equation}
 P_\beta=\e^{\beta z_1}P\e^{-\beta z_1}
 =i\partial_t+\Delta_z-2\beta\partial_{z_1}+\beta^2. \label{eq:Pbeta}
\end{equation}
We construct a right inverse $T_\beta$ whose bounds are independent of
$\beta$.

Take the spatial Fourier transform.  If $P_\beta w=f$, then
\begin{equation}
 i\partial_t\widehat w
 +(\beta^2-|\xi|^2-2i\beta\xi_1)\widehat w
 =\widehat f.                                      \label{eq:Pbeta-ODE}
\end{equation}
Let $a(\xi)=\beta^2-|\xi|^2$.  For $\xi_1<0$ solve forward from $-\infty$;
for $\xi_1>0$ solve backward from $+\infty$:
\begin{equation}
 \widehat{T_\beta f}(t,\xi)=
 \begin{cases}
 -i\displaystyle\int_{-\infty}^t
   \e^{(ia(\xi)+2\beta\xi_1)(t-s)}\widehat f(s,\xi)\,\dd s,
   &\xi_1<0,\\[1.2em]
 i\displaystyle\int_t^{\infty}
   \e^{(ia(\xi)+2\beta\xi_1)(t-s)}\widehat f(s,\xi)\,\dd s,
   &\xi_1>0.
 \end{cases}                                      \label{eq:Tbeta-def}
\end{equation}
The exponential has modulus at most one in the selected integration
direction.

\begin{lemma}[Kernel estimates for $T_\beta$]\label{lem:Tbeta-kernel}
Let $K_\beta(\tau)$ be the spatial convolution operator appearing in
\eqref{eq:Tbeta-def}.  For every $\tau\ne0$,
\begin{align}
 \norm{K_\beta(\tau)}_{L^2_z\to L^2_z}&\le1,        \label{eq:K-L2}\\
 \norm{K_\beta(\tau)}_{L^1_z\to L^\infty_z}
 &\le C|\tau|^{-1},                                 \label{eq:K-disp}\\
 \norm{K_\beta(\tau)}_{L^{4/3}_z\to L^4_z}
 &\le C|\tau|^{-1/2}.                              \label{eq:K-43-4}
\end{align}
The constants are independent of $\beta>0$.
\end{lemma}

\begin{proof}
The Fourier multiplier of $K_\beta(\tau)$ is supported where
$\tau\xi_1<0$ and has modulus
$\e^{-2\beta|\tau\xi_1|}$, proving \eqref{eq:K-L2}.

For the kernel bound, factor the inverse Fourier integral into the $\xi_2$
and $\xi_1$ variables.  The full Gaussian integral in $\xi_2$ has modulus
$C|\tau|^{-1/2}$ by \eqref{eq:complex-gaussian}, interpreted as the
boundary value from positive Gaussian damping.  After reflecting the selected half-line if necessary, the
remaining factor is
\begin{equation}
 \int_0^\infty
 \e^{i x_1\eta-i\tau\eta^2}
 \e^{-2\beta|\tau|\eta}\,\dd\eta.                 \label{eq:half-Fresnel}
\end{equation}
Scale $\eta=|\tau|^{-1/2}y$.  It remains to show that
\[
 \left|\int_0^\infty
   \e^{i(ay\pm y^2)}\e^{-by}\,\dd y\right|\le C
 \qquad(a\in\R,\ b\ge0).                          \label{eq:uniform-Fresnel}
\]
We prove the needed second-derivative estimate directly.  Suppose
$\phi\in C^2([0,R])$ and $|\phi''|\ge\lambda>0$.  The derivative $\phi'$ is
monotone.  The set $E=\{|\phi'|\le\sqrt\lambda\}$ is therefore an interval
of length at most $2\lambda^{-1/2}$, and
\[
 \left|\int_E\e^{i\phi}A\right|
 \le2\lambda^{-1/2}\norm A_\infty.
\]
The complement has at most two interval components.  On either one,
integration by parts with
$\e^{i\phi}=(i\phi')^{-1}(\e^{i\phi})'$ gives boundary terms bounded by
$2\lambda^{-1/2}\norm A_\infty$, an $A'$ term bounded by
$\lambda^{-1/2}\int|A'|$, and a final term bounded by
\[
 \norm A_\infty\int\frac{|\phi''|}{|\phi'|^2}
 =\norm A_\infty\operatorname{Var}\!\left(\frac1{\phi'}\right)
 \le2\lambda^{-1/2}\norm A_\infty.
\]
Summing the components proves
\[
 \left|\int_0^R\e^{i\phi(y)}A(y)\,\dd y\right|
 \le C\lambda^{-1/2}
 \left(\norm A_\infty+\int_0^R|A'(y)|\,\dd y\right).
\]
For the phase in \eqref{eq:uniform-Fresnel}, $\lambda=2$.  If
$b>0$, take $A(y)=\e^{-by}$; then
$\norm A_\infty+\int_0^\infty|A'|\le2$, and absolute convergence allows
$R\to\infty$.  When $b=0$, the preceding estimate gives uniform boundedness
of the partial integrals, but we also record convergence of the improper
integral.  Choose $R_0>2(1+|a|)$.  On $[R_0,\infty)$ the derivative of
$ay\pm y^2$ has fixed sign and magnitude comparable to $y$.  One integration
by parts gives, for $R_2>R_1\ge R_0$,
\[
 \left|\int_{R_1}^{R_2}\e^{i(ay\pm y^2)}\,\dd y\right|
 \le C\left(R_1^{-1}+\int_{R_1}^\infty y^{-2}\,\dd y\right)
 \le CR_1^{-1}.
\]
Thus the partial integrals are Cauchy and their limit obeys the same uniform
bound.  This proves \eqref{eq:uniform-Fresnel} for every $b\ge0$.  Hence the
$\xi_1$ factor is $O(|\tau|^{-1/2})$.

For completeness, the oscillatory kernel is obtained without treating a
conditionally convergent Fourier integral as an ordinary Lebesgue integral.
Insert the multiplier cutoff $\e^{-\delta|\xi|^2}$, $\delta>0$.  Its
inverse transform is a continuous convolution kernel.  The full Gaussian
factor has modulus at most $C|\tau|^{-1/2}$ by
\eqref{eq:complex-gaussian}, while the half-line argument above applies to
the amplitude $A(y)=\e^{-by-cy^2}$: it is decreasing and
$\norm A_\infty+\int|A'|\le2$.  Thus the regularized kernels obey the
same $C|\tau|^{-1}$ bound, uniformly in $\delta$.  For Schwartz input the outputs of the
regularized multipliers converge in $L^2$, and the kernels converge in
tempered distributions to the inverse Fourier transform of the multiplier
of $K_\beta(\tau)$; the improper Fresnel calculation gives
the pointwise kernel limit.  Passing to the limit first on Schwartz
functions and then by $L^1$ density yields the bounded convolution
realization and proves \eqref{eq:K-disp}.

The $L^1\to L^\infty$ kernel realization and the $L^2$ multiplier
realization agree on Schwartz functions, so the compatible-operator
Riesz--Thorin theorem applies and gives \eqref{eq:K-43-4}.
\end{proof}

\begin{lemma}[Homogeneous bounds for the kernel family]
\label{lem:K-homogeneous}
For every $s\in\R$, $g\in L^2(\R^2)$, and $\beta\ge0$,
\begin{align}
 \norm{K_\beta(t-s)g}_{L^4_{t,z}(\R^{1+2})}
 &\le C\norm g_2,                                  \label{eq:K-homogeneous}\\
 \norm{K_\beta(t-s)^*g}_{L^4_{t,z}(\R^{1+2})}
 &\le C\norm g_2.                                  \label{eq:K-adjoint-homogeneous}
\end{align}
The constants are independent of $s$ and $\beta$.
\end{lemma}

\begin{proof}
Time translation reduces the first estimate to $s=0$.  Define
$A_\beta g(t)=K_\beta(t)g$.  The multiplier of
$A_\beta A_\beta^*$ from time $s$ to time $t$ is
$m_t(\xi)\overline{m_s(\xi)}$, where $m_\tau$ is the multiplier of
$K_\beta(\tau)$.  It vanishes when $t$ and $s$ have opposite signs.  When
they have the same sign, it is supported on one half-line in $\xi_1$, has
oscillatory factor
\[
 \exp\big(i(\beta^2-|\xi|^2)(t-s)\big),
\]
and damping factor
\[
 \exp\big(-2\beta(|t|+|s|)|\xi_1|\big).
\]
Consequently its $L^2_z$ operator norm is at most one.  The factorization and
half-Fresnel argument in \cref{lem:Tbeta-kernel}, now with
$|t-s|$ as the oscillatory scale and
$2\beta(|t|+|s|)$ as a nonnegative damping parameter, gives
\[
 \norm{K_\beta(t)K_\beta(s)^*}_{L^1_z\to L^\infty_z}
 \le C|t-s|^{-1}
\]
for $t\ne s$.  Interpolation therefore gives
\[
 \norm{K_\beta(t)K_\beta(s)^*}_{L^{4/3}_z\to L^4_z}
 \le C|t-s|^{-1/2}.
\]
The time-line HLS estimate yields
\[
 \norm{A_\beta A_\beta^*F}_{L^4_{t,z}}
 \le C\norm F_{L^{4/3}_{t,z}}.
\]
Hence, by duality,
\begin{align*}
 \norm{A_\beta^*F}_2^2
 &=\big|\pair{A_\beta A_\beta^*F}{F}\big|\\
 &\le \norm{A_\beta A_\beta^*F}_4\norm F_{4/3}
 \le C\norm F_{4/3}^2.
\end{align*}
Thus $A_\beta^*:L^{4/3}\to L^2$ and, by duality again,
$A_\beta:L^2\to L^4$, proving \eqref{eq:K-homogeneous}.  Replacing every
kernel by its adjoint gives the same composition estimates and proves
\eqref{eq:K-adjoint-homogeneous}.
\end{proof}

\begin{proposition}[Uniform linear-weight Carleman estimate]
\label{prop:linear-carleman}
For every $\beta\in\R$ and every $\varphi\in C_c^\infty(\R^{1+2})$,
\begin{equation}
 \norm{\e^{\beta z_1}\varphi}_{X'}
 \le C\norm{\e^{\beta z_1}P\varphi}_{X},           \label{eq:linear-carleman}
\end{equation}
where $C$ is absolute and independent of $\beta$.
\end{proposition}

\begin{proof}
By reflection it is enough to take $\beta>0$; the case $\beta=0$ is obtained
from the same formulas.  Set
$w=\e^{\beta z_1}\varphi$ and
$f=\e^{\beta z_1}P\varphi=P_\beta w$.  Because $w$ is compactly supported in
time, solving the scalar ODE \eqref{eq:Pbeta-ODE} in the selected direction
shows that $w=T_\beta f$.

We verify all four mappings needed for the sum-to-intersection estimate.  If
$f_1\in L^1_tL^2_z$, then \eqref{eq:K-L2} and Minkowski give
\[
 \norm{T_\beta f_1}_{L^\infty_tL^2_z}
 \le\norm{f_1}_{L^1_tL^2_z},
\]
while \cref{lem:K-homogeneous} and Minkowski give
\[
 \norm{T_\beta f_1}_{L^4_{t,z}}
 \le\int_\R
   \norm{K_\beta(t-s)f_1(s)}_{L^4_{t,z}}\,\dd s
 \le C\norm{f_1}_{L^1_tL^2_z}.
\]
If $f_2\in L^{4/3}_{t,z}$, then \eqref{eq:K-43-4} and
\cref{lem:HLS-time} give
\[
 \norm{T_\beta f_2}_{L^4_{t,z}}
 \le C\norm{f_2}_{L^{4/3}_{t,z}}.
\]
For the missing energy bound, fix $t$ and $g\in L^2_z$.  By spacetime
H\"older and the adjoint estimate \eqref{eq:K-adjoint-homogeneous},
\begin{align*}
 \big|\pair{T_\beta f_2(t)}{g}\big|
 &\le\int_\R
   \big|\pair{f_2(s)}{K_\beta(t-s)^*g}\big|\,\dd s\\
 &\le\norm{f_2}_{L^{4/3}_{s,z}}
       \norm{K_\beta(t-s)^*g}_{L^4_{s,z}}
 \le C\norm{f_2}_{4/3}\norm g_2.
\end{align*}
Taking the supremum over unit $g$ gives
$\norm{T_\beta f_2}_{L^\infty_tL^2_z}\le C\norm{f_2}_{4/3}$.

For a decomposition $f=f_1+f_2$, both components of the $X'$ norm are
therefore bounded by
$C(\norm{f_1}_{L^1L^2}+\norm{f_2}_{4/3})$.  Taking the infimum over all
decompositions proves \eqref{eq:linear-carleman}.
\end{proof}

\section{Extension to the graph class}
The exterior product is neither smooth nor compactly supported in space.  We
therefore prove the closure step at the level of the explicit inverse.  Every
approximation below is spatial only; compact time support is preserved.

\begin{lemma}[Spatial mollification in the source spaces]
\label{lem:spatial-mollification-source}
Let $\rho_\eps$ be a standard approximate identity on $\R^2$ and let
$J_0\Subset\R$ be bounded.
\begin{enumerate}
\item If $f_1\in L^1_tL^2_z$ is supported in $J_0$ in time, then
$\rho_\eps*_zf_1\to f_1$ in $L^1_tL^2_z$.
\item If $f_2\in L^{4/3}_{t,z}$ is supported in $J_0$ in time, then
$\rho_\eps*_zf_2\to f_2$ in $L^{4/3}_{t,z}$ and, for fixed $\eps>0$,
$\rho_\eps*_zf_2\in L^1_tL^2_z$.
\end{enumerate}
Consequently spatial mollification converges strongly in $X$.
\end{lemma}

\begin{proof}
The first two convergence statements are the standard approximate-identity
theorem, applied for almost every time and then integrated; equivalently,
approximate first by compactly supported simple functions.  For the additional
mapping, Young's inequality in $z$ gives
\[
 \norm{\rho_\eps*f_2(t)}_2
 \le\norm{\rho_\eps}_{4/3}\norm{f_2(t)}_{4/3}.
\]
Since $J_0$ has finite length,
$L^{4/3}(J_0)\hookrightarrow L^1(J_0)$, proving the claim.
\end{proof}

\begin{lemma}[Almost-everywhere parameter slicing of a first-order equation]
\label{lem:parameterwise-distribution-ODE}
Let $A:\R^d\to\C$ be measurable and bounded on compact sets.  Suppose
$y\in L^2_{\loc}(\R_t\times\R^d_\xi)$ and
$g\in L^1_{\loc}(\R_t;L^2_{\loc}(\R^d_\xi))$ satisfy
\begin{equation}
 \partial_ty-A(\xi)y=g
 \quad\text{in }\D'(\R_t\times\R^d_\xi).          \label{eq:parameter-distribution-equation}
\end{equation}
Then there is a single null set $N\subset\R^d$ such that, for every
$\xi\notin N$,
\[
 y(\cdot,\xi)\in L^2_{\loc}(\R),\qquad
 g(\cdot,\xi)\in L^1_{\loc}(\R),
\]
and
\begin{equation}
 \partial_ty(\cdot,\xi)-A(\xi)y(\cdot,\xi)
 =g(\cdot,\xi)
 \quad\text{in }\D'(\R_t).                       \label{eq:parameterwise-ODE}
\end{equation}
\end{lemma}

\begin{proof}
Fubini first removes one null set so that the two asserted section
integrability properties hold.  For a time test $\eta\in C_c^\infty(\R)$,
define, for almost every $\xi$,
\[
 R_\eta(\xi)
 =-\int_\R y(t,\xi)\eta'(t)\,\dd t
  -A(\xi)\int_\R y(t,\xi)\eta(t)\,\dd t
  -\int_\R g(t,\xi)\eta(t)\,\dd t.
\]
On each compact frequency set this is an $L^1$ function of $\xi$; the first
two terms follow from Cauchy--Schwarz on the compact time support, and the
last from Minkowski.  Testing \eqref{eq:parameter-distribution-equation}
against $\eta(t)\phi(\xi)$ gives
\[
 \int_{\R^d}R_\eta(\xi)\phi(\xi)\,\dd\xi=0
 \qquad(\phi\in C_c^\infty(\R^d)).
\]
The scalar fundamental lemma therefore gives $R_\eta=0$ almost everywhere.

To avoid taking an uncountable union of exceptional sets, choose for every
integer $m\ge1$ a countable family dense in
$C_c^\infty((-m,m))$ for the $C^1$ norm; rational piecewise polynomials
followed by convolution with rational smooth kernels provide such a family.
Remove the union of the null sets associated with all members of all these
families.  For a remaining frequency $\xi$, the map
\[
 \eta\longmapsto R_\eta(\xi)
\]
is continuous in the $C^1$ norm on each fixed compact time interval, with
bound
\[
 |R_\eta(\xi)|
 \le \norm{y(\cdot,\xi)}_{L^1(K)}\norm{\eta'}_\infty
   +|A(\xi)|\norm{y(\cdot,\xi)}_{L^1(K)}\norm\eta_\infty
   +\norm{g(\cdot,\xi)}_{L^1(K)}\norm\eta_\infty.
\]
Density extends $R_\eta(\xi)=0$ to every time test $\eta$, which is exactly
\eqref{eq:parameterwise-ODE}.
\end{proof}

\begin{lemma}[Scalar compact-support ODE]
\label{lem:compact-support-ODE}
Let $\lambda\in\C$, $y\in L^2(\R)$ have compact support, and
$g\in L^1(\R)$.  If
\[
 y'-\lambda y=g
\]
in distributions, then $y$ has a locally absolutely continuous
representative.  If $\Re\lambda<0$, then
\[
 y(t)=\int_{-\infty}^t\e^{\lambda(t-s)}g(s)\,\dd s,
\]
while if $\Re\lambda>0$, then
\[
 y(t)=-\int_t^\infty\e^{\lambda(t-s)}g(s)\,\dd s.
\]
\end{lemma}

\begin{proof}
On every bounded interval, $\lambda y+g\in L^1$, so the distributional
identity implies $y\in W^{1,1}_{\loc}$ and hence gives an absolutely
continuous representative.  Multiplying by $\e^{-\lambda t}$ and applying
the fundamental theorem of calculus gives
\[
 \e^{-\lambda t}y(t)-\e^{-\lambda a}y(a)
 =\int_a^t\e^{-\lambda s}g(s)\,\dd s.
\]
If $\Re\lambda<0$, take $a$ to the left of the compact support of $y$; then
$y(a)=0$ and the first formula follows.  If $\Re\lambda>0$, integrate from
$t$ to a point to the right of the support.  The signs are exactly those
displayed.
\end{proof}

\begin{lemma}[Fourier-section identification]
\label{lem:fourier-section-identification}
Let $w\in L^\infty_tL^2_z$ have compact support in a bounded time interval,
and suppose
$P_\beta w=f$ in distributions, where
$f\in L^1_tL^2_z$.  Then
\begin{equation}
 w=T_\beta f\quad\text{in distributions and in }X'. \label{eq:section-Tbeta}
\end{equation}
\end{lemma}

\begin{proof}
Because the time support is bounded,
$w\in L^2_{t,z}$.  Spatial Plancherel and Fubini give a full-measure set of
frequencies on which $\widehat w(\cdot,\xi)\in L^2_t$.  Moreover,
Minkowski gives
\[
 \left\|\int_\R|\widehat f(t,\xi)|\,\dd t\right\|_{L^2_\xi}
 \le\int_\R\norm{\widehat f(t)}_{L^2_\xi}\,\dd t
 =\norm f_{L^1_tL^2_z},
\]
so $\widehat f(\cdot,\xi)\in L^1_t$ for almost every $\xi$.
The spatial Fourier transform of the distributional equation is
\[
 \partial_t\widehat w
 -(ia(\xi)+2\beta\xi_1)\widehat w
 =-i\widehat f
 \quad\text{in }\D'(\R_t\times\R^2_\xi).
\]
Apply \cref{lem:parameterwise-distribution-ODE} with
$A(\xi)=ia(\xi)+2\beta\xi_1$.  It supplies one full-measure set of
frequencies on which this is an ordinary distributional equation in $t$.
For $\xi_1<0$ the real part of the coefficient is negative, and for
$\xi_1>0$ it is positive.  Apply \cref{lem:compact-support-ODE}.  The
resulting formulas are exactly \eqref{eq:Tbeta-def}.  The hyperplane
$\xi_1=0$ has Lebesgue measure zero and does not affect the inverse Fourier
transform.  The formulas and Plancherel identify $w=T_\beta f$ in
$L^2_{\loc}$ and hence in distributions.  Since the mapping estimates of
\cref{prop:linear-carleman} put $T_\beta f$ in $X'$, the distributional
equality permits us to use $T_\beta f$ as the scalar representative of
$w$.  Consequently $w\in X'$ and \eqref{eq:section-Tbeta} is an equality
of $X'$ equivalence classes.
\end{proof}

\begin{lemma}[Graph closure of the Carleman estimate]\label{lem:carleman-closure}
Suppose $U\in X'$, $U$ has compact support in time, and
$U(t,z)=0$ for $z_1>H$ for some $H$.  Fix $\beta>0$ and assume that the
weighted distribution $\e^{\beta z_1}PU$ belongs to $X$.  Then
\begin{equation}
 \norm{\e^{\beta z_1}U}_{X'}
 \le C\norm{\e^{\beta z_1}PU}_{X}.                \label{eq:graph-carleman}
\end{equation}
\end{lemma}

\begin{proof}
Set
\[
 w=\e^{\beta z_1}U,
 \qquad f=\e^{\beta z_1}PU.
\]
The upper support bound makes the weight bounded on the support of $U$, so
$w\in X'$.  It has the same compact time support as $U$, and distributional
conjugation gives
\[
 P_\beta w=f.                                      \label{eq:weighted-graph-equation}
\]
Choose a decomposition $f=f_1+f_2$ with
$f_1\in L^1_tL^2_z$ and $f_2\in L^{4/3}_{t,z}$.  Because $f$ has compact
time support as a distribution, choose a bounded interval $J_0$ containing
that support and replace the decomposition by
$\1_{J_0}f_1+\1_{J_0}f_2$.  Outside $J_0$ the equality
$f_1+f_2=f=0$ holds almost everywhere, so the new components still sum to
$f$; multiplication by an indicator does not increase either component norm.
Thus both components may be assumed supported in one bounded interval without
changing the infimum norm.

Convolve \eqref{eq:weighted-graph-equation} in $z$ with $\rho_\eps$:
\[
 P_\beta w_\eps=f_\eps,
 \qquad
 w_\eps=\rho_\eps*_zw,
 \qquad f_\eps=\rho_\eps*_zf.
\]
The constant-coefficient operator commutes with spatial convolution.
By \cref{lem:spatial-mollification-source},
$f_\eps\in L^1_tL^2_z$ for every $\eps>0$ and
$f_\eps\to f$ in $X$.  Also $w_\eps\in L^\infty_tL^2_z$, with compact time
support.  Therefore \cref{lem:fourier-section-identification} gives
\[
 w_\eps=T_\beta f_\eps.
\]
The bounded map $T_\beta:X\to X'$ proved in
\cref{prop:linear-carleman} yields
$T_\beta f_\eps\to T_\beta f$ in $X'$.  On the other hand,
$w_\eps\to w$ in distributions because spatial mollifiers approximate every
locally integrable function and $w\in X'\subset L^1_{\loc}$.  Uniqueness of
the distributional limit gives
\[
 w=T_\beta f.
\]
Applying the established operator bound and then taking the infimum over the
initial decomposition of $f$ proves \eqref{eq:graph-carleman}.
\end{proof}

\section{Exponential separation with a small potential}

\begin{lemma}[One-slab separation]\label{lem:one-slab}
There is an absolute $\eps_*>0$ with the following property.  Suppose
$U\in X'$, $U$ is compactly supported in time, $U=0$ for $z_1>H$, and
\begin{equation}
  PU=VU+R                                             \label{eq:one-slab-eq}
\end{equation}
in distributions, with $V\in Y$ and $\norm V_Y\le\eps_*$.  Assume that for
all sufficiently large $\beta$ the weighted distribution
$\e^{\beta z_1}R$ belongs to $X$ and
\begin{equation}
 \norm{\e^{\beta z_1}R}_X\longrightarrow0
 \quad(\beta\to+\infty).                            \label{eq:R-decay}
\end{equation}
This condition holds, for example, when $R\in X$ is supported in
$\{z_1\le-d\}$ for some $d>0$.  Then $U=0$ on $\{z_1>0\}$.
\end{lemma}

\begin{proof}
Apply \cref{lem:carleman-closure} and the multiplier estimate:
\begin{align*}
 \norm{\e^{\beta z_1}U}_{X'}
 &\le C\norm{V\e^{\beta z_1}U}_{X}
      +C\norm{\e^{\beta z_1}R}_{X}\\
 &\le C\norm V_Y\norm{\e^{\beta z_1}U}_{X'}
      +C\norm{\e^{\beta z_1}R}_{X}.
\end{align*}
Choose $\eps_*=(2C)^{-1}$ and absorb the first term.  On $z_1>0$ the weight
is at least one, hence
\[
 \norm{U}_{X'(\{z_1>0\})}
 \le2C\norm{\e^{\beta z_1}R}_{X}.
\]
Let $\beta\to\infty$.
\end{proof}

\begin{lemma}[Exponentially suppressed moment errors]\label{lem:moment-error}
Let $d>0$, let $A\in L^1_t$, and let
$U\in L^\infty_tL^2_z$ have compact time support.  Then
\begin{equation}
 \norm{\e^{\beta z_1}A(t)z_1U\1_{\{z_1\le-d\}}}_{L^1_tL^2_z}
 \longrightarrow0\quad(\beta\to\infty).           \label{eq:moment-error-decay}
\end{equation}
\end{lemma}

\begin{proof}
For $z_1\le-d$,
$|z_1|\e^{\beta z_1}\le C_d\e^{-\beta d/2}$ for all sufficiently large
$\beta$.  Therefore the displayed norm is bounded by
$C_d\e^{-\beta d/2}\norm A_{L^1_t}\norm U_{L^\infty_tL^2_z}$.
\end{proof}

\section{Moving the interface}
We now show how to manufacture the hypotheses of \cref{lem:one-slab}.

\begin{lemma}[Moving translation and Galilean gauge in distributions]
\label{lem:moving-galilean-distribution}
Let $J$ be bounded, let
$Z\in L^1_{\loc}(J\times\R^2)$, and suppose
\[
 PZ=WZ
 \quad\text{in }\D'(J\times\R^2),
\]
where $WZ\in L^1_{\loc}$.  Let $w\in C^2(J;\R)$ and define
\[
 V(t,y)=Z(t,y_1+w(t),y_2),
 \qquad a(t)=\frac{w'(t)}2,
\]
\begin{equation}
 U(t,y)=
 \exp\!\left(-ia(t)y_1-i\int_{t_*}^t a(\tau)^2\,\dd\tau\right)V(t,y).
 \label{eq:moving-gauge-definition}
\end{equation}
Then, in distributions,
\begin{align}
 PV&=W(t,y_1+w(t),y_2)V+iw'(t)\partial_{y_1}V,
                                                        \label{eq:moving-translation-distribution}\\
 PU&=\left(W(t,y_1+w(t),y_2)+\frac{w''(t)}2y_1\right)U. \label{eq:moving-gauge-distribution}
\end{align}
Spatial translation and the unimodular factor preserve the
$L^\infty_tL^2_y$ and $L^4_{t,y}$ norms.  They also carry the support
condition $Z(t,z)=0$ for $z_1>B$ to
$V(t,y)=U(t,y)=0$ for $y_1>B-w(t)$.
\end{lemma}

\begin{proof}
For a test function $\Phi\in C_c^\infty(J\times\R^2)$, make the
measure-preserving substitution $z_1=y_1+w(t)$ in the distributional pairing.
The pulled-back test is
$\widetilde\Phi(t,z)=\Phi(t,z_1-w(t),z_2)$ and satisfies
\[
 \partial_t\widetilde\Phi
 =\Phi_t(t,z_1-w,z_2)-w'(t)\Phi_{y_1}(t,z_1-w,z_2),
 \qquad
 \Delta_z\widetilde\Phi=\Delta_y\Phi(t,z_1-w,z_2).
\]
Insert these identities in $\pair{PZ}{\widetilde\Phi}$ and move the
$y_1$ derivative back onto $V$.  This gives
\eqref{eq:moving-translation-distribution}; all pairings are legitimate
because the test support is compact and $V,WV$ are locally integrable.

Write the phase in \eqref{eq:moving-gauge-definition} as $\e^{-i\Theta}$,
where $\Theta(t,y_1)=a(t)y_1+\int_{t_*}^ta(\tau)^2\,\dd\tau$.
Multiplication of a distribution by this smooth function obeys the ordinary
Leibniz rule.  Since
\[
 \Theta_{y_1}=a,
 \qquad \Theta_t=a'y_1+a^2,
\]
a direct expansion gives
\begin{align*}
 P(\e^{-i\Theta}V)
 &=\e^{-i\Theta}
   \bigl(PV+a'y_1V-2ia\partial_{y_1}V\bigr).
\end{align*}
Use $w'=2a$ in \eqref{eq:moving-translation-distribution}; the first-order
terms cancel and $a'=w''/2$, proving
\eqref{eq:moving-gauge-distribution}.  The norm and support statements are
immediate from the change of variables and the modulus-one phase.
\end{proof}

\begin{lemma}[Quantitative one-step propagation]\label{lem:one-step-propagation}
There is an absolute constant $c_0>0$ with the following property.  Let
$J=(a,b)$ have length $L$, let $0<\delta<L/8$, and let
$0<h\le c_0\sqrt\delta$.  Suppose
\begin{equation}
 Z\in C(J;L^2)\cap L^4(J\times\R^2),
 \qquad PZ=WZ,
 \qquad \norm W_{L^1_tL^\infty_z(J)}\le c_0,        \label{eq:step-hyp}
\end{equation}
and $Z=0$ on $J\times\{z_1>B\}$.  Then
\begin{equation}
 Z=0\quad\text{on }
 (a+4\delta,b-4\delta)\times\{z_1>B-h\}.           \label{eq:step-conclusion}
\end{equation}
\end{lemma}

\begin{proof}
\textbf{1. Choice of compatible cutoffs.}
Choose $\omega\in C^\infty([a,b])$ such that
\begin{align*}
 \omega(t)&=2
 &&\text{on }[a,a+2\delta]\cup[b-2\delta,b],\\
 \omega(t)&=-1
 &&\text{on }[a+3\delta,b-3\delta],\\
 -1\le\omega&\le2,
 \qquad |\omega'|\le C\delta^{-1},
 \qquad |\omega''|\le C\delta^{-2}.
\end{align*}
The derivatives may be supported in the two transition intervals of length
$\delta$.  Choose also $\eta\in C_c^\infty(J)$ with
\begin{align*}
 \eta&=0
 &&\text{on }[a,a+\delta]\cup[b-\delta,b],\\
 \eta&=1
 &&\text{on }[a+2\delta,b-2\delta],\\
 |\eta'|&\le C\delta^{-1},
\end{align*}
and with $\supp\eta'$ contained where $\omega=2$.  Set
$w(t)=B+h\omega(t)$.  Thus the moving boundary is $B+2h$ wherever the time
cutoff changes and is $B-h$ on the central interval.

Introduce translated coordinates $y_1=z_1-w(t)$, $y_2=z_2$, and
\[
 V(t,y)=Z(t,y_1+w(t),y_2).
\]
The known half-space vanishing implies
\begin{equation}
 V(t,y)=0\quad\text{when }y_1>B-w(t)=-h\omega(t).   \label{eq:V-upper-support}
\end{equation}
Since $-h\omega(t)\le h$, we have $V=0$ for $y_1>h$ at every time.  On
$\supp\eta'$, where $\omega=2$, we have the stronger vanishing
$V=0$ for $y_1>-2h$.

\textbf{2. Galilean gauge.}
Let $a(t)=w'(t)/2$ and define
\[
 U(t,y)=\exp\!\left(-ia(t)y_1-i\int_a^t a(s)^2\,\dd s\right)V(t,y).
\]
Apply \cref{lem:moving-galilean-distribution}.  It gives, in the exact
rough distributional class of $Z$,
\begin{equation}
 PU=\left(W(t,y_1+w(t),y_2)+\frac{w''(t)}2y_1\right)U. \label{eq:gauged-moving}
\end{equation}
It also proves that translation and the unimodular gauge preserve the
$L^\infty_tL^2_y$ and $L^4_{t,y}$ norms and the stated support inequalities.

\textbf{3. A bounded near coefficient and a weighted far error.}
Put $Y=\eta U$.  Then $Y\in X'$, it is compactly supported in time, and it
vanishes for $y_1>h$.  On the support of $U$ split the affine coefficient as
\[
 \frac{w''}2y_1=A_{\mathrm{near}}+A_{\mathrm{far}},
\]
where
\[
 A_{\mathrm{near}}
 =\frac{w''}2y_1\1_{\{-h<y_1\le h\}},
 \qquad
 A_{\mathrm{far}}
 =\frac{w''}2y_1\1_{\{y_1\le-h\}}.
\]
The value of the coefficient for $y_1>h$ is irrelevant because $U$ is zero
there.  Since $w''=h\omega''$ is supported on a set of total length
$O(\delta)$,
\begin{equation}
 \norm{A_{\mathrm{near}}}_{L^1_tL^\infty_y}
 \le \frac h2\int_J|w''(t)|\,\dd t
 \le C h^2\delta^{-1}.                            \label{eq:near-affine-small}
\end{equation}
The equation for $Y$ is
\begin{equation}
 PY=\widetilde VY+R,                                \label{eq:Y-moving-eq}
\end{equation}
where
\[
 \widetilde V(t,y)=W(t,y_1+w(t),y_2)+A_{\mathrm{near}}(t,y)
\]
and
\[
 R=\eta A_{\mathrm{far}}U+i\eta' U.
\]
Translation preserves the $L^1_tL^\infty$ norm.  Choose the absolute
constant $c_0$ so small that
$c_0+Cc_0^2\le\eps_*$.  The hypotheses
$\norm W_{L^1L^\infty}\le c_0$ and $h^2/\delta\le c_0^2$ then give
$\norm{\widetilde V}_Y\le\eps_*$.

\textbf{4. Exponential decay of both errors.}
For every $\beta>0$, the weighted far term belongs to $L^1_tL^2_y$.
Indeed, \cref{lem:moment-error} applies with
$A(t)=\eta(t)w''(t)/2$, because $A\in L^1_t$, $U\in L^\infty_tL^2_y$, and
the term is supported in $y_1\le-h$.  It also gives
\[
 \norm{\e^{\beta y_1}\eta A_{\mathrm{far}}U}_{L^1_tL^2_y}
 \longrightarrow0.
\]
On $\supp\eta'$ we have $U=0$ for $y_1>-2h$, and therefore
\[
 \norm{\e^{\beta y_1}\eta'U}_{L^1_tL^2_y}
 \le \e^{-2\beta h}\norm{\eta'}_{L^1_t}
      \norm U_{L^\infty_tL^2_y}
 \longrightarrow0.
\]
Thus $R$ satisfies \eqref{eq:R-decay}.

\textbf{5. Propagation across one slab.}
Apply \cref{lem:one-slab} with upper support bound $H=h$.  It gives
$Y=0$ for $y_1>0$.  On
$(a+4\delta,b-4\delta)$ we have $\eta=1$ and $w=B-h$, so $y_1>0$ is exactly
$z_1>B-h$.  This proves \eqref{eq:step-conclusion}.
\end{proof}

\section{Half-space propagation}

\begin{theorem}[Half-space uniqueness for integrably bounded potentials]
\label{thm:halfspace-UC}
Let $J$ be an open interval and let
\begin{equation}
 Z\in C(J;L^2(\R^2))\cap L^4_{\loc}(J\times\R^2),
 \qquad PZ=WZ,                                      \label{eq:halfspace-class}
\end{equation}
where $W\in L^1_{\loc}(J;L^\infty(\R^2))$.  If for some $B\in\R$
\begin{equation}
 Z=0\quad\text{almost everywhere on }J\times\{z_1>B\}, \label{eq:halfspace-zero}
\end{equation}
then $Z=0$ on $J\times\R^2$.
\end{theorem}

\begin{proof}
Fix $t_0\in J$ and a target distance $D>0$.  By absolute continuity of the
$L^1_tL^\infty_z$ norm, choose a symmetric interval
$I=(t_0-L/2,t_0+L/2)\Subset J$ on which
$\norm W_{L^1_tL^\infty_z(I)}\le c_0/2$.  We shall propagate the boundary
from $B$ to $B-D$ while keeping $t_0$ inside the surviving time interval.

Choose an integer $N$ and put
\[
 \delta=\frac{L}{32N},
 \qquad h=\frac{c_0}{2}\sqrt\delta.
\]
At each application of \cref{lem:one-step-propagation}, the time interval
loses at most $8\delta$ in total, so after $N$ applications it still contains
the central interval of length at least $3L/4$, and in particular contains
$t_0$.  The potential norm on every smaller interval is no larger than on
$I$.  The relation $h\le c_0\sqrt\delta$ is built into the choice.

After $N$ steps, $Z$ vanishes almost everywhere on a time neighborhood
of $t_0$ and on the region $z_1>B-Nh$.  Continuity into $L^2$ implies the
same vanishing at the exact time $t_0$.  But
\[
 Nh=\frac{c_0}{2}\sqrt{\frac{NL}{32}},
\]
which tends to infinity with $N$.  Choose $N$ so large that $Nh>D$.
Therefore $Z(t_0,z)=0$ for $z_1>B-D$.  Since $D$ was arbitrary,
$Z(t_0)=0$ in $L^2(\R^2)$.  Since $t_0$ was arbitrary, $Z=0$ throughout
$J\times\R^2$.
\end{proof}

\begin{corollary}[Half-space uniqueness in the $X,Y$ formulation]
\label{cor:halfspace-XY}
The conclusion of \cref{thm:halfspace-UC} remains valid if
$W\in Y_{\loc}$ and, around every time, there is a sufficiently short
interval on which $W$ admits a decomposition with $Y$ norm smaller than
$c_0/2$.  The case $W\in L^1_{\loc}L^\infty$ used in the main proof is the
first summand of this formulation.
\end{corollary}

\begin{proof}
The proof of \cref{lem:one-step-propagation} uses the potential only through
\cref{lem:XY-multiplier} and the smallness needed in
\cref{lem:one-slab}.  Replace the $L^1L^\infty$ norm by the corresponding
small $Y$ norm.
\end{proof}

\begin{checkpoint}
The half-space propagation theorem has been proved from an explicit
Fourier-kernel inverse for the conjugated free operator, HLS, norm
absorption, and a quantitative moving-interface construction.  No external
unique-continuation theorem is used.
\end{checkpoint}

\chapter{Completion of the phase-retrieval proof}\label{chap:completion}

\section{Vanishing of the exterior product}
Fix a bounded open interval $I$.  By \cref{prop:G-equation},
\[
 iG_t+G_{rr}+G_{ss}=QG,
 \qquad G=0\text{ for }s<0.
\]
The solution belongs to
$C(I;L^2_{r,s})\cap L^4(I\times\R^2)$, and by
\cref{lem:Q-bounds},
$Q\in L^1(I;L^\infty_{r,s})$.

Reflect the normal variable:
\begin{equation}
 z_1=-s,\qquad z_2=r,
 \qquad Z(t,z_1,z_2)=G(t,z_2,-z_1),                 \label{eq:Z-reflection}
\end{equation}
and set
\[
 W(t,z_1,z_2)=Q(t,z_2,-z_1).
\]
Then
\[
 (i\partial_t+\Delta_z)Z=WZ,
 \qquad Z=0\quad\text{for }z_1>0.
\]
All hypotheses of \cref{thm:halfspace-UC} hold, so $Z=0$ on
$I\times\R^2$.  Consequently $G=0$ for $s>0$.  Since $F$ is odd in $s$, it
vanishes for almost every $(t,r,s)\in I\times\R^2$.  The continuity
$F\in C(I;L^2)$ then upgrades this almost-everywhere-in-time statement to
\begin{equation}
  F=0\quad\text{in }C(I;L^2(\R^2)).                \label{eq:F-zero-I}
\end{equation}
The interval $I$ was arbitrary, hence
\begin{equation}
  F(t)=0\quad\text{in }L^2(\R^2)
  \quad\text{for every }t\in\R.                   \label{eq:F-zero-all}
\end{equation}

\section{The Hilbert-space exterior-product lemma}

\begin{lemma}[Zero exterior product means linear dependence]
\label{lem:Hilbert-wedge}
Let $f,g\in L^2(\R)$ and assume
\begin{equation}
 f(x)g(y)-g(x)f(y)=0
 \quad\text{in }L^2(\R^2).                         \label{eq:Hilbert-wedge-zero}
\end{equation}
If $f\ne0$, there is a unique $c\in\C$ with $g=cf$ in $L^2$.  If, in
addition, $\norm f_2=\norm g_2$, then either $f=g=0$, or this scalar satisfies
$|c|=1$.  Equivalently, under the equal-norm hypothesis there always exists
$\zeta\in\C$ with $|\zeta|=1$ and $g=\zeta f$.
\end{lemma}

\begin{proof}
Assume first that $f\ne0$.  The contraction map
\[
 \mathcal C_f:L^2(\R^2)\to L^2(\R),
 \qquad
 (\mathcal C_fH)(y)=\int_\R\overline{f(x)}H(x,y)\,\dd x
\]
is bounded with norm at most $\norm f_2$, by Cauchy--Schwarz and Fubini.
Applying it to \eqref{eq:Hilbert-wedge-zero} yields
\[
 \norm f_2^2g-\left(\int_\R\overline f g\right)f=0.
\]
Thus $g=cf$, where
$c=\pair f g/\norm f_2^2$; uniqueness follows from $f\ne0$.  If the norms are
equal and nonzero, then
$|c|\norm f_2=\norm g_2=\norm f_2$, so $|c|=1$.
If $f=0$ and the norms are equal, then $g=0$ as well, and one may take
$\zeta=1$.
\end{proof}

\section{Proof of the main theorem}

\begin{proof}[Proof of \cref{thm:main}]
Evaluate \eqref{eq:F-zero-all} at $t=0$.  If $u(0)=0$, then
\eqref{eq:all-time-density} gives $v(0)=0$, and uniqueness of the global flow
shows $u=v=0$.

Assume $u(0)\ne0$.  By \cref{lem:Hilbert-wedge}, there is $\zeta\in\C$ with
$v(0)=\zeta u(0)$.  Equal densities imply equal $L^2$ norms, so
$|\zeta|=1$.  Gauge covariance \eqref{eq:gauge-cov} and uniqueness of the
global mild flow give
\[
 v(t)=\Phi_t(v(0))
     =\Phi_t(\zeta u(0))
     =\zeta\Phi_t(u(0))
     =\zeta u(t)
\]
for every $t\in\R$.
\end{proof}

\section{Immediate corollaries}

\begin{corollary}[Equality on any time interval]\label{cor:interval-version}
Let $u,v$ be global $L^2$ solutions.  If $|u|=|v|$ almost everywhere on
$J\times\R$ for one nonempty open time interval $J$, then there is
$\zeta\in\T$ such that $v=\zeta u$ on all of $\R\times\R$.
\end{corollary}

\begin{proof}
Repeat the proof on $J$.  It gives $F(t)=0$ for every $t\in J$, hence
$v(t_0)=\zeta u(t_0)$ at one $t_0\in J$.  Global uniqueness and gauge
covariance propagate the relation both forward and backward.
\end{proof}

\begin{corollary}[The finite-energy case]\label{cor:H1-case}
The conclusion holds in particular for solutions in $C(\R;H^1(\R))$.
No separate one-dimensional lateral unique-continuation theorem is required.
\end{corollary}

\begin{proof}
An $H^1$ solution is the $L^2$ mild solution with the same data by
persistence and uniqueness, so \cref{thm:main} applies.
\end{proof}

\begin{corollary}[Rigidity of a common real potential]\label{cor:linear-common-potential}
Let $V\in L^1_{\loc}(\R;L^\infty(\R))$ be real.  Suppose
$u,v\in C(\R;L^2)\cap L^4_{\loc}L^\infty$ are mild solutions of
\[
 iu_t+u_{xx}=Vu,
 \qquad iv_t+v_{xx}=Vv,
\]
have the same modulus almost everywhere, and possess the local
half-derivative smoothing needed for the quadratic current.  Then
$v=\zeta u$ for one constant unimodular $\zeta$.
\end{corollary}

\begin{proof}
The continuity equation is unchanged because $V$ is real.  The Wronskian,
exterior-product, trace, and half-space arguments then give, at any fixed
time $t_0$, a scalar $\zeta\in\T$ with $v(t_0)=\zeta u(t_0)$.  The difference
$w=v-\zeta u$ satisfies the mild equation
\[
 w(t)=-i\int_{t_0}^t S(t-s)V(s)w(s)\,\dd s.
\]
Unitarity gives
\[
 \norm{w(t)}_2
 \le\int_{t_0}^t\norm{V(s)}_\infty\norm{w(s)}_2\,\dd s
\]
with the integration orientation adjusted when $t<t_0$.  The elementary
Gronwall argument on compact intervals yields $w=0$ in both time directions.
\end{proof}

\begin{checkpoint}
The main result is proved.  The remaining chapters are verification aids:
they isolate approximation details, state the exact foundational boundary,
and give a theorem-by-theorem formalization order.
\end{checkpoint}

\appendix

\chapter{Approximation and distributional bookkeeping}\label{app:bookkeeping}

\section{Why smooth approximation is simultaneous}
The proof uses one sequence of smooth solutions to justify several identities.
Here is the construction in a form suitable for formalization.  Choose
$u_{0,n},v_{0,n}\in\Ssch(\R)$ with convergence in $L^2$.  Let $u_n,v_n$ be
the global solutions.  On every compact interval $I$, the stability estimates
proved in \cref{chap:wellposed} give
\begin{align*}
 u_n&\to u,
 &v_n&\to v
 &&\text{in }C_tL^2_x\cap L^6_{t,x},\\
 u_n&\to u,
 &v_n&\to v
 &&\text{in }L^4_tL^\infty_x,\\
 |u_n|^2u_n&\to|u|^2u,
 &|v_n|^2v_n&\to|v|^2v
 &&\text{in }L^1_tL^2_x,\\
 \chi u_n&\to\chi u,
 &\chi v_n&\to\chi v
 &&\text{in }L^2_tH^{1/2}_x.
\end{align*}
The last convergence holds after every compact spatial cutoff $\chi$.
The first line follows from the contraction argument; the second from the
endpoint Duhamel estimate; the third from the cubic difference inequality;
and the fourth from forced local smoothing.  Thus the continuity equation,
quadratic current, exterior-product equation, and normal trace all pass to the
limit along the same approximating sequence.  For the exterior equation one
passes first to the unfactored four-term identity
\eqref{eq:F-unfactored}; equality of the two limiting densities is used only
after that passage to factor its right-hand side as $QF$.

\section{Density of tensor-product tests}

\begin{lemma}[Finite-derivative density of separated tests]
\label{lem:tensor-test-density}
Let $K\Subset\R^{m+n}$ be compact, let $M\in\N$, and let
$\Psi\in C_c^\infty(\R^m\times\R^n)$ have support in $K$.  There are finite
sums
\[
 \Psi_N(t,x)=\sum_{j=1}^{J_N}\eta_{N,j}(t)\phi_{N,j}(x),
 \qquad
 \eta_{N,j}\in C_c^\infty(\R^m),\quad
 \phi_{N,j}\in C_c^\infty(\R^n),
\]
supported in one fixed compact rectangle, such that
\begin{equation}
 \max_{|\alpha|\le M}
 \norm{\partial^\alpha(\Psi_N-\Psi)}_\infty\longrightarrow0.
 \label{eq:tensor-test-density}
\end{equation}
\end{lemma}

\begin{proof}
Choose open product rectangles
\[
 K\subset R_0\Subset R_1\Subset R_2
\]
and product cutoffs
$\kappa(t,x)=\kappa_t(t)\kappa_x(x)$ such that
$\kappa=1$ on $R_0$ and $\supp\kappa\subset R_1$.  The original function
$\Psi$ is identically zero near the boundary of $R_2$.  Regard its
restriction to $R_2$ as a smooth function on the flat torus obtained by
identifying opposite faces of $R_2$; the periodic extension, denoted
$\widetilde\Psi$, agrees with $\Psi$ throughout $R_2$.

Write its Fourier series as
\[
 \widetilde\Psi(t,x)
 =\sum_{k\in\mathbb Z^m,\,\ell\in\mathbb Z^n}
 c_{k,\ell}\,
 \e^{2\pi i k\cdot t/L_t}\e^{2\pi i\ell\cdot x/L_x},
\]
with the evident componentwise side lengths.  For every integer $N_0$,
integration by parts $N_0$ times in a coordinate for which
$|(k,\ell)|$ is maximal gives
\begin{equation}
 |c_{k,\ell}|
 \le C_{N_0}(1+|k|+|\ell|)^{-N_0}.               \label{eq:tensor-fourier-decay}
\end{equation}
Choose $N_0>M+m+n+1$.  After applying any derivative of order at most $M$,
the resulting Fourier series is absolutely and uniformly convergent by
\eqref{eq:tensor-fourier-decay}.  Hence the rectangular partial sums $P_N$
converge to $\widetilde\Psi$ in the $C^M$ norm on $R_2$.

Finally set $\Psi_N=\kappa P_N$, extending it by zero outside $R_2$.  Since
$\supp\kappa\Subset R_2$, this extension is smooth.  Each Fourier monomial
in $P_N$ is a product of a function of $t$ and a function of $x$, and the
cutoff is itself a product.  Thus $\Psi_N$ is a finite sum of separated
compactly supported smooth functions, all supported in the fixed rectangle
$\supp\kappa$.  The product rule and $C^M$ convergence of $P_N$ give
\eqref{eq:tensor-test-density}; indeed
\[
 \kappa\widetilde\Psi=\kappa\Psi=\Psi,
\]
because $\kappa=1$ on the support of $\Psi$.
\end{proof}

This lemma is the only tensor-test density input used in the proof.  For the
current identity take $m=n=1$ and $M=1$; for the zero-extension equation take
$m=2,n=1$ and $M=2$.  No nuclear-space tensor-product theorem is required.

\section{Slicing an \texorpdfstring{$L^1$}{L1} family of distributions}
Fix a compact spatial interval $K$.  Suppose $t\mapsto\Lambda_t$ is a
measurable family of distributions for which
\[
 |\Lambda_t(\phi)|\le A(t)
   \bigl(\norm\phi_\infty+\norm{\phi'}_\infty\bigr),
 \qquad A\in L^1(I),
\]
for every smooth test supported in $K$.  Choose a countable family dense in
the $C^1$ norm among such tests.  If
\[
 \int_I\eta(t)\Lambda_t(\phi)\,\dd t=0
\]
for every $\eta\in C_c^\infty(I)$ and every member of that family, then the
scalar fundamental lemma gives one null set for each test.  Remove their
countable union and extend by the displayed continuity bound.  Thus
$\Lambda_t=0$ as a distribution for almost every $t$.  Countable compact
exhaustion yields one common full-measure set.  This is the mechanism used
for the logarithmic derivative and Wronskian identities; it does not require
an unjustified $H^{1/2}$-dual bound on the test variable.

\section{Distributional products appearing in the exterior equation}
The product $QF$ is an actual $L^2$ function on compact time intervals, not a
formal distributional multiplication.  A typical term is
\[
 |u(t,x)|^2u(t,x)v(t,y).
\]
Its $L^2_{t,x,y}$ norm factors exactly as
\[
 \norm v_{L^\infty_tL^2_y}
 \left(\int_I\norm{u(t)}_6^6\,\dd t\right)^{1/2}.
\]
Therefore every passage to the limit in \eqref{eq:F-equation} can be made in
$L^2$, which is stronger than distributional convergence.

\section{Why the zero extension has no hidden time or tangential jumps}
The multiplier $\1_{\{s>0\}}$ is independent of $t$ and $r$, so
$\partial_t$ and $\partial_r^2$ commute with it in distributions.  Only
$\partial_s^2$ produces jump terms.  Scalarization by $\psi(t,r)$ reduces the
statement to \cref{lem:jump-formula}.  Because the scalarized function lies
in $H^2_{\loc}(s)$, both traces are genuine numbers and there is no ambiguity
about the coefficients of $\delta_0$ and $\delta_0'$.

\chapter{Detailed norm calculations}\label{app:norms}

\section{Interpolation to \texorpdfstring{$L^8_tL^4_x$}{L8t L4x}}
For completeness, complex or real interpolation of the two mixed-norm
endpoints gives
\[
 \norm u_{L^8_tL^4_x}
 \le\norm u_{L^\infty_tL^2_x}^{1/2}
      \norm u_{L^4_tL^\infty_x}^{1/2}.
\]
This can also be checked pointwise without an interpolation theorem:
$\norm{u(t)}_4^4\le\norm{u(t)}_2^2\norm{u(t)}_\infty^2$.  Squaring and
integrating in time gives
\[
 \int_I\norm{u(t)}_4^8\,\dd t
 \le\norm u_{L^\infty_tL^2_x}^4
      \int_I\norm{u(t)}_\infty^4\,\dd t.
\]

\section{The ridge strip identity}
For fixed $s$ and $t$, the substitution
$x=(r+s)/\sqrt2$ gives $\dd r=\sqrt2\,\dd x$.  Therefore
\begin{align*}
 &\int_a^{a+h}\int_\R
 \left|\rho\!\left(t,\frac{r+s}{\sqrt2}\right)\right|^2
 \,\dd r\,\dd s\\
 &\qquad=\int_a^{a+h}\sqrt2\int_\R|u(t,x)|^4\,\dd x\,\dd s
 =\sqrt2h\norm{u(t)}_4^4.
\end{align*}
Integrating in time gives the exact identity behind \eqref{eq:Q-strip}.

\section{The exterior forcing}
Expanding $QF$ yields eight tensor monomials, but only two forms occur up to
renaming:
\[
 |u(x)|^2u(x)v(y),
 \qquad
 |u(x)|^2v(x)u(y).
\]
Because $|u|=|v|$, both have the same norm.  For the second,
\begin{align*}
 &\iint_{\R^2}|u(x)|^4|v(x)|^2|u(y)|^2\,\dd x\,\dd y\\
 &\qquad=\norm u_2^2\int_\R|u(x)|^6\,\dd x.
\end{align*}
Thus the common estimate uses only mass and $L^6_{t,x}$.

\section{The moving-interface smallness scale}
In \cref{lem:one-step-propagation}, a boundary displacement of size $h$ over
a transition time $\delta$ has
\[
 |w'|\lesssim h\delta^{-1},
 \qquad |w''|\lesssim h\delta^{-2}.
\]
The near affine coefficient is restricted to $|y_1|\le h$, and its time
support has length $O(\delta)$.  Hence
\[
 \int |w''(t)y_1|\,\dd t
 \lesssim (h\delta^{-2})h\delta
 =h^2\delta^{-1}.
\]
Choosing $h\lesssim\sqrt\delta$ makes this perturbation uniformly small.
After $N$ steps with $\delta\simeq L/N$, the total distance is
$Nh\simeq\sqrt{NL}$, which tends to infinity.  This is the quantitative
reason the loss of time at each step does not obstruct propagation through an
arbitrary spatial distance.


\chapter{Resolution of the formalization-readiness audit}
\label{app:audit-resolution}

The earlier draft was audited against a stronger standard than ordinary
journal exposition: a later translator should not have to invent a missing
analytic argument or silently choose a non-equivalent representation.  The
following table records where each blocking item is closed in this edition.
The table is an index only; the cited numbered results, not this table, carry
the proofs.

\begin{longtable}{p{0.075\textwidth}p{0.285\textwidth}p{0.535\textwidth}}
\toprule
Item & Earlier obstruction & Closure in this edition\\
\midrule
\endhead
B01 & No pinned formal boundary. & Lean/Mathlib are pinned in
\cref{app:formalization-contract}; the Fourier rescaling and bootstrap
entry points are stated explicitly.\\
B02 & Mixed norms, sums, restrictions, mild solutions, and traces had no
fixed representation. & Scalar mixed norms, their quotient convention,
completeness, restriction, sum/intersection norms, the completed retarded
operator, local Sobolev cutoffs, and scalarized traces are defined in
\cref{chap:mixed-norms,lem:retarded-continuous,chap:traces}.\\
B03 & The Fourier--Gagliardo equivalence was only summarized. &
\Cref{thm:fourier-gagliardo} prove both directions,
all changes of variables, mollification, cutoff removal, and density.\\
B04 & The kernel and $L^2$ propagators, Gaussian formula, and interpolation
were not connected. & \Cref{lem:complex-gaussian-kernel,lem:propagator-compatibility,lem:three-lines-explicit,lem:riesz-thorin-used} give the compatible
realizations and prove the interpolation input.\\
B05 & Endpoint $(4,\infty)$ duality used an invalid blanket Bochner-dual
identification. & \Cref{lem:measurable-esssup-section,prop:endpoint-norming,thm:hom-strichartz} use jointly measurable scalar
sections and a countable norming class; no full $L^\infty$ dual is invoked.\\
B06 & $H^1$ and smooth persistence were sketches. &
\Cref{lem:Wp-difference-quotient,lem:H1-persistence,prop:Hk-persistence}
construct the differentiated fixed points, isolate the affine top derivative,
and carry out the induction.\\
B07 & The mass derivative in the Gelfand triple was asserted. &
\Cref{lem:gelfand-chain-rule,lem:mass-H1} prove the chain rule by time
mollification and apply it with all pairings displayed.\\
B08 & Continuous spatial slices lacked representative and subsequence
bookkeeping. & \Cref{lem:joint-representative-L2,lem:mixed-subsequence-uniform,lem:H1-C0,lem:C0-limit-identification,cor:C0-slices} provide the complete construction.\\
B09 & The zero-set regularizer used a false infinite-measure $L^2$
argument. & \Cref{lem:Teps} now uses
$|T_\eps(f)|\le |f|$ and dominated convergence, followed by dominated
convergence of the Gagliardo integrand.\\
B10 & The logarithmic form was tested against a nonsmooth multiplier without
a continuity theorem. & \Cref{def:multiplier-test-algebra,lem:quadratic-form-extension} give the precise test norm, continuity bound,
and simultaneous convergence statement.\\
B11 & Scalar $h,h''\in L^2_{\loc}\Rightarrow h\in H^2_{\loc}\subset C^1$
was imported informally. & \Cref{lem:interior-first-derivative,thm:local-H2-characterization} prove the interior estimate, mollified weak
limit, and one-dimensional continuous representatives.\\
B12 & The Neumann trace limit omitted localization, multiplier, and Bochner
steps. & \Cref{prop:Neumann-trace} fixes nested cutoffs, proves strong
$L^2_tH^{-1/2}$ convergence, controls translated tests, and identifies the
limit with the sliced Wronskian.\\
B13 & Tensor tests, slicing, and jump assembly lacked a concrete topology. &
\Cref{lem:tensor-test-density,thm:current-recovery,prop:G-equation} use
finite derivative sup-seminorms and countable dense families, with all
continuity bounds displayed.\\
B14 & The Carleman graph closure was one paragraph. &
\Cref{lem:spatial-mollification-source,lem:parameterwise-distribution-ODE,lem:compact-support-ODE,lem:fourier-section-identification,lem:carleman-closure} prove the source decomposition, frequency slicing,
compact-support inverse, and limit passage in sequence.\\
B15 & Endpoint mixed-norm measurability recurred unresolved throughout. &
\Cref{chap:mixed-norms} fixes one scalar representation globally and proves
all operations used later, including scalar Minkowski and endpoint
completeness.\\
B16 & The earlier readiness claims exceeded the supplied detail. & The front
matter now distinguishes a mathematical proof specification from a Lean
kernel certificate, and \cref{app:formalization-contract} states a
three-way no-invention gate for imported, defined, and proved declarations.\\
\bottomrule
\end{longtable}

Three additional omissions found during the second pass are also closed.
First, \cref{lem:retarded-continuous} distinguishes the Strichartz-completion
operator used in the fixed point from the a posteriori $L^2$-valued Bochner
integral.  Second, the proof of \cref{thm:LWP} now establishes uniqueness in
the entire solution class, not merely inside the contraction ball.  Third,
\cref{lem:complex-gaussian-kernel,lem:three-lines-explicit} supply the
Gaussian and three-lines calculations that interpolation and the dispersive
kernel require.

\chapter{Formalization contract and exact foundational boundary}
\label{app:formalization-contract}

This chapter is part of the proof package.  It fixes the version boundary,
the representations of all nontrivial objects, and the declarations that are
allowed to remain in the Mathlib layer.  It is intended to prevent an
autoformalizer from replacing a missing analytic theorem by an axiom.

\section{Pinned baseline}

The target baseline is
\[
 \text{Lean }4.32.0,
 \qquad
 \text{Mathlib tag }\texttt{v4.32.0},
 \qquad
 \text{release commit }\texttt{81a5d25}.
\]
A project may use a top-level \texttt{import Mathlib}, but it must pin that
revision in \texttt{lakefile.toml} and pin
\texttt{leanprover/lean4:v4.32.0} in \texttt{lean-toolchain}.  The proof is
mathematically independent of declaration names, but a reproducible
translation is not independent of versions.

The two project files are therefore fixed to the following literal contents.
The package name is inessential, but the toolchain and revision are not.
\begin{verbatim}
# lean-toolchain
leanprover/lean4:v4.32.0
\end{verbatim}
\begin{verbatim}
# lakefile.toml
name = "CubicNLSPhaseRetrieval"
version = "0.1.0"
defaultTargets = ["CubicNLSPhaseRetrieval"]

[[lean_lib]]
name = "CubicNLSPhaseRetrieval"

[[require]]
name = "mathlib"
git = "https://github.com/leanprover-community/mathlib4.git"
rev = "v4.32.0"
\end{verbatim}
\begin{verbatim}
-- CubicNLSPhaseRetrieval.lean
import Mathlib
\end{verbatim}
Here and below the word ``imported'' means available after that single import
at the pinned revision; it never means postulated by a new axiom.

At that baseline the intended Fourier entry points are the existing Mathlib
declarations
\begin{verbatim}
MeasureTheory.Lp.fourierTransform_l_i
MeasureTheory.Lp.norm_fourier_eq
MeasureTheory.Lp.inner_fourier_eq
MeasureTheory.Lp.fourier_toTemperedDistribution_eq
TemperedDistribution.MemSobolev
TemperedDistribution.memSobolev_two_iff_fourier
TemperedDistribution.MemSobolev.lineDerivOp
TemperedDistribution.MemSobolev.laplacian
\end{verbatim}
where the first displayed identifier is written in Lean with the unicode
subscript characters used by Mathlib.  These are bootstrap entry points, not
axioms: the project must run \texttt{\#check} on every identifier at the
pinned revision before introducing any wrapper theorem.  The rescaling
from Mathlib's Fourier character to the convention of this manuscript is
\eqref{eq:mathlib-fourier-rescaling}.

\section{Allowed foundational layer}

The following are the only theorem families treated as imported
foundations.  They are not PDE-specific.
\begin{enumerate}
\item real and complex normed-field algebra, conjugation, finite-dimensional
      Euclidean spaces, and elementary inequalities;
\item Lebesgue measure, product measure, almost-everywhere equality,
      measurability, simple-function approximation, and completeness of
      scalar $L^p$ spaces;
\item Lebesgue and Bochner integration, Fubini--Tonelli, Fatou, dominated and
      monotone convergence, H\"older, Minkowski, and Young convolution
      inequalities;
\item complete normed spaces, continuous linear maps, quotient and product
      Banach spaces, Hilbert-space duality, weak compactness in Hilbert space,
      and Banach's fixed-point theorem;
\item the fundamental theorem for one-dimensional $W^{1,1}$ functions,
      standard approximate identities, and density of smooth compactly
      supported functions in finite-exponent $L^p$ spaces;
\item the $L^1$ and $L^2$ Fourier transforms, Plancherel, Fourier inversion
      on Schwartz functions, and the maximum-modulus principle for holomorphic functions on rectangles.
      The needed three-lines estimate is proved as
      \cref{lem:three-lines-explicit}.
\end{enumerate}
Every use of these items in the manuscript is accompanied by the hypotheses
needed to apply it.  In particular, Strichartz, local smoothing, Sobolev
multiplication, trace regularity, Carleman estimates, and unique continuation
are not in the imported layer.

\section{Project-defined scalar mixed norms}

The Lean project should not define $L^4_tL^\infty_x$ as a Bochner space.
Instead it should bundle the scalar construction of
\cref{chap:mixed-norms}.  A suitable schematic interface is
\begin{verbatim}
structure ScalarMixed (p : ENNReal) (q : ENNReal) where
  toFun : Time * Space -> C
  aestronglyMeasurable : AEStronglyMeasurable toFun
  finiteNorm : mixedNorm p q toFun < infinity

-- Quotient by product-a.e. equality.
-- Sections are toFun (t, x), not values in an L-infinity Banach space.
-- For q = infinity, mixedNorm uses t |-> essSup_x |toFun (t, x)|.
\end{verbatim}
The implementation must prove, before Strichartz is started:
\begin{enumerate}
\item measurability of the section essential supremum
      (\cref{lem:measurable-esssup-section});
\item completeness (\cref{lem:scalar-endpoint-complete});
\item restriction and scalar Minkowski
      (\cref{lem:scalar-mixed-minkowski});
\item the countable norming identity
      (\cref{prop:endpoint-norming});
\item the almost-everywhere uniform subsequence lemma
      (\cref{lem:mixed-subsequence-uniform}).
\end{enumerate}
No theorem about the full dual of a Bochner $L^\infty$ space is required.

\section{Distributions used in the proof}

For an open set $\Omega\subset\R^n$, a complex distribution is a
complex-linear functional on $C_c^\infty(\Omega)$ such that for every compact
$K\Subset\Omega$ there are $m$ and $C_K$ with
\[
 |T(\phi)|\le C_K\max_{|\alpha|\le m}\norm{\partial^\alpha\phi}_\infty
 \qquad(\supp\phi\subset K).
\]
This compact-by-compact definition is sufficient for every distribution in
the manuscript.  The relevant orders and bounds are explicit:
\begin{longtable}{p{0.29\textwidth}p{0.17\textwidth}p{0.44\textwidth}}
\toprule
Distribution & Order needed & Continuity bound\\
\midrule
\endhead
$f_x$, $f\in H^{1/2}_{\loc}$ & one &
$|\pair{f_x}{\phi}|\le C_K\norm{\chi f}_{H^{1/2}}
\norm\phi_{H^{1/2}}$\\
$\overline f f_x$ & one &
\eqref{eq:quadratic-extension-bound} for smooth tests, then extension\\
$\partial_t|u|^2+2\partial_x\mathfrak j_u$ & one &
\eqref{eq:current-test-continuity} and $C_tL^2$ control\\
$PF-QF$ & two &
$F,QF\in L^2_{\loc}$ paired after moving derivatives to the test\\
$H_\psi''$ & two in $s$ &
\eqref{eq:Hpsi-second} and Cauchy--Schwarz\\
$P(\1_{s>0}F)-Q\1_{s>0}F$ & two &
scalar jump identity plus tensor-test density\\
\bottomrule
\end{longtable}

A Lean implementation may use Mathlib's test-function and distribution
structures, but it should expose lemmas with the preceding concrete bounds.
This avoids relying on an unspecified topological tensor-product theorem.

\section{Tensor tests and slicing}

The finite-derivative approximation theorem required by the distributional
arguments is already proved as \cref{lem:tensor-test-density}.  The Lean
project should instantiate that theorem with the derivative order appearing
in the concrete continuity bound of the functional at hand.  This avoids an
unspecified topological tensor-product theorem and makes the approximation
metric explicit.

For slicing, the project should use the countable argument of the proof of
\cref{thm:current-recovery}: on each compact interval choose a countable
family dense in the finite test-function seminorm appearing in
\eqref{eq:sliced-current-test-bound}, remove the union of scalar null sets,
and extend by that concrete continuity estimate.  No identification with the
dual of $H^{1/2}$ and no choice of a distribution-valued representative at
every time is needed.

\section{Local Sobolev objects and traces}

Define $H^s_{\loc}$ by cutoffs: a distribution $f$ belongs locally when
$\chi f\in H^s$ for every $\chi\in C_c^\infty$.  All statements are invariant
under the cutoff because two cutoffs equal to one near a test support produce
the same distributional pairing.  The scalar trace object is not a trace of
an arbitrary $L^2$ function.  It is the map
\[
 \psi\longmapsto H_\psi\in H^2_{\loc}(\R_s)
\]
constructed from the equation.  Dirichlet and Neumann traces are then the
continuous linear functionals
\[
 \psi\longmapsto H_\psi(0),
 \qquad
 \psi\longmapsto H_\psi'(0).
\]
The local $H^2$ theorem needed to define them is
\cref{thm:local-H2-characterization}; the Neumann identification is the
strong difference-quotient proof in \cref{prop:Neumann-trace}.  Thus the
project need not import a general hypersurface trace theorem.

\section{Carleman graph objects}

The spaces $X$, $X'$, and $Y$ are sums and intersections of scalar spaces,
not opaque notation.  Their completeness follows from
\cref{lem:sum-intersection-complete}.  The conjugated inverse $T_\beta$ is
first defined on compactly supported smooth forcing by
\eqref{eq:Tbeta-def} and then extended using the four operator estimates in
\cref{prop:linear-carleman}.  The graph-closure theorem must be implemented
through the following compiled sequence, with no abstract uniqueness axiom:
\begin{align*}
 &\text{source mollification}
   \Longrightarrow \text{Fourier sections}
   \Longrightarrow \text{scalar compact-support ODE},\\
 &w_\eps=T_\beta f_\eps
   \Longrightarrow w=T_\beta f.
\end{align*}
These are precisely
\cref{lem:spatial-mollification-source,lem:compact-support-ODE,lem:fourier-section-identification,lem:carleman-closure}.

\section{No-invention gate}

Before a module is accepted, every declaration in it must fall into one of
three classes:
\begin{enumerate}
\item an exact theorem imported from the pinned Mathlib revision;
\item a definition whose representation is fixed in this chapter;
\item a theorem proved from earlier accepted declarations and corresponding
to a numbered result in this manuscript.
\end{enumerate}
Declarations that merely restate a central analytic step, such as an axiom
named \texttt{endpoint\_strichartz}, \texttt{diagonal\_trace}, or
\texttt{halfspace\_unique\_continuation}, fail the gate.  Temporary
\texttt{sorry} terms may be used during interactive development only if the
build target used for certification rejects them.

\chapter{A formalization ledger relative to Mathlib}\label{app:ledger}

This appendix records the dependency order after the exact version and
representation contract of \cref{app:formalization-contract}.  It adds no
assumptions and separates definitions from theorem obligations.

\section{Foundational imports}
The project is pinned as in \cref{app:formalization-contract}.  It may use
Mathlib's $L^2$ Fourier transform and tempered-distribution Sobolev
infrastructure, after applying the rescaling
\eqref{eq:mathlib-fourier-rescaling}.  Scalar endpoint mixed norms,
Schr\"odinger propagators, local Sobolev quadratic forms, diagonal traces,
and the Carleman spaces remain project-defined because their exact
representations are part of the theorem statement and proof.

\section{Suggested module order}
\begin{longtable}{p{0.10\textwidth}p{0.28\textwidth}p{0.52\textwidth}}
\toprule
No. & Module & Main declarations and proof obligations\\
\midrule
\endhead
1 & \texttt{FourierSobolev} & Fourier transform convention; multipliers
$\langle D\rangle^s$, $|D|^s$; $H^{\pm1/2}$; translation difference
quotients; smooth multiplication.\\
2 & \texttt{FreeSchrodinger} & $S(t)$ as a unitary multiplier group;
Schwartz differentiation; dispersive kernel.\\
3 & \texttt{FractionalIntegration} & Maximal function, Hedberg inequality,
one-dimensional HLS, and the Christ--Kiselev decomposition.\\
4 & \texttt{Strichartz1D} & Homogeneous $(6,6)$ and $(4,\infty)$ estimates;
retarded $L^{6/5}_{t,x}\to L^6_{t,x}$ and $L^{6/5}_{t,x}\to L^\infty_tL^2_x$ estimates.\\
5 & \texttt{CubicFlow} & Fixed point in
$C_tL^2\cap L^6_{t,x}$; uniqueness; stability; mass conservation; global
iteration; gauge covariance.\\
6 & \texttt{EndpointRegularity} & Cubic forcing in $L^1_tL^2_x$;
$L^4_tL^\infty_x$; $L^8_tL^4_x$; smooth approximation convergence.\\
7 & \texttt{LocalSmoothing1D} & Pointwise-in-space Fourier identity;
cutoff commutator; forced local smoothing.\\
8 & \texttt{CriticalQuadraticForms} & Gagliardo product estimate;
Lipschitz composition; $T_\eps$ regularizer; definition and stability of
$\overline f f_x$.\\
9 & \texttt{DensityCurrent} & Continuity equation by approximation;
weak localized mass; equal density implies equal logarithmic form.\\
10 & \texttt{Wronskian} & Critical zero-set-safe Wronskian theorem and
almost-everywhere slicing.\\
11 & \texttt{ExteriorProduct} & Tensor continuity; two-particle equation;
coordinate rotation; ridge-potential estimates.\\
12 & \texttt{DiagonalTrace} & Scalarized $H^2$ regularity; translation
quotient computation; jump formula; zero extension.\\
13 & \texttt{Carleman2D} & Explicit conjugated inverse; kernel estimates;
$X$--$X'$ Carleman inequality; graph closure; one-slab absorption.\\
14 & \texttt{HalfspacePropagation} & Moving translation and Galilean gauge;
quantitative iteration.\\
15 & \texttt{PhaseRetrieval} & Reflection, vanishing exterior product,
Hilbert contraction, unimodular scalar, gauge propagation.\\
\bottomrule
\end{longtable}

\section{No-axiom dependency graph}
The proof graph is acyclic:
\[
\begin{gathered}
 \text{Fourier}+\text{HLS}
 \Longrightarrow \text{Strichartz}
 \Longrightarrow \text{global cubic flow},\\
 \text{Fourier smoothing}+\text{flow stability}
 \Longrightarrow \text{critical current},\\
 \text{critical calculus}+\text{current}
 \Longrightarrow \text{Wronskian},\\
 \text{Wronskian}+\text{tensor equation}
 \Longrightarrow \text{zero diagonal Cauchy data},\\
 \text{Carleman kernel}+\text{moving interface}
 \Longrightarrow \text{half-space uniqueness},\\
 \text{zero extension}+\text{half-space uniqueness}
 \Longrightarrow F=0
 \Longrightarrow \text{global phase}.
\end{gathered}
\]
No arrow points to an unproved PDE theorem.  The only leaves are the
foundational analytic operations listed in the front matter.

\section{Statements that should not be encoded as axioms}
A Lean project claiming the full theorem should not introduce declarations
such as
\begin{verbatim}
axiom strichartz_1d : ...
axiom local_smoothing : ...
axiom halfspace_unique_continuation : ...
\end{verbatim}
Those are central mathematical content, not harmless interfaces.  The module
order above is designed so that each can be implemented as a theorem derived
from earlier modules.  Abstract interfaces are appropriate only after the
concrete theorem has been proved and registered as an instance or wrapper.

\chapter{Final audit of self-containment}\label{app:audit}

\section{Checklist}
The following items are proved within the document rather than cited.  In
this expanded edition the list includes the representation and closure
lemmas that were previously compressed.
\begin{enumerate}
\item The free dispersive estimate and every Strichartz estimate used in the
      nonlinear construction.
\item Local and global $L^2$ well-posedness, stability, mass conservation,
      and gauge covariance of the cubic flow.
\item The endpoint $L^4_tL^\infty_x$ estimate and convergence of smooth
      approximants in that norm.
\item Local smoothing by one half derivative, including the Duhamel term.
\item Critical $H^{1/2}$ multiplication, Lipschitz composition, mollification,
      and the zero-set regularizer.
\item Recovery of the current from the density at $L^2$ regularity.
\item The Wronskian identity without division by the solution.
\item The exterior-product equation and every integrability estimate for its
      coefficient and forcing.
\item Distribution-valued Dirichlet and Neumann traces on the diagonal.
\item The zero-extension jump formula.
\item The uniform exponential-weight Carleman estimate from an explicit
      Fourier inverse and kernel bounds.
\item Moving-interface half-space propagation with an explicit quantitative
      iteration.
\item The Hilbert-space step from zero exterior product to one unimodular
      scalar.
\end{enumerate}

\section{What is not claimed}
This TeX file is not itself executable Lean code, and compiling the PDF is not
a kernel check.  Its self-containment claim is mathematical: all
PDE-specific propositions used in the main proof occur with proofs in the
same source.  A machine-checked version requires implementing the modules in
\cref{app:ledger} and running Lean against a pinned Mathlib revision.  The
manuscript has been organized so that such an implementation does not need to
invent missing mathematical arguments or hide them behind axioms.

\section{End of proof}
Combining the chapters in the stated dependency order proves
\cref{thm:main}.  The proof uses the common density first as a conservation-law
observable, then as the real coefficient of a two-particle linear equation.
Local smoothing provides the missing half derivative at the diagonal;
Carleman propagation converts zero Cauchy data there into a vanishing exterior
product; Hilbert-space algebra supplies the single global phase.

\backmatter
\chapter*{Index of principal notation}
\addcontentsline{toc}{chapter}{Index of principal notation}
\begin{longtable}{p{0.20\textwidth}p{0.68\textwidth}}
$S(t)$ & Free one-dimensional Schr\"odinger group, multiplier
$\e^{-it\xi^2}$.\\
$\mathcal X(I)$ & Nonlinear fixed-point space
$C_tL^2_x\cap L^6_{t,x}$.\\
$\rho$ & Common density $|u|^2=|v|^2$.\\
$\mathfrak j_u$ & Endpoint current quadratic form
$\Ima\pair{u_x}{\phi\overline u}$.\\
$F$ & Exterior product $u(x)v(y)-v(x)u(y)$.\\
$(r,s)$ & Tangential and normal coordinates to the diagonal $x=y$.\\
$Q$ & Ridge potential in the exterior-product equation.\\
$G$ & Zero extension $\1_{\{s>0\}}F$.\\
$P$ & Free two-dimensional spacetime operator $i\partial_t+\Delta_z$.\\
$X,X',Y$ & Source, solution, and multiplier spaces in the Carleman argument.\\
$T_\beta$ & Explicit inverse of the conjugated free operator $P_\beta$.\\
$\zeta$ & Final constant unimodular phase.\\
\end{longtable}

\end{document}
